\documentclass[11pt,a4paper,oneside,openany]{book}
\usepackage[T1]{fontenc}
\usepackage{lmodern,amsmath,amssymb,amsthm,booktabs,geometry,xcolor,fancyhdr,needspace,fvextra}
\geometry{left=24mm,right=24mm,top=23mm,bottom=23mm,headheight=14pt,headsep=8mm,footskip=11mm}
\definecolor{ink}{RGB}{26,71,77}\definecolor{soft}{RGB}{73,88,90}
\usepackage[unicode,bookmarksnumbered,colorlinks=true,linkcolor=ink,urlcolor=ink]{hyperref}
\hypersetup{pdftitle={The Least Mills Constant: Cubic Pisot Candidates, Prime Gaps and Reciprocal Sums},pdfsubject={Consolidated mathematical research notes}}
\input{glyphtounicode}\pdfgentounicode=1
\setlength{\parindent}{0pt}\setlength{\parskip}{4.5pt}\setlength{\emergencystretch}{3em}\raggedbottom
\newcommand{\Q}{\mathbb Q}\newcommand{\C}{\mathbb C}\newcommand{\Z}{\mathbb Z}\newcommand{\F}{\mathbb F}\newcommand{\Tr}{\operatorname{Tr}}\newcommand{\Nm}{\operatorname{N}}\newcommand{\ord}{\operatorname{ord}}\newcommand{\K}{K}\newcommand{\OO}{\mathcal O_K}\newcommand{\unit}{\mathcal U}
\newcommand{\lead}[1]{\textbf{#1}\enspace}\let\para\lead
\newcommand{\sub}[1]{\par\medskip\textbf{#1}\enspace}
\newcommand{\localhead}[1]{\needspace{5\baselineskip}\section*{\color{ink}#1}}
\let\sect\localhead\let\secx\localhead\let\head\localhead
\makeatletter
\def\@makechapterhead#1{\vspace*{4pt}{\parindent\z@\raggedright\normalfont{\sffamily\small\bfseries\color{ink}\MakeUppercase{\@chapapp}\enspace\thechapter\par}\vspace{8pt}{\color{ink}\huge\bfseries #1\par}\vspace{9pt}{\color{ink}\hrule height .6pt}\vspace{13pt}}}
\makeatother
\pagestyle{fancy}\fancyhf{}\fancyhead[L]{\sffamily\footnotesize\color{soft}THE LEAST MILLS CONSTANT}\fancyhead[R]{\sffamily\footnotesize\color{soft}\nouppercase{\leftmark}}\fancyfoot[C]{\small\thepage}
\fancypagestyle{plain}{\fancyhf{}\fancyfoot[C]{\small\thepage}\renewcommand{\headrulewidth}{0pt}}
\renewcommand{\chaptermark}[1]{\markboth{#1}{}}\renewcommand{\headrulewidth}{.25pt}
\setcounter{tocdepth}{0}
\begin{document}
\frontmatter
\thispagestyle{empty}\vspace*{18mm}
{\color{ink}\sffamily\small\bfseries NUMBER THEORY\quad /\quad COLLECTED RESEARCH NOTES\par}
\vspace{8mm}{\fontsize{35}{40}\selectfont\bfseries The Least Mills Constant\par}
\vspace{4mm}{\LARGE\color{soft}Cubic Pisot candidates, prime gaps\par and reciprocal sums\par}
\vspace{9mm}{\color{ink}\hrule height 1.2pt}\vspace{8mm}
Let $M$ be the smallest real number greater than one for which $p_n=\lfloor M^{3^n}\rfloor$ is prime at every $n\ge1$. This volume studies when the observed sequence could be the trace of repeated cubes of a single algebraic integer.

A three-term arithmetic test characterizes an algebraic tail. Local trace arguments exclude whole families of cubic candidates and restrict the degree in one fixed norm-one field. A separate argument gives transcendence results for reciprocal sums along the same prime sequence.

\vfill
{\large\bfseries Scope of the results}\par
The unconditional transcendence of $M$ remains open. The degree bound $3^{1237}$ applies under the field and norm hypotheses of Chapter~\ref{ch:frob}. Chapter~\ref{ch:seven} forces a composite floor every fourth step in another field for the specified norms.
\vspace{7mm}{\color{ink}\hrule}\vspace{3mm}
{\sffamily\small Consolidated English edition \hfill 15 November 2026}\clearpage
\pdfbookmark[0]{Contents}{volcontents}\tableofcontents
\clearpage\phantomsection\addcontentsline{toc}{chapter}{Definitions and scope}\markboth{Definitions and scope}{}\section*{\color{ink}Definitions and scope}
A \emph{cubic Pisot number} is a real algebraic integer $\beta>1$ of degree three whose other two conjugates have absolute value less than one. Saito's Theorem 1.7 supplies a candidate $\beta=M^{3^t}$ if $M$ is algebraic. Both small conjugates are real. The one of larger absolute value is negative. For all sufficiently large $j$, the original floor therefore equals the field trace,
\[
 \lfloor\beta^{3^j}\rfloor=\operatorname{Tr}_{\Q(\beta)/\Q}(\beta^{3^j}).
\]
The real trace identity permits finite-field calculations, since a fixed prime dividing infinitely many of the traces must divide floors larger than itself. For arbitrary Pisot numbers in the wider problem the eventual floor may instead be the trace minus one. The individual chapters specify which situation they treat.

\textbf{A criterion on the actual minimum.} Chapter~\ref{ch:initial} proves that $M$ is transcendental exactly when its three-prime integer test fails at infinitely many indices. A failure on a finite, independently specified prime chain has the additional identification requirement described in Chapter~\ref{ch:gaps}.

\textbf{Two different cubic fields.} Chapters~\ref{ch:res}--\ref{ch:frob} treat discriminant squareclasses, with special attention to the field defined by $X^3+19X^2-19X-19$. In that field at norm $+1$, the successive conditional degree bounds are $3^{364}$, $3^{1154}$ and $3^{1237}$. Chapter~\ref{ch:seven} treats $X^3-X^2-27X+3$. Its divisor-seven theorem holds for both unit norms and for the signed norm family described there.

\textbf{Corrections carried into the collection.} The complete inert-prime coverage at $59$ already appears in corpus [075], R16. Chapter~\ref{ch:frob} gives two representatives exhibiting that coverage. The example defined by $X^3-6X^2-3X+2$ is already excluded at $3$ by [042]. The gap quotient in Chapter~\ref{ch:seven} has $3$-adic centers $-3$ and $+9$. The statement requires a norm prime to $3$.

\textbf{Editorial conventions.} The original eight notes form the eight chapters. Their local section, theorem and reference numbers have been retained. A three-digit bracket such as [075] identifies a manuscript in the Verifiable Mathematical Reasoning Corpus. Purely typographical changes, a common notation font and updated reproduction commands unify the text. The first note was translated from the Korean PDF and its formulas were reset. The computational material is included in this same source.
\mainmatter

% Chapter initial
\chapter{Foundations and reciprocal sums}\label{ch:initial}\markboth{Foundations and reciprocal sums}{}
Fix the least real number $M>1$ for which every $p_n=\lfloor M^{3^n}\rfloor$, $n\ge1$, is prime. Write $g_n=p_{n+1}-p_n^3$. This chapter studies that fixed prime sequence and the sums formed from it. Saito proved that $M$ is irrational and that the density hypothesis for the Riemann zeta function implies its transcendence. The tests below are unconditional. Whether the integer test in Theorem~1 fails infinitely often on the actual least Mills sequence remains open.

\subsection*{The canonical cubic and the integer test}
If $M$ is algebraic, Saito's Theorem~1.7 gives a totally real cubic Pisot power of $M$. Descend to the least cubic level $\beta=M^{3^t}$ as in corpus [031]. The degree formula there gives $d=[\mathbb Q(M):\mathbb Q]=3^{t+1}$ and $\beta=M^{d/3}$. Write its conjugates as $\beta,-\rho,\gamma$, where $0<|\gamma|<\rho<1$ and $\rho\le\beta^{-17/40}$. Taking odd roots when descending through cubic levels preserves these inequalities. Corpus [031] gave the lower bound $100$ for $\beta$; the finite exclusion below gives $\beta>10\,000$.

For $n\ge2$ the prime $p_n$ is odd. When $3\mid g_n$, define
\[
 h_n=g_n/3,\qquad e_n\equiv h_n\pmod{p_n},\quad |e_n|<p_n/2,
 \qquad b_n=(e_n-h_n)/p_n.
\]
Call stage $n$ a pass if these quantities exist at $n,n+1$ and
\[
 e_{n+1}=e_n^3,\qquad
 b_{n+1}=b_n^3-3p_nb_ne_n+3e_n^2.
\]
The test uses precisely $p_n,p_{n+1},p_{n+2}$.

\textbf{Theorem 1.} The least Mills constant is algebraic if and only if all sufficiently late stages pass.

\textit{Proof.} In the algebraic case put $q=3^j$. The sum of the $q$th powers of the two small conjugates belongs to $(-1,0)$ for all sufficiently large $j$, so the floor of $\beta^q$ equals the trace. If $p,b,e$ are the elementary symmetric functions of the three powers, their next trace is $p^3-3pb+3e$. As $|e|/p\to0$, centered reduction recovers the actual $e$. Newton's identities after cubing give the other two recurrences. This is the reconstruction from corpus [010].

Conversely suppose every stage from $K$ onward passes. Let $x,y,z$ be the three roots of $X^3-p_KX^2+b_KX-e_K$. Inductively, the three symmetric functions at stage $K+j$ are those of $x^{3^j},y^{3^j},z^{3^j}$. Indeed the trace recurrence follows from the definition of $h_n$, and the other recurrences are exactly the pass conditions. Put $H=M^{3^K}$ and $R=\max(|x|,|y|,|z|)$. Since $p_{K+j}^{1/3^j}\to H$, the triangle inequality gives $R\ge H$.

If $R>H$, at least two roots attain modulus $R$: a unique largest term forces the trace growth rate to be $R$. Name two such roots $x,y$. Since $|xyz|=|e_K|<p_K/2<H$, we have $|z|<H/R^2<1$. Thus $xy$ is the unique largest pair product in modulus, giving $|b_{K+j}|^{1/3^j}\to R^2$. On the other hand $p_n^3<p_{n+1}<(p_n+1)^3$ yields $g_n=O(p_n^2)$ and $|b_n|=O(p_n)$. The pair-product growth rate is at most $H$, a contradiction. Hence $R=H$. The modulus of an algebraic number is algebraic, so $H$ and its positive $3^K$th root $M$ are algebraic.\hfill$\square$

\textbf{Finite-stage degree bound.} Assume $M$ algebraic and $k\ge\max(2,t)$. Put $X=M^{3^k}>8$. The two smaller conjugates of $X$ have absolute value at most $X^{-17/40}<1/2$. Their sum is in $(-1,0)$, so the trace equals the floor. Their product times $X$ has absolute value $<p_k/2$, and the second symmetric coefficient $b_k$ is negative. To see the sign, opposite signs give two negative contributions $X(u+v)$ and $uv$; if both are negative, $X(|u|+|v|)>|uv|$. Thus $g_k=3(e_k-p_k b_k)>3p_k/2$ and $3\mid g_k$. A failure of either requirement at $k\ge2$ implies that $M$ is transcendental or $[\mathbb Q(M):\mathbb Q]\ge3^{k+2}$. A finite greedy prime tower can be used here only after identifying its terms with the actual least-Mills sequence.

\subsection*{Finite exclusion of cubic Pisot candidates}
Suppose $\beta<B=10\,000$. Write $f(X)=X^3-aX^2-d_0X-c$ for its minimal polynomial and $u,v$ for its small conjugates. Necessarily $\beta>2$. Otherwise, using $|uv|\ge1/\beta$ and $(1-u^2)(1-v^2)\le(1-|uv|)^2$, one obtains
\[
 0<|f(1)f(-1)|\le(\beta^2-1)(1-1/\beta)^2\le3/4,
\]
which is impossible for an integer. Since $-2<u+v<0$, we have $1\le a<B$. We also obtain
\[
 0<|c|=\beta\rho|\gamma|<\beta^{3/20},\qquad
 0<d_0=-(\beta u+\beta v+uv)<2\beta^{23/40}+1.
\]
The sign of $d_0$ follows as for $-b_k$ above. We may take $|c|\le3$ and $d_0\le401$. The catalogue in [031] gives additional necessary tests. If $c<0$, impose $d_0+|c|\le a$. If $c>0$, impose $d_0\le2a+2$, $d_0\le a+c$, positive discriminant, and $f(B)>0$. These follow by evaluating $f(-1)$ and using the root intervals (both small roots are negative for $c>0$). The enumeration includes $a=B$ for safety.

Start from $(a,-d_0,c)$. For each trial prime $\ell$, iterate
\[
 (T,b,e)\longmapsto(T^3-3Tb+3e,\ b^3-3Tbe+3e^2,\ e^3)\pmod\ell.
\]
A zero trace on the eventual cycle makes $\ell$ divide arbitrarily large terms. An algebraic least Mills candidate would instead have prime traces there once the terms exceed $\ell$. Beginning with $18\,064\,626$ triples gives these checkpoints:
\[
\begin{array}{c|rrrr}
\text{Last prime tested} &2&47&97&107\\\hline
\text{Surviving triples} &7\,533\,373&257&1&0
\end{array}
\]
The last survivor is $(209,-165,1)$. Its orbit modulo $107$ has period $126$ starting at the initial state; the traces at indices $58$ and $78$ are zero. The original note reports independent Python and C++ agreement at all 28 checkpoints; the standard-library Python certificate is included in this volume. Equality $\beta=10\,000$ is impossible since $\beta$ has degree three. As $\beta>10\,000$ implies $2\rho<1$, its integer trace $a>\beta-1>9\,999$. By the degree formula the minimal polynomial of $M$ is $f(X^{3^t})$. Consequently its maximum absolute coefficient is at least $10\,000$, independently of its degree.

\subsection*{Counting large prime gaps after cubes}
Let $D(m)$ be the distance from $m^3$ to the first prime strictly greater than $m^3$. Corpus [010] shows that an algebraic least Mills constant gives, for all sufficiently large $n$,
\[
  g_n=D(p_n),\qquad g_n=(p_n^3)^{\kappa+o(1)},\qquad \kappa>1/2.
\]
Minimality gives the eventual first-prime choice: every sufficiently large prime admits a continuation above its cube by the short-interval bound. Choosing a smaller admissible prime child and continuing indefinitely would construct a smaller prime-generating constant. The strict exponent inequality uses $\rho>|\gamma|$ and the positive integer $\beta\rho|\gamma|$.

\textbf{Theorem 2.} For every $\epsilon>0$,
\[
 \#\{P\le m\le2P:D(m)\ge m^{3/2}\}\ll_{\epsilon} P^{21/100+\epsilon}.
\]
\textit{Proof.} The gap $G$ between consecutive primes containing $m^3$ has $G\ge D(m)$. The Baker--Harman--Pintz bound gives $G=O(P^{63/40})$, while consecutive cubes in this interval are separated by at least $3P^2$. For large $P$ each gap contains at most one of these cubes. Partition the preceding primes into a bounded number of dyadic intervals $[X,2X]$ with $X\asymp P^3$. The selected gaps satisfy $G\ge m^{3/2}\ge\sqrt X$. J\"arviniemi bounds the sum of such gaps by $O_{\epsilon}(X^{0.57+\epsilon})$. Divide by the minimum selected length $P^{3/2}$ and adjust $\epsilon$; the resulting exponent is $3(0.57)-1.5=0.21$.\hfill$\square$

An algebraic $M$ would follow this exceptional set at every sufficiently late stage. Tower primes grow so quickly that the count can still allow one exceptional cube in each occupied dyadic interval. If $D(p_n)\le p_n^{3/2}$ for infinitely many $n$, then $M$ is transcendental. Corpus [042] constructs cubic Pisot numbers avoiding any prescribed finite set of prime divisors along the relevant odd powers, motivating tests adaptive to a candidate or to the least-prime requirement.

\subsection*{Reciprocal powers on one fixed Mills sequence}
For integers $j\ge1$ define
\[
 S_j=\sum_{n\ge1}p_n^{-j},\qquad Q_n=\prod_{i=1}^np_i,
 \qquad s_{j,n}=\sum_{i=1}^np_i^{-j}.
\]
Here every $S_j$ uses the full sequence. The distinction from splitting the sequence into odd and even indices matters at $j=1$.

\textbf{Theorem 3.} For integers $1\le u<v<2u$, at least one of $S_u,S_v$ is transcendental.

\textit{Proof.} The relative errors in replacing $p_i$ by $M^{3^i}$ are summable. Summing the exponents in the product yields
\[
  \frac{Q_n^2}{p_{n+1}}\longrightarrow
  L=M^{-3}\prod_{i\ge1}\left(\frac{p_i}{M^{3^i}}\right)^2>0.
\]
For fixed $j$, the first omitted term dominates the tail, giving $S_j-s_{j,n}\asymp Q_n^{-2j}$. First prove that $1,S_u,S_v$ are linearly independent over $\mathbb Q$. If they satisfy an integer relation, substitute the partial sums and clear denominators by $q=Q_n^v$. The corresponding linear combination of the scaled tails is an integer tending to zero because $v<2u$. If the coefficient of $S_u$ is nonzero, its tail dominates; otherwise any nonzero coefficient of $S_v$ gives a nonvanishing tail. This contradicts the integrality unless the relation is trivial.

Suppose both values are algebraic. Apply Schmidt's subspace theorem to the integral vectors $(q,qs_{u,n},qs_{v,n})$ with independent forms $X_0,S_uX_0-X_1,S_vX_0-X_2$. The product of absolute values is
\[
 q^3|S_u-s_{u,n}|\,|S_v-s_{v,n}|\asymp Q_n^{v-2u}.
\]
Their height is comparable to $q$, so this saves a fixed power of the height. Infinitely many vectors lie in a fixed proper rational subspace. Divide by $q$ and pass to the limit to obtain the excluded linear relation.\hfill$\square$

In particular, at least one of $S_2,S_3$ is transcendental. Consecutive indices $j$ for which $S_j$ is algebraic must increase by a factor of at least two. There are at most $1+\lfloor\log_2 N\rfloor$ such indices up to $N$. Transcendental $S_j$ therefore have natural density one among positive integer exponents. No integral $v$ satisfies $1<v<2$, so a separate argument is needed for the first power.

\subsection*{The first reciprocal power and the normalized limit}
\textbf{Theorem 4.} At least one of $S_1,L$ is transcendental, and $L$ is irrational.

\textit{Proof.} Write $p=p_{n+1}$, $q=Q_n^2$, $Z=q/p$. The positive sequence $Z$ decreases strictly to $L$, because its next term is $Zp^3/p_{n+2}$ and $p_{n+2}>p^3$. Eventual minimality and Baker--Harman--Pintz give $p_{n+2}-p^3=O(p^{63/40})$. Summing successive decreases, using at least cubic growth from stage to stage, gives
\[
 0<Z-L=O(p^{-57/40}),\qquad q-Lp=O(p^{-17/40}).
\]
If $L$ were rational with fixed denominator, the strictly positive distance $Z-L$ would be at least a fixed positive multiple of $1/p$. Thus $L$ is irrational. Put $A=qs_{1,n}\in\mathbb Z$. After omitting the term $1/p$ as well, the remaining tail multiplied by $q$ is $O(p^{-2})$. Consequently $qS_1-A-L=O(p^{-57/40})$.

If $S_1,L$ were both algebraic, apply the subspace theorem to $(q,p,A,1)$ and the independent algebraic forms
\[
 X_0,\qquad X_0-LX_1,\qquad S_1X_0-X_2-LX_3,\qquad X_3 .
\]
Their product of absolute values is $O(p^{-17/20})$, and the vector height is comparable to $p$. There are fixed integers $a,b,c,d$, not all zero, such that $aq+bp+cA+d=0$ infinitely often. Dividing by $q$ and taking limits gives $a+b/L+cS_1=0$. Subtract its multiple by $q$ from the original relation and take limits again. Since $p-q/L\to0$ and $A-qS_1\to-L$, we obtain $d=cL$. Irrationality forces $c=d=0$, then $a+b/L=0$ forces $a=b=0$. This contradiction proves the theorem.\hfill$\square$

Thus algebraicity of either $S_1$ or $L$ would imply transcendence of the other. This disjunction leaves the individual status of $S_1$ open. For $M$ itself, Theorem~1 reduces unconditional transcendence to infinitely many failures of an exact integer test.

\subsection*{The odd and even subsequences}
For completeness, write $T_{\rm odd}=\sum_{k\ge1}1/p_{2k-1}$ and $T_{\rm even}=\sum_{k\ge1}1/p_{2k}$ for the sums that motivated the question about $S_1$. Each is transcendental for this same fixed $M$. Indeed, let $(r_k)$ be either subsequence and $D_m=\prod_{k=1}^m r_k$. Since $\log r_k\sim 3^{2k-d}\log M$, where $d=0$ or $1$, geometric summation gives $\log D_m/\log r_m\to9/8$. The tail after $m$ terms is $r_{m+1}^{-1}(1+o(1))=D_m^{-8+o(1)}$. The partial sum has reduced denominator $D_m$: at a prime $r_i\mid D_m$ just one numerator summand $D_m/r_i$ is nonzero. The positive tail is eventually smaller than $1/(bD_m)$ for any fixed integer $b>0$, which rules out a rational limit of denominator $b$. Roth's theorem now rules out an algebraic irrational limit because infinitely many of these reduced fractions approximate the sum with exponent exceeding $2$ by a fixed amount.
  
The full sum $S_1=T_{\rm odd}+T_{\rm even}$ is also irrational: $Q_n(S_1-s_{1,n})$ is positive and of order $1/Q_n$, whereas a rational $S_1$ of fixed denominator $b$ would make $bQ_n(S_1-s_{1,n})$ a positive integer. Its transcendence remains open in this analysis. Theorem~4 gives the disjunction involving $L$; the individual parity results give no algebraic relation reconstructing $M$.

\subsection*{Sources for this chapter}
The corpus inputs are [010] for trace reconstruction and gaps, [031] for coefficient restrictions and degrees, [042] for finite-prime avoidance, and [110] for reciprocal-series truncations. External sources are K. Saito, \emph{Mills' constant is irrational}, Mathematika 71 (2025), e70027; K. Saito, \emph{Transcendency of variants of Mills' constant}, Ramanujan Journal 70 (2026), article 69, Theorems 1.7 and 1.8 (numbering of \href{https://arxiv.org/abs/2508.16068v3}{arXiv v3}); R. C. Baker, G. Harman, and J. Pintz, \emph{The difference between consecutive primes, II} (2001); O. J\"arviniemi, \emph{On large differences between consecutive primes}, \href{https://arxiv.org/abs/2212.10965v2}{arXiv:2212.10965v2}; and J.-H. Evertse, \emph{An improvement of the quantitative subspace theorem}, Section 1.


% Chapter gaps
\chapter{Power residues in the gaps of the least Mills constant}\label{ch:gaps}\markboth{Power residues in prime gaps}{}
Let $M$ be the least real number greater than $1$ for which $p_n=\lfloor M^{3^n}\rfloor$ is prime for every $n\ge1$. Write $g_n=p_{n+1}-p_n^3$. We derive two necessary conditions for algebraicity and test the first on a published finite prime chain. The unproved step is to force infinitely many violations on the orbit of this specific minimum.

\sect{1. A finite-stage degree obstruction}
\textbf{Theorem.} If $M$ is algebraic, there is an integer $t\ge0$ such that $[\Q(M):\Q]=3^{t+1}$ and for every $n\ge t$, $n\ge1$, with $p_n\ge6$,
\[
 \boxed{\quad z^{3^{n-t}}\equiv p_n\pmod{g_n}
              \quad\text{has an integer solution.}\quad}
\]
Consequently, if $z^3\equiv p_n\bmod g_n$ has no solution at one such size $p_n\ge6$, then $M$ is transcendental or its degree is at least $3^{n+1}$. Infinitely many failures on the actual orbit imply transcendence. For a failed $3^r$-th power test, the analogous lower bound is $3^{n-r+2}$ whenever $1\le r\le n+1$.

\sub{The canonical cubic.}
Saito's Theorem 1.7 supplies, under algebraicity, a power $M^{3^m}$ which is a totally real cubic Pisot number with conjugates having the shape
\[
 \beta,\ -\rho,\ \gamma,\qquad
 0<\sigma=|\gamma|<\rho\le\beta^{-17/40}<1.
\]
Following corpus [031], take cube roots in its totally real cubic field $K$ as long as the positive cube root is in $K$ and the exponent index is positive. The odd real cube roots preserve the conjugate inequalities and still have degree three since their cubes generate $K$. Denote the resulting power by $\beta=M^{3^t}$. If $t>0$, then $\beta$ is not a cube in $K$. The polynomial $X^{3^t}-\beta$ is irreducible over $K$: start with the irreducible cubic $X^3-\beta$. At each step the norm of the preceding root down to $K$ is $\beta$ (the degree is odd), so that root cannot be a cube either. The tower law gives degree $3^{t+1}$, including $t=0$.

\sub{Where the trace identity starts.}
Put $j=n-t$, $q=3^j$, $X=\beta^q$, and write $\alpha_1,\alpha_2,\alpha_3$ for the conjugates of $\beta$. If $p_n\ge6$, the sum of the two small powered conjugates lies in $(-1,0)$ since $2X^{-17/40}<1$. Thus $p_n=\sum_i\alpha_i^q$ already here and at every later stage. Newton's identity gives
\[
 g_n=-3\prod_{i<k}(\alpha_i^q+\alpha_k^q),\qquad
 \gcd(p_n,g_n)=\gcd(p_n,p_{n+1})=1 .
\]
The floors increase and $p_n^3$ is composite, so $g_n>0$.



\sect{2. A root modulo every prime power}
\sub{Odd primes other than three.}
Take $\ell\ge5$ dividing $g_n$ and reduce the three original conjugates into a finite splitting field in characteristic $\ell$. Two of their $q$-th powers sum to zero by the product identity. The remaining powered root equals $p_n\bmod\ell$ and is nonzero. If the original cubic were irreducible modulo $\ell$, Frobenius would cycle its three roots. Applying it twice to the opposite-pair relation makes every pair opposite. Odd characteristic would then force all three powered roots to vanish. Hence the cubic is reducible. Even if it has multiple roots, all its roots are in $\F_{\ell^2}$, and $p_n=w^q$ there for some nonzero $w$.

Taking the norm to $\F_\ell$ gives $p_n^2=\Nm(w)^q$. Choose integers $u,v$ with $2u+qv=1$. Then $p_n=(\Nm(w)^u p_n^v)^q$ in $\F_\ell$. The resulting unit root lifts to every $\ell$-power by Hensel's lemma since $\ell\nmid q$. For $\ell=2$, raising to the odd power $q$ is an automorphism of the units modulo every power of two.

\sub{The prime three.}
We give the prime-tail argument from [031,075] to justify its use even at the stated finite stage. For $j\ge0$ let $A_j=\Tr(\beta^{3^j})$, and let $B_j,C_j$ be the second and third elementary symmetric functions of the powered roots. The map sending $(A_j,B_j,C_j)$ to the next triple is
\[
 (a,b,c)\longmapsto(a^3-3ab+3c,\ b^3-3abc+3c^2,\ c^3).
\]
It fixes triples modulo $3$ and improves a congruence modulo $3^s$ to one modulo $3^{s+1}$ for $s\ge1$. Thus $A_k\equiv A_j\bmod 3^{j+1}$ for every $k>j$.

For every sufficiently late $k$, the prime $P=A_k$ divides neither $6$, the norm, nor the discriminant of the fixed cubic. Let $Q$ be the order of the group generated by all ratios of the reduced roots. If $v_3(Q)\le k$, choose $L\ge1$ with $3^L\equiv1$ modulo the prime-to-three part of $Q$. At each index $k+sL$, $s\ge1$, the three rooted powers are a common nonzero multiple of their values at index $k$. Therefore $P$ divides $A_{k+sL}$, contradicting the increasing prime tail. Hence $3^{k+1}\mid Q$. An irreducible reduction would give $Q\mid P^2+P+1$ (all conjugate ratios have norm one), whose $3$-adic valuation is at most one. Otherwise $Q\mid P^2-1$. It follows that $A_k\equiv\varepsilon\bmod 3^{k+1}$ at all sufficiently late $k$, with a fixed sign $\varepsilon$ because $A_k\bmod3$ is constant. The trace congruence now gives $A_j\equiv\varepsilon\bmod 3^{j+1}$ for \emph{every} $j\ge0$.

The $3$-adic logarithm identifies $1+3\Z_3$ with $3\Z_3$. Its $3^j$-th powers are precisely $1+3^{j+1}\Z_3$. The sign is its own odd power. Hence $p_n=A_j$ has a $q$-th root in $\Z_3^\times$ at every stage of the theorem. The Chinese remainder theorem joins the roots modulo all prime powers dividing $g_n$. A failed $3^r$-th power test forces $n-t<r$ whenever $n\ge t$ and gives the stated degree bound.



\sect{3. A test on certified finite data}
For $n\ge2$ with $p_n>3$, the ordinary cube test fails precisely when at least one of the following conditions holds:
\begin{enumerate}
\item $9\mid g_n$ and $p_n\not\equiv1,-1\bmod9$.
\item A prime $\ell\equiv1\bmod3$ divides $g_n$ and $p_n^{(\ell-1)/3}\not\equiv1\bmod\ell$.
\end{enumerate}
This follows from cyclicity of $\F_\ell^\times$, unit-root lifting at $\ell\ne3$, and the cube image in $\Z_3^\times$. Every witness in the second case satisfies $\ell\nmid g_{n-1}$, since otherwise $p_n\equiv p_{n-1}^3\bmod\ell$.

Let $b_1=2$ and take $b_{n+1}$ to be the least prime above $b_n^3$ through the first ten certified terms of Caldwell and Cheng. Write $d_n=b_{n+1}-b_n^3$. Their reported offsets and direct enumeration give the following data.
\begin{center}\small
\begin{tabular}{r r r l}
\toprule
$n$ & $d_n$ & $b_n\bmod d_n$ & $x^3\equiv b_n\bmod d_n$\\
\midrule
1 & 3 & 2 & soluble\\
2 & 30 & 11 & soluble\\
3 & 6 & 5 & soluble\\
4 & 80 & 7 & soluble\\
5 & 12 & 7 & soluble\\
6 & 450 & 49 & insoluble modulo 9\\
7 & 894 & 859 & soluble\\
8 & 3636 & 2911 & insoluble modulo 9\\
9 & 70756 & 57367 & insoluble modulo 7 and 19\\
\bottomrule
\end{tabular}
\end{center}
At stages 6 and 8 the residues are $4\bmod9$. At stage 9, $d_9=4\cdot 7^2\cdot19^2$, $b_9\equiv2\bmod7$, and $b_9\equiv6\bmod19$. The two obstruction values are $2^2\equiv4\bmod7$ and $6^6\equiv11\bmod19$. The following Python uses only the nine published offsets. It independently reproduces the local computations without retesting the published primality certificates.
\begin{Verbatim}[fontsize=\footnotesize,breaklines,breakanywhere]
gaps = [3, 30, 6, 80, 12, 450, 894, 3636, 70756]
b = 2
failures = []
for n, gap in enumerate(gaps, start=1):
    if not any(pow(x, 3, gap) == b % gap for x in range(gap)):
        failures.append(n)
    if n == 9:
        assert (b % 7, b % 19) == (2, 6)
    b = b**3 + gap
assert failures == [6, 8, 9]
\end{Verbatim}
If the ninth transition is also the ninth transition of the actual minimum, it rules out algebraic degree below $3^{10}=59049$. At that index $3\nmid d_9$, so the older gap divisibility test also excludes low degrees. The local tests certify independent congruence failures in the stated finite data. Identification of that prefix with the minimum requires an additional continuation argument, and a fixed finite prefix only supplies finite degree bounds.



\sect{4. One-sided approximation from successive gaps}
Put $h_n=g_n/3$, $F_n=g_{n+1}/g_n$ and
\[
 T_n=\frac{h_n^2-F_n}{3p_n^3}.
\]
These rational numbers use three consecutive outputs. Under algebraicity let $p=p_n$, $b=B_j$, $e=C_j$ as in Section~2. Then
\[
 \boxed{\quad \lceil T_n\rceil=e,\qquad
 \frac{p_n^6}{h_n^3}\bigl(\lceil T_n\rceil-T_n\bigr)\longrightarrow1.\quad}
\]
The equality holds for all sufficiently large $n$. This identifies the signed integer norm $e=\Nm(\beta)^{3^{n-t}}$ by rounding a rational number formed from the outputs. Its sign equals the sign of the cubic norm, even when the norm is negative.

\sub{Proof.}
The symmetric recurrences give $h=e-pb$ and
\[
 T_n-e=(b/p)^3-\frac{4eb}{p^2}+\frac{3e^2}{p^3}.
\]
For instance, write $p'=p^3+3h$, $b'=b^3-3pbe+3e^2$. Substitution in $F=(e^3-p'b')/h$ gives this identity by expanding. Write $q=3^{n-t}$ and $\sigma=|\gamma|$. From the three conjugates we have
\[
 p\sim\beta^q,\qquad b\sim-(\beta\rho)^q,\qquad
 |e|=(\beta\rho\sigma)^q,\qquad h/p^2\sim\rho^q .
\]
After division by $|(b/p)^3|$, the absolute sizes of the last two terms of the identity are $O((\sigma/\rho)^q)$ and $O((\sigma^2/(\beta\rho))^q)$. Both tend to zero. Therefore $e-T_n\sim\rho^{3q}\sim(h/p^2)^3$. It is positive and less than one eventually, as required. Moreover $h=e-pb$ and $|e|/p\sim(\rho\sigma)^q\to0$, so this ceiling must also be the unique integer in $(-p/2,p/2)$ congruent to $h$ modulo $p$. Consequently three outputs suffice to compare the two exact recoveries of $e$. In particular, if $\{T_n\}\le1-\delta$ for a fixed $\delta>0$ along infinitely many actual indices, $M$ is transcendental.

\sub{Minimality and the remaining implication.}
The Baker--Harman--Pintz bound gives a prime between $v^3$ and $(v+1)^3-1$ for every sufficiently large prime $v$. Every such prime starts an infinite nested chain. A prime strictly between $p_n^3$ and $p_{n+1}$ could therefore replace $p_{n+1}$ and give a smaller admissible real number at a sufficiently late index. Minimality forces the eventual least-prime rule. The available argument does not force any of the above violations infinitely often on this specified chain.

\sub{Sources.}
K.~Saito, \href{https://doi.org/10.1007/s11139-026-01443-0}{\emph{Transcendency of variants of Mills' constant}}, Ramanujan J.\ 70 (2026), 69, Theorems 1.7--1.8. C.~K.~Caldwell and Y.~F.~Cheng, \href{https://arxiv.org/pdf/1010.4883}{\emph{Determining Mills' constant and a note on Honaker's problem}}, J. Integer Seq.\ 8 (2005), 05.4.1, Section 3. R.~C.~Baker, G.~Harman and J.~Pintz, \href{https://doi.org/10.1112/plms/83.3.532}{\emph{The difference between consecutive primes, II}}, Proc. London Math. Soc.\ 83 (2001), 532--562. \href{https://github.com/QEEDDD/Verifiable-Mathematical-Reasoning-Corpus/tree/8a3f6895d0b9b9f125e43157e4a003a2a2b51704}{Verifiable Mathematical Reasoning Corpus}, manuscripts [010] (pair sums and symmetric identities), [031] (canonical field and trace limit), [075] (projective return). No claim of literature-wide priority is made for these deductions.


% Chapter res
\chapter{Quadratic discriminants excluding cubic Pisot powers}\label{ch:res}\markboth{Quadratic discriminants}{}
Let $M$ be the least real number $>1$ for which $\lfloor M^{3^n}\rfloor$ is prime for every $n\ge1$. If $M$ is algebraic, Saito's Theorem~1.7 gives a power $M^{3^t}$ which is a totally real cubic Pisot number [1]. We give an unconditional exclusion for entire discriminant squareclasses of such candidates. A Pisot number is a real algebraic integer $>1$ whose other conjugates all have modulus $<1$.

\secx{1. A uniform exclusion}
Let $\beta$ be a totally real cubic Pisot number, with minimal polynomial
\[
 f(X)=X^3-aX^2+bX-C,\quad \Delta=\operatorname{disc}(f)=D m^2,
 \quad F_j=\lfloor\beta^{3^j}\rfloor\quad(j\ge0).
\]
Here $C=\Nm(\beta)$, $m$ is a positive integer and $D$ is positive and squarefree. For an odd prime $p$ write $\chi_p(x)=(x/p)$, extended by zero at multiples of $p$. For odd squarefree $d$ put $\chi_d=\prod_{p\mid d}\chi_p$, with $\chi_1=1$.

For a prime $\ell\equiv1\bmod12$, let $H_\ell$ be the subgroup of $\F_\ell^\times$ of elements of order prime to $3$. Say that $\ell$ satisfies $(*)$ if
\[
 \boxed{\ \text{every cycle of }z\mapsto z^3\text{ in }H_\ell
   \text{ contains a }z\text{ with }\chi_\ell(1+2z)\in\{0,1\}.\ }\tag{*}
\]
\lead{Theorem.} Let $d$ be a product of an even number of distinct primes of the form $2\cdot3^s+1$, $s\ge1$. Allow $d=1$. Set $D=d$ or $D=d\ell$ for a prime $\ell$ satisfying $(*)$, and take $D>1$. Suppose
\[
 C\equiv1\bmod3,\qquad \chi_d(C)=1,\qquad
 \chi_\ell(C)=1\ \text{if the factor }\ell\text{ is present}.
\]
There are integers $J\ge0$ and $L\ge1$ such that $F_{J+kL}$ is composite for every $k\ge0$. In particular the composite indices have positive lower density. The progression and the density bound may depend on $\beta$.

\lead{Application.} Every norm-$+1$ cubic Pisot unit in these squareclasses satisfies the hypotheses. For such $\beta$, the identity $\beta=M^{3^t}$ for any integer $t\ge0$ would make all sufficiently late $F_j$ prime. Thus all these candidates are excluded. The finite sets in Section~4 give \textbf{4,863 distinct squareclasses} to which this applies, without bounding the coefficients of $f$. The known reduction to a cubic does not force $C=1$ or any of these $D$.

\lead{Relation to the corpus.} Manuscript [009] gives a quadratic-character cycle test for an individual cubic in the exceptional trace branch. Manuscript [042] proves a uniform norm-class result for cyclic cubics. Here the squareclass is nonsquare, and we determine when ramified local cycles force the character test for every eligible generator. Positive signs can result from cancellation across distinct prime factors of $D$.



\secx{2. Floor shifts and prime returns}
Denote the three real roots of $f$ by $\alpha_1=\beta,\alpha_2,\alpha_3$. For $q=3^j$ let $A_j=\sum_i\alpha_i^q$, $B_j=\sum_{i<k}(\alpha_i\alpha_k)^q$ and $C_j=C^q$. The two smaller roots have different absolute values. Otherwise they are negatives of one another and $\beta=a$ is rational. Their odd power sum has a fixed sign and eventually has modulus $<1$. Consequently $F_j=A_j-h$ eventually, with $h=0$ for a negative sum and $h=1$ for a positive sum.

\lead{The $3$-adic limit.} Newton's identity sends the integer state $(A_j,B_j,C_j)$ to its successor by
\[
 \Phi(x,y,c)=(x^3-3xy+3c,\ y^3-3xyc+3c^2,\ c^3).
\]
The state is fixed modulo $3$. If two states agree modulo $3^v$, $v\ge1$, their images agree modulo $3^{v+1}$ by the difference of cubes and the factors $3$ in the other terms. Thus $A_j$ has a limit $\lambda\in\Z_3$ with $A_j\equiv\lambda\bmod3^{j+1}$. Initial states equal modulo $3$ have $j$th traces equal modulo $3^{j+1}$.

\lead{Prime returns.} If $\lambda-h\ne1,-1$, the two valuations $v_3(F_j\pm1)$ are bounded for large $j$. If there are arbitrarily late prime terms, take one $p=F_j$ avoiding $6C\Delta$. The splitting degree $e$ of $f\bmod p$ is $1,2$ or $3$. For all three, $v_3(p^e-1)$ is bounded uniformly for these primes. For $e=2$ factor $p^2-1$. For $e=3$ the extra factor $p^2+p+1$ has valuation $1$ when $p\equiv1\bmod3$ and $0$ otherwise. For sufficiently late $j$, all $\alpha_i^{3^j}$ modulo $p$ have order prime to $3$ in $\F_{p^e}^\times$. Cubing acts as a permutation on its subgroup of order prime to $3$. Hence the three powered roots return simultaneously after some $L\ge1$ iterations, giving $p\mid F_{j+kL}$. These terms eventually exceed $p$. If there are only finitely many prime terms, the whole tail is composite.

\lead{Exceptional limits.} Since $C\equiv1\bmod3$, the shift $h=1$ always falls under the preceding argument. Indeed, $\lambda=0$ would give $a\equiv0\bmod3$ and $A_1\equiv3\bmod9$. For $\lambda=2$, first $a\equiv2\bmod3$ and $A_1\equiv2+3b\bmod9$ imply $b\equiv0\bmod3$. Two iterates of $(2,0,1)$ give a trace of $1235\equiv20\bmod27$, excluding $2$.

For $h=0$, $\lambda=1$ implies $a\equiv b\equiv1\bmod3$ by $A_1\bmod9$. Then $f=(X-1)(X^2+1)\bmod3$ has discriminant $2\bmod3$. But $D\equiv1\bmod3$, as each of its prime factors is $1\bmod3$. Thus $\Delta=Dm^2$ is $0$ or $1\bmod3$. Only $h=0,\ \lambda=-1$ remains. In this case $F_j=A_j\equiv-1\bmod3$ for large $j$.



\secx{3. Ramified cycles and the discriminant character}
For each prime $v\mid D$ the hypotheses give $v\nmid C$ and $v\mid\Delta$. The cubic modulo $v$ has the form $(X-r)^2(X-s)$ with $r,s\in\F_v^\times$, including $r=s$. A repeated irreducible factor of a cubic must be linear. For $q=3^j$,
\[
 C\equiv r^2s,\qquad A_j\equiv 2r^q+s^q=s^q\bigl(1+2(r/s)^q\bigr)\pmod v.
\]
\lead{The factors of $d$.} If $v=2\cdot3^u+1$, then for $j\ge u$ each of $r^q,s^q$ is its quadratic character, identified with $\pm1\bmod v$. In particular $s^q=\chi_v(C)=:\epsilon$. Since $v\equiv7\bmod12$, $\chi_v(-1)=\chi_v(3)=-1$. The two possibilities for $2r^q+\epsilon$ give
\[
 \chi_v(A_j)=-\epsilon=-\chi_v(C).
\]
For $\epsilon=1$ the choices are $3,-1$, and for $\epsilon=-1$ they are $1,-3$. Multiplication over $v\mid d$ gives $\chi_d(A_j)=\chi_d(C)=1$ at all sufficiently late indices.

\lead{The optional factor $\ell$.} Here $\chi_\ell(s)=\chi_\ell(C)=1$. The eventual values $z=(r/s)^{3^j}$ form a cycle in $H_\ell$. Condition $(*)$ supplies a residue class of indices modulo its period on which $\chi_\ell(A_j)$ is always $0$ or always $1$. If $\ell$ is absent every late index is eligible. Thus $\chi_D(A_j)$ is $0$ or $1$ along a fixed progression.

\lead{Finishing the exceptional branch.} A zero gives a fixed divisor $\ell$ along the progression. For sign $+1$, if only finitely many terms there are prime, its tail is composite. Otherwise choose a sufficiently late prime $p=F_j=A_j$ there avoiding $6C\Delta$. As $D\equiv1\bmod4$, quadratic reciprocity gives $(\Delta/p)=(D/p)=\chi_D(p)=1$. For a separable cubic modulo $p$, the Frobenius permutation of the three roots has sign $(\Delta/p)$: Frobenius multiplies the Vandermonde product by its permutation sign. Thus the roots split over $\F_{p^3}$, since an even permutation on three letters has order $1$ or $3$. Now $p\equiv-1\bmod3$, so $3\nmid p^3-1$. Cubing permutes $\F_{p^3}^\times$ and the prime-return argument of Section~2 gives a composite progression. This completes the theorem.

\lead{A hand check at $13$ and $37$.} More generally, if $\ell=4\cdot 3^u+1$ is prime with $u\equiv1,2\bmod4$, then $H_\ell=\{1,-1,i,-i\}$ with $i^2=-1$. The fixed points give $3,-1$, both squares. On $i\leftrightarrow-i$ one has $(1+2i)(1-2i)=5$ and $\chi_\ell(5)=-1$, since $\ell\equiv2,3\bmod5$. Thus at least one point has character $+1$. The choices $u=1,2$ give $13,37$.



\secx{4. Finite certificates and the first failed cycle}
Among the 70 primes $\ell<2000$ congruent to $1\bmod12$, the following 37 satisfy $(*)$. The enumeration concerns \emph{all} cycles for each prime.
\[
\begin{gathered}
13,37,97,157,181,229,277,349,373,397,433,457,541,\\
613,709,733,757,853,877,937,997,1021,1033,1117,1213,\\
1237,1297,1549,1597,1621,1669,1693,1741,1777,1789,1933,1993 .
\end{gathered}
\]
Choose the factors of $d$ from
\[
7,\ 19,\ 163,\ 487,\ 1459,\ 39367,\ 86093443,\ 258280327 .
\]
These are $2\cdot3^s+1$ with $s=1,2,4,5,6,9,16,17$ respectively. There are $2^7$ subsets of even cardinality and $38$ choices for the optional prime, including absence. The lists are disjoint, so unique factorization gives $2^7\cdot38-1=4863$ different $D>1$. This counts squareclasses covered by the criterion. It gives no proportion of all candidates.

\lead{Reproducible calculation (Python 3, standard library).}
{\small
\begin{Verbatim}[fontsize=\footnotesize,breaklines,breakanywhere]
from math import isqrt

def prime(p):
    return p >= 2 and all(p % d for d in range(2, isqrt(p) + 1))

def all_cycles_good(p):
    h = p - 1
    while h % 3 == 0:
        h //= 3
    seen = set()
    for x in range(1, p):
        if x in seen or pow(x, h, p) != 1:
            continue
        z, good = x, False
        while z not in seen:
            seen.add(z)
            good |= pow(1 + 2*z, (p-1)//2, p) != p-1
            z = pow(z, 3, p)
        if not good:
            return False
    return True

primes = [p for p in range(13, 2000, 12) if prime(p)]
print(len(primes), [p for p in primes if all_cycles_good(p)])
assert all(prime(2*3**s+1) for s in (1,2,4,5,6,9,16,17))
\end{Verbatim}
}
\lead{Where the test stops.} At $\ell=61$ the cycle $11\leftrightarrow50$ gives $1+2z=23,40\bmod61$, both nonsquares. Even local norm $1$ is possible: take repeated root $11$ and simple root $-1$, so that $11^2(-1)=1\bmod61$. The two traces are $21,38$, also nonsquares. Failure of $(*)$ says nothing by itself about primality of the integer floors.



\secx{5. Witnesses inside Saito's conjugate bounds}
These two minimal polynomials have $C=1$. The table gives exact divisibility progressions for their larger roots.
\[
\begin{array}{clll}
D & f(X)& \Delta/D & 7\mid F_j\\ \midrule
13&X^3-13913X^2-1297X-1 &4921410^2&j\equiv1\bmod4\\
37&X^3-647X^2-181X-1    &18492^2 &j\equiv1\bmod6
\end{array}
\]
Use $k=20$ for $D=13$ and $k=4$ for $D=37$. The signs of $f$ are $-,+,-,-,+$ at $-1,-1/k,0,a,a+1$. All three roots are located and the only possible rational roots $\pm1$ are excluded. The two small roots are negative and their sum is $a-\beta\in(-1,0)$, so $F_j=A_j$ even at $j=0$. The pairs $(A_j,B_j)\bmod7$, starting at $j=0$, are
\[
\begin{array}{cl}
13 &(4,5),(0,5),(3,2),(5,0),(2,3),(0,5),\ldots\\
37 &(3,1),(0,2),(3,4),(1,3),(2,0),(4,3),(3,1),\ldots .
\end{array}
\]
Also $f\bmod3=(X+1)^2(X-1)$ for both. The state $(-1,-1,1)$ is fixed by $\Phi$ and agrees with their initial states modulo $3$, so $\lambda=-1$.

\lead{Infinitely many units in the strict Saito range.} In each field $K=\Q(\beta)$ multiply the seed by totally positive units congruent to $1$ modulo $21\mathcal O_K$. They form a finite-index subgroup of the units. By Dirichlet's unit theorem its log-absolute-value vectors form a full lattice in $x_1+x_2+x_3=0$. A translate of this lattice comes within bounded distance of $T(1,-0.46,-0.54)$ for arbitrarily large $T$. In the three embeddings, the values of the resulting unit are $B,-\rho,-\sigma$ with
\[
 0<\sigma<\rho<\min(\sigma^{17/23},B^{-17/40})<1.
\]
Indeed $0.46>17/40$ and $0.46>0.54\cdot17/23$. Each $B$ generates the cubic field, since its embedding values differ. Its discriminant has the same squareclass. The characteristic-polynomial coefficients stay unchanged modulo $21$, by reduction of the multiplication matrix in an integral basis. Both the exceptional $3$-adic limit and the indicated divisibility by $7$ persist for all sufficiently late indices in the progression.

\lead{A global witness to failure at 61.} Let $\theta$ solve $X^3+19X^2-19X-19=0$. This cubic is Eisenstein at $19$ and has discriminant $61\cdot114^2$. Its roots lie in $(-21,-19),(-1,0),(1,2)$. For $w=-17\theta-60$ the signs are $(+,-,-)$ and its minimal polynomial is $X^3-143X^2-33451X-1182913$. Its norm is $1\bmod183$. Modulo $61$ the roots of the original cubic are $3,3,-25$, hence those of $w$ are $11,11,-1$. Modulo $3$ its polynomial is $(X+1)^2(X-1)$. Multiplying $w$ by totally positive units equal to $1\bmod183\mathcal O_K$ preserves these residues. Its log vectors form a lattice translate in the plane of sum $\log|\Nm(w)|$. The same ray direction gives infinitely many strict-range Pisot numbers exhibiting the failed cycle of Section~4. Their common norm is $1182913$.



\secx{6. Scope and sources}
\lead{Target.} Saito's Problem~1.1 [2] asks for infinitely many composite terms of $\lfloor\beta^{3^j}\rfloor$ for \emph{every} cubic Pisot number. For the least Mills constant, Saito's stronger reduction [1] requires only a totally real cubic with the indicated sign and size conditions, if $M$ is algebraic. Section~1 proves compositeness on an arithmetic progression in the specified norm and discriminant classes. Under the Density Hypothesis of the Riemann zeta function, transcendence of $M$ is already due to Saito [1, Theorem~1.8].

\lead{What remains.} Even for local norm $+1$, the character cycle need not contain a square at every discriminant. Saito's cubic may also have other norms. The universal implication from algebraicity to contradiction on the actual orbit remains open. The even- and odd-indexed reciprocal subseries do not supply an algebraic relation with $M$ in this argument.

\lead{G1--G6 audit.} We follow the meanings of the six gates in corpus [135]. G1 identifies the earlier problem in [2]. G2 retains its original base and exponent sequence and records the restriction on $\beta$. G3 is supported by the published problem and the conditional result in [1]. G4 is established only for the subfamily in Section~1. For G5, the additional mathematical step relative to [009] is to reduce the test to the orbit of the ratio of the repeated and simple roots and to combine deterministic signs at distinct ramified primes. For G6, we have no verification of historical priority. This note is a research draft, without external peer review.

\medskip
\lead{References}
\begin{enumerate}
\item K. Saito, \emph{Transcendency of variants of Mills' constant}. \href{https://arxiv.org/abs/2508.16068}{arXiv:2508.16068}.
\item K. Saito, \emph{Intervals without primes near an iterated linear recurrence sequence}, Problem~1.1. \href{https://arxiv.org/abs/2504.14968}{arXiv:2504.14968}. \href{https://doi.org/10.1142/S1793042126500685}{doi:10.1142/S1793042126500685}.
\item K. Saito, \emph{Mills' constant is irrational}, Mathematika 71 (2025), e70027.\newline\href{https://doi.org/10.1112/mtk.70027}{doi:10.1112/mtk.70027}.
\item QEEDDD, \emph{Verifiable Mathematical Reasoning Corpus}, [009], \emph{Composite Terms in Sparse Cubic Pisot Floors}, section on the quadratic-character cycle sieve. See also [042], \emph{Cubic Pisot Floors}, section on cyclic cubic norm classes, and [135], \emph{Certified Prime Extensions and Open Problem Audit}. \href{https://github.com/QEEDDD/Verifiable-Mathematical-Reasoning-Corpus/tree/8a3f6895d0b9b9f125e43157e4a003a2a2b51704}{Repository snapshot, commit 8a3f689}.
\end{enumerate}


% Chapter d61
\chapter{Composite floors in discriminant squareclass 61}\label{ch:d61}\markboth{Discriminant squareclass 61}{}
Let $\theta$ be a root of $f=X^3+19X^2-19X-19$, and let $K=\Q(\theta)$.
We distinguish its three real embeddings by increasing $\theta$-value. Write
\[
 C=1\,182\,913,\quad w=-17\theta-60,\quad
 F_j(\beta)=\lfloor\beta^{3^j}\rfloor,\quad A_j(\beta)=\Tr_{K/\Q}(\beta^{3^j}).
\]
Here $\beta$ denotes a cubic algebraic integer in $K$, with a specified real embedding. It is Pisot if it is $>1$ in that embedding and the other two absolute values are $<1$.

\sect{1. An entire norm class}
\lead{Theorem 1.} Every cubic Pisot number in $K$ of norm $C$ has integers $J\ge0$ and $L\ge1$ for which $F_{J+kL}(\beta)$ is composite for every $k\ge0$. This holds in all three real embeddings. In the originally exceptional family $w\varepsilon$, with $\varepsilon$ a totally positive unit, the divisor can be chosen to be $7$ and $L=4$.

\lead{Why this class.} The polynomial is Eisenstein at $19$ and has discriminant $61\cdot114^2$. Modulo $61$ its roots are $3,3,-25$. Hence those of $w$ are $11,11,-1$. Cubing alternates the two traces $21,38\bmod61$, both nonsquares. This obstructs the individual quadratic-character test in corpus [009]. Totally positive units congruent to $1\bmod183\OO$ preserve this local obstruction. Their logarithms form a full lattice in the norm-zero log plane. Translates of that lattice come within bounded distance of the ray $t(1,-0.46,-0.54)$ in their affine plane, so there are infinitely many Pisot numbers in this family with conjugates $B,-\rho,-\sigma$ satisfying
\[
 0<\sigma<\rho<\min(B^{-17/40},\sigma^{17/23}).
\]
These are the strict versions of Saito's bounds [1, Theorem 1.7]. Theorem 1 excludes every one of them, and also every other Pisot generator of the same positive norm in $K$.

\lead{Connection to the original constant.} Suppose $M$ is the least real $>1$ with all $\lfloor M^{3^n}\rfloor$ prime, $n\ge1$. If $M$ is algebraic, Saito gives a totally real cubic Pisot power $M^{3^t}$ with the bounds above, allowing weak inequalities at the upper endpoints. Such a power cannot be a norm-$C$ element of the field $K$ specified here. Saito's theorem does not prescribe $K$ or the norm. His Problem 1.1 [2] asks about \emph{every} cubic Pisot number.

\lead{Further results.} Within $K$, the norm-$1$ case reduces to one real sign class. Section 6 extends the exact-norm exclusion to infinitely many squarefree norms. A finite-field certificate excludes the norm-one Saito candidates up through $6.34\cdot10^{108}$. We also construct infinitely many strict-range norm-$1$ candidates whose eventual floors avoid every prime through $3000$. Sections 4--5 distinguish these two quantifiers.



\sect{2. The norm equation and the two local tests}
Direct evaluation gives $\Nm(w)=23\cdot51431=C$. Both factors are prime (trial division up to $226$ suffices for $51431$), and each is prime to the polynomial discriminant. The reductions of $f$ have factorizations
\[
\begin{array}{c|l}
23&(X-10)(X^2+6X-5)\\
51431&(X+21181)(X^2-21162X+11138).
\end{array}
\]
The quadratic factors have nonsquare discriminants in the indicated fields. Thus each prime has a unique degree-one prime ideal in $K$. The ideal of every integral element of norm $C$ is the product of these two prime ideals, because $C$ is squarefree. All such elements are $w$ times a unit of norm $1$.

\lead{Signs modulo squares.} Put
\[
 E=2\theta^2+37\theta-56,\qquad T=2\theta^2+27\theta-46.
\]
Their norms are $1$, and their signs in the three ordered embeddings are $(-,-,+)$ and $(+,-,-)$ respectively. For exact sign checks, the three roots lie in $(s,s+1)/10^{12}$, where the respective lower numerators are $-19906514237356,-623731016520,1530245253875$. Substitution in $f,E,T$ and rational interval bounds certify the signs. Dirichlet's theorem gives the norm-one units rank two. Their torsion is trivial ($-1$ has norm $-1$). Their quotient by squares therefore has order four. The four signs exhibited by $E,T$ exhaust it. Every such unit takes the form
\[
 E^aT^b u^2,\qquad a,b\in\{0,1\},\quad \Nm(u)=1.
\]
In particular every totally positive unit is a square. This argument needs no fundamental-unit computation.

\lead{At 3.} There is a simple root $1$ of $f\bmod3$. Hensel's lemma gives a linear factor over $\Q_3$, and the squareclass $61$ is a square in $\Q_3$. Thus the complementary quadratic splits, and there are three embeddings into $\Q_3$. Their reductions of $\theta$ are $1,-1,-1$. In this order the signs of the reduced units $E,T,w$ are $(+,-,-),(+,+,+),(+,-,-)$, where signs stand for $1,-1$ in $\F_3$. Unit squares reduce to $1$. For a $3$-adic unit $x$ of residue $\delta=\pm1$, $x^{3^j}$ tends $3$-adically to $\delta$. Thus for $\beta=wE^aT^bu^2$ the trace limit $\lambda$ is $-1$ if $a=0$ and $3$ if $a=1$.

\lead{At 7.} Since $f=(X+3)(X^2+2X+3)\bmod7$ with the quadratic irreducible, the residue algebra is $\F_7\times\F_{49}$. Modulo common nonzero scalars its unit group is cyclic of order $48$, by $(x,y)\mapsto y/x$. Use square brackets for this quotient. Here $[E]$ has order $16$: $E^8=(0,5,6)$ and $E^{16}=(2,0,0)$ in the power basis modulo $7$. After one cubing the square of any projective unit lies in $\langle[E]^2\rangle$ of order $8$. Cubing permutes that subgroup. Thus for any $u$, the class of $(wu^2)^3$ equals that of $(wE^{2k})^3$ for some $0\le k<8$.



\sect{3. Four sign classes, with no residual case}
In the following table the entry $c$ means $7\mid A_j(wE^{2k})$ for all $j\ge1$ with $j\equiv c\bmod4$.
\[
\begin{array}{c|rrrrrrrr}
 k&0&1&2&3&4&5&6&7\\\midrule
 c&1&3&2&2&3&1&0&0
\end{array}
\]
Compute the first four cubed states in $\F_7[X]/(f)$ and take the trace $3x-19y+399z$ for $x+y\theta+z\theta^2$. Every initial class has order dividing $48$, and cubing then acts with period dividing $4$ because $3^4=81\equiv1\bmod16$. Zero trace is invariant under nonzero scalars. The table therefore applies to all $wu^2$.

\lead{Floor and trace.} For totally real cubic Pisot $\beta$, the small conjugates have unequal absolute values, since otherwise they are negatives of each other and $\beta$ is rational. Their odd power sum has fixed sign and tends to zero. Consequently $F_j=A_j-h$ eventually, with $h=0$ if the sum is negative, $h=1$ if positive. Positive norm forces the two small conjugates to have the same sign.

\lead{Prime-return lemma [009].} If the $3$-adic limit $z=\lambda-h$ is different from both $1$ and $-1$, the floors are composite on a progression. Indeed, either only finitely many floors are prime, or arbitrarily late prime floors $p=F_j$ avoid the norm and discriminant and $p>3$. The valuations $v_3(p-1),v_3(p+1)$ are then bounded independently of $j$. The splitting degree $e$ of the cubic modulo $p$ is $1,2$ or $3$, so $v_3(p^e-1)$ is also bounded (use the factors of $p^2-1$ and $p^3-1$). For sufficiently late $j$, all three powered roots modulo $p$ have order prime to $3$. Cubing permutes that subgroup. Their simultaneous return preserves $A_j-h\bmod p$ on a progression. The growing later floors exceed $p$ and are composite.

\lead{Proof of Theorem 1.} Write $\beta=wE^aT^bu^2$ as above. If $h=1$, the limit $\lambda-1$ is $-2$ or $2$, and the lemma applies. For $h=0$ and $a=1$, every sufficiently late floor is divisible by $3$. For $a=b=0$ the table supplies the prime $7$. The possibility $a=0,b=1,h=0$ is empty: $w$ and $T$ have the same real signs $(+,-,-)$, making $wTu^2$ totally positive. This finishes every parity, every choice of embedding, and both floor shifts.

\lead{The norm-one case.} For $\beta=E^aT^bu^2$ the trace limit is $3$ when $a=0$ and $-1$ when $a=1$. The same lemma covers $h=1$, and $h=0,a=0$ gives the divisor $3$. For $a=1,b=0$, observe $[w^3]/[E^3]$ belongs to $\langle[E]^2\rangle$. Indeed $w^3/E^3$ is a projective square of order dividing $16$ (direct reduction gives $[w^3]=[E]^5$). Thus the same eight classes after the first cubing apply. Only $\beta=ETu^2$ remains. It has signs $(-,+,-)$, so the middle $\theta$-embedding is the only possible dominant one.



\sect{4. An explicit lower bound for surviving units}
Use the following 22 primes. They are inert in $K$. The modular certificate in Section 7 verifies the additional property used here.
\[
\begin{gathered}
5,41,97,103,127,131,167,199,271,293,773,827,\\
859,911,1021,1163,1483,1583,1709,1723,1811,1993 .
\end{gathered}
\]
Let $P$ be their product. For each prime $p$ on the list, the projective unit group of $\F_{p^3}$ has order $S=p^2+p+1$. Its subgroup $H$ of order $m=S/\gcd(S,3)$ has prime order. Every nonidentity class of $H$ visits a zero-trace class under repeated cubing. This is a finite computation for each $p$, including when $3$ does not generate $(\Z/m\Z)^\times$.

If $\Nm(\beta)=1$, its projective class already lies in $H$. Indeed the norm-one subgroup in $\F_{p^3}^\times$ has order $S$ and intersects $\F_p^\times$ in $\gcd(3,p-1)=\gcd(3,S)$ elements. A nonidentity class therefore yields $p\mid A_j$ on a progression. For a possible $h=0$ prime tail, $\beta$ must reduce to a scalar at all 22 primes. The Chinese remainder theorem gives $\beta=t+P\eta$ for $t\in\Z$ and $\eta\in\OO$. Consequently,
\[
 \operatorname{disc}(\beta)=P^6\operatorname{disc}(\eta)\ \ge\ 61P^6.
\]
Indeed $\eta$ generates the same totally real cubic field and its discriminant has squareclass $61$. In Saito's range $|\beta_i|\le\beta^{-17/40}$ for both small conjugates. Their separation is at most $2\beta^{-17/40}$. The other two differences each have absolute value $<2\beta$. Taking the square root of the discriminant gives
\[
 \sqrt{61}\,P^3<8\beta^{63/40},\qquad
 \boxed{\ \beta>(\sqrt{61}\,P^3/8)^{40/63}>6.34\cdot10^{108}.\ }
\]
This bounds the hypothetical \emph{cubic power} $\beta=M^{3^t}$. There is no bound here on the exponent $t$.

\lead{The 61 test within this class.} In the norm-one case, the residual roots at $61$ are $r,r,s$ with $s=r^{-2}$. In the previously identified exceptional cycle $(r/s)^{3^j}=\pm11$ for all $j\ge1$. The powered roots lie in the subgroup of order $20$. Cubing permutes that subgroup, and the cube roots there of $\pm11$ are $\mp11$. Thus the individual repeated roots alternate $11,-11$ and the simple root equals $-1$. A hypothetical prime tail therefore alternates $21,38\bmod61$ and every sufficiently late difference $F_{j+1}-F_j^3$ is $49\bmod61$. These constraints and the large lower bound coexist with infinitely many candidates, as the next construction shows.



\sect{5. A genuine obstruction inside the norm-one sign class}
Let $Z=ET=3564-2150\theta-117\theta^2$, of signs $(-,+,-)$, and let $\mathrm{odd}(m)$ denote the largest odd divisor of $m$. Define
\[
 n=\mathop{\rm lcm}_{p\le3000,\ p\ {\rm prime}}\mathrm{odd}\bigl(p(p^2-1)(p^3-1)\bigr),
 \qquad B_0=Z^n .
\]
This positive odd exponent has 2,874 decimal digits. The residue of $Z$ at the double root $\theta=3\bmod61$ is $26$, of order $60$. Since $15\mid n$ and $n$ is odd, the corresponding residue of $B_0$ is $\pm11$, of order four. The simple residue is $-1$ by norm one. The trace residues alternate $21,38$. At $3$ the $3$-adic trace limit is $-1$.

\lead{Finite assertion.} None of the 430 primes $p\le3000$ divides $\Tr(B_0^{3^j})$ for any $j\ge0$. The verifier below uses the actual exponent $n$, reduces in $\F_p[X]/(f)$, then iterates cubing to the first repeated state. Across the 430 primes there are 1,485 distinct states counted with their prime. This verifies all indices, including at ramified primes.

\lead{Passage to genuine Pisot units.} Multiply $B_0$ by the totally positive units congruent to $1$ modulo $Q\OO$, where $Q=\prod_{p\le3000}p$. The resulting units retain every trace congruence above and the same signs. Their log vectors form a full lattice translate. Follow the ray $t(-0.46,1,-0.54)$ in the plane of coordinate sum zero with a bounded lattice approximation. For all sufficiently large $t$, the middle embedding gives a cubic Pisot number $B$ with arbitrarily large height and
\[
 0<|\beta_3|<-\beta_2<\min(B^{-17/40},|\beta_3|^{17/23}).
\]
The powers of the two negative conjugates have sum in $(-1,0)$ for all $j\ge0$ once both conjugates have absolute value $<1/2$. Thus infinitely many strict-range units meet the original bad $61$ phase and every floor avoids every prime $\le3000$. The construction gives no assertion that a floor is prime. It identifies why a sieve of this size cannot settle the last sign class.

\lead{What is open.} Theorems 1 and 2 give complete exclusions at their respective norms in this field. The unbounded norm-one sign class above, norms beyond these results in this field, and other cubic fields of discriminant squareclass $61$ are not settled in general. The result neither establishes the unconditional transcendence of $M$ nor establishes a historical priority claim.

\lead{G1--G6.} G1 identifies Saito's existing Problem 1.1. G2 keeps the original floors and makes the field and norm restriction explicit. G3 rests on the published problem and on the unconditional alternative of [1]. G4 holds for the fixed field and the families of exact norms in Theorems~1 and~2, with a bounded-height result for a second norm. The contribution relevant to G5 is the interplay of norm ideals, sign classes modulo squares, trace cycles at two different primes, and preservation of their local data in narrow ray classes. G6 is open. The earlier corpus supplied the local tests and the unit-lattice construction. This research note has no external peer review.



\sect{6. Infinitely many exact norms with the same obstruction}
\lead{Theorem 2.} In this same field there are infinitely many distinct positive squarefree norms $C'$ such that \emph{every} cubic Pisot number of norm $C'$ has a composite progression. For each such norm there are infinitely many strict-range Pisot numbers with the exact obstructing trace cycle $21,38\bmod61$. One can choose $\gg X/\log X$ such norm values $C'\le X$ for sufficiently large $X$, with a constant depending on the fixed field and modulus.

\lead{Proof.} Put $m=3\cdot7\cdot61=1281$ and prescribe positivity at all three real places in addition to the finite modulus $m\OO$. Use $s_+=43$ and $s_-=715$. Both are $1\bmod21$, are nonsquares modulo $61$, and satisfy $s_+s_-\equiv1\bmod m$.

For each sign, use the narrow ray class of the principal ideal $(s_\pm)$. The prime ideal theorem in ray classes (equivalently, class field theory and Chebotarev [4]) supplies prime ideals here with positive prime-ideal density and counting function asymptotic to $\operatorname{Li}(x)/h_m^+$, where $h_m^+$ is the narrow ray class number. Ideals above rational primes with residue degree greater than one contribute $O(\sqrt{x})$ up to norm $x$. Discard these and primes over the finite set dividing $mC\operatorname{disc}(f)$. Each remaining ideal has norm a rational prime $q_\pm$, and has a generator $\pi_\pm$ which is totally positive and satisfies $\pi_\pm\equiv s_\pm\bmod m\OO$. This is exactly the defining equivalence of ideals in a narrow ray class. Thus $q_\pm=\Nm(\pi_\pm)\equiv s_\pm^3\bmod m$ and $(q_\pm/61)=-1$.

As $61\equiv1\bmod4$, quadratic reciprocity makes the Frobenius permutation on the three roots of $f$ at $q_\pm$ odd. An odd permutation on three letters is a transposition. Hence each $q_\pm$ has splitting type $(1,2)$ and a unique degree-one prime ideal. Fix one $q_+$ and vary $q_-\ne q_+$ in the other class. Put
\[
 w'=w\pi_+\pi_-,\qquad C'=Cq_+q_- .
\]
The four rational factors of $C'$ are distinct and each has a unique degree-one prime ideal. As in Section 2, \emph{every} element of norm $C'$ has the form $w'$ times a norm-one unit. Now $w'\equiv w\bmod1281\OO$ and the real signs of $w'$ equal those of $w$. The four cases in Section 3 use precisely these data, so the proof of Theorem 1 applies to the entire norm class $C'$. Totally positive ray units preserve the residues at $61$ and provide the strict-range examples as in Section 1 of this chapter. Because $q_+$ is fixed and distinct $q_-$ give distinct values of $C'$, the ray-class prime count yields the claimed lower bound.

\lead{Difference between the quantifiers.} The prime ideal theorem supplies infinitely many \emph{norm values}. The sign and reduction argument then handles \emph{all generators} for each of them. There is no inference from a proportion of generators to the particular least Mills constant.



\sect{7. Reproducible finite-field certificates}
The two listings below are independent. The first is C++17. Compile with optimization and run. The program checks all 22 stated primes, their required orders, irreducibility, and every relevant projective orbit. Coordinates are in the power basis of $\theta$. Trace zero is unchanged by multiplying a representative by a nonzero scalar.
{\footnotesize
\begin{Verbatim}[fontsize=\footnotesize,breaklines,breakanywhere]
#include <array>
#include <cassert>
#include <vector>
using A = std::array<int,3>;
int p;
int mod(long long v) { return (v%p+p)%p; }
bool prime(int n) {
  if (n<2) return false;
  for(int d=2; d<=n/d; ++d) if(n%d==0) return false;
  return true;
}
A mul(A a,A b) {
  long long c[5]={};
  for(int i=0;i<3;++i) for(int j=0;j<3;++j) c[i+j]+=1LL*a[i]*b[j];
  for(int k=4;k>=3;--k) {
    c[k-1]-=19*c[k]; c[k-2]+=19*c[k]; c[k-3]+=19*c[k];
  }
  return {mod(c[0]),mod(c[1]),mod(c[2])};
}
int main() {
  int primes[]={5,41,97,103,127,131,167,199,271,293,773,827,
                859,911,1021,1163,1483,1583,1709,1723,1811,1993};
  for(int v:primes) {
    p=v; assert(prime(p));
    for(long long x=0;x<p;++x) assert(mod(x*x*x+19*x*x-19*x-19));
    int S=p*p+p+1, m=S%3?S:S/3;
    assert(prime(m) && m%3);
    A e={mod(-56),mod(37),2}, g={1,0,0}, a={1,0,0};
    for(int i=0;i<S/m;++i) g=mul(g,e);
    assert(g[1] || g[2]);
    std::vector<bool> hit(m,false);
    for(int i=0;i<m;++i) {
      if(mod(3*a[0]-19*a[1]+399*a[2])==0)
        for(int j=i;!hit[j];j=(3*j)%m) hit[j]=true;
      a=mul(a,g);
    }
    assert(a[0] && !a[1] && !a[2] && !hit[0]);
    for(int i=1;i<m;++i) assert(hit[i]);
  }
}
\end{Verbatim}
}
Here $g$ is a generator of $H$ since its class is nonidentity of prime order $m$. Marking an orbit forward from a zero marks exactly the classes that visit it, since cubing is invertible on $H$.



\sect{8. Small-prime barrier certificate and sources}
Python 3 with the standard library suffices. The computation of $B_0\bmod p$ below uses the full exponent and does not assume a description of the unit group of the residue algebra. Raising a possible repeated root causes no issue.
{\footnotesize
\begin{Verbatim}[fontsize=\footnotesize,breaklines,breakanywhere]
from math import isqrt, lcm
def prime(p):
    return p >= 2 and all(p % d for d in range(2, isqrt(p)+1))
def mul(a, b, p):
    c = [0]*5
    for i in range(3):
        for j in range(3):
            c[i+j] += a[i]*b[j]
    for k in (4,3):
        c[k-1] -= 19*c[k]; c[k-2] += 19*c[k]; c[k-3] += 19*c[k]
    return tuple(x % p for x in c[:3])
def power(a, n, p):
    z = (1,0,0)
    while n:
        if n % 2: z = mul(z,a,p)
        a = mul(a,a,p)
        n //= 2
    return z
primes = [p for p in range(2,3001) if prime(p)]
n = 1
for p in primes:
    m = p*(p*p-1)*(p**3-1)
    while m % 2 == 0: m //= 2
    n = lcm(n,m)
states = 0
for p in primes:
    z = power((3564,-2150,-117),n,p)
    seen = set()
    while z not in seen:
        seen.add(z)
        assert (3*z[0]-19*z[1]+399*z[2]) % p
        z = mul(mul(z,z,p),z,p)
    states += len(seen)
assert (len(primes),len(str(n)),states) == (430,2874,1485)
\end{Verbatim}
}
\lead{Sources.}
{\small
[1] K. Saito, \emph{Transcendency of variants of Mills' constant}, Theorems 1.7--1.8. \href{https://arxiv.org/abs/2508.16068}{arXiv:2508.16068}.

[2] K. Saito, \emph{Intervals without primes near an iterated linear recurrence sequence}, Problem 1.1. \href{https://arxiv.org/abs/2504.14968}{arXiv:2504.14968}.

[3] QEEDDD, \emph{Verifiable Mathematical Reasoning Corpus}, [009] on character cycles and prime returns, [042] on unit lattices and finite sieves, [075] on projective trace orbits, and [135] on the G1--G6 audit. \href{https://github.com/QEEDDD/Verifiable-Mathematical-Reasoning-Corpus/tree/38348353416f1f7fc50cecb93c3b554a00260c13}{Corpus snapshot}.

[4] J. S. Milne, \emph{Class Field Theory}, ray class fields and the Chebotarev density theorem. \href{https://www.jmilne.org/math/CourseNotes/CFT.pdf}{Author's course notes}.}


% Chapter ray
\chapter{Recurring divisors and exponent lattices}\label{ch:ray}\markboth{Recurring divisors and lattices}{}
Let $K=\Q(\theta)$ with $\theta^3+19\theta^2-19\theta-19=0$, with the three embeddings ordered by increasing value of $\theta$. We consider only the remaining norm-one unit class, of signature $(-,+,-)$, and take the positive, middle embedding as the real value of $\beta$. Write
\[
 E=2\theta^2+37\theta-56,\quad T=2\theta^2+27\theta-46,\quad
 Z=ET,\quad \unit=\{v\in\mathcal O_K^\times:\Nm(v)=1\}.
\]
The surviving sign class is $Z\unit^2$. Indeed $E,T$ have signatures $(-,-,+),(+,-,-)$ and norm one. Dirichlet's theorem makes $\unit$ torsion-free of rank two, so these two independent sign vectors give all four squareclasses. If $\beta$ is Pisot in our class, both small conjugates are negative. Eventually $F_j=\lfloor\beta^{3^j}\rfloor$ equals $A_j=\Tr(\beta^{3^j})$.

\sect{1. The previously constructed infinite family}
Write $Q=\prod_{p\le3000}p$, over primes, and put
\[
 n=\mathop{\rm lcm}_{p\le3000}\mathrm{odd}\!\left(p(p^2-1)(p^3-1)\right),
 \qquad  B_0=Z^n ,
\]
where $\mathrm{odd}$ denotes the largest odd divisor. The earlier construction [3] produces infinitely many strict Saito-range Pisot numbers in the family $B_0\varepsilon$, where $\varepsilon$ runs over totally positive units congruent to $1\bmod Q\mathcal O_K$. Their traces avoid every prime up to $3000$ at every index. The prior certificate verified 1,485 states across the 430 primes. Their traces alternate between $21$ and $38$ modulo $61$.

\textbf{Theorem A.} Every Pisot member of this entire family has a progression of indices on which $21503$ divides $F_j$. More generally, this holds for every Pisot $\beta\in Z\unit^2\cap\unit^{91}$. In this statement the period can be chosen to divide $211456$. Thus all sufficiently late blocks of $211456$ consecutive indices contain a composite floor.

The progression can depend on the unit. Both the old avoidance property and this new recurrence concern all sufficiently late indices. The algebraic unit $B_0$ serves as a starting point; the proof deals with its unit multiples giving genuine Pisot numbers, without requiring that $B_0$ itself be Pisot.

\textbf{Theorem B.} A hypothetical original Mills power in the remaining field, norm and sign class must have $\log_{10}\beta>2\cdot10^{186}$. If the least Mills constant $M$ is algebraic and has this type of cubic power, then
\[
                      [\Q(M):\Q]\ \ge\ 3^{364}.
\]
The bound and the conditioning on this particular field and norm are proved in Sections 4--5.



\sect{2. Why the old avoidance condition forces a new divisor}
\textbf{Two independent power-residue characters.} In each row of the table, the polynomial of $\theta$ splits into three distinct roots modulo $q$. At the indicated first two roots, take a unit residue to the $(q-1)/l$ power and express it in the order-$l$ group with generator $g$. The last two entries are the resulting exponent pairs for $(E,T)$ at the respective roots.
\[
\begin{array}{c|r|r|c|cc}
l&q&g&\text{roots}&\text{root 1}&\text{root 2}\\\midrule
7&1933&1812&71,112,1731&(1,4)&(3,4)\\
13&937&911&501,666,688&(9,10)&(9,1)
\end{array}
\]
These evaluations act on all of $\mathcal O_K$ because the two primes are prime to $\operatorname{disc}(f)=61\cdot114^2$. Both matrices have nonzero determinant modulo $l$. They therefore detect every class of $\unit/\unit^l$, a two-dimensional $\F_l$-space. If $\varepsilon$ is a totally positive unit equal to $1$ modulo $Q\mathcal O_K$, it maps to zero under both pairs. Hence $\varepsilon\in\unit^7\cap\unit^{13}=\unit^{91}$. Also $91\mid n$ (for instance $p=13$ supplies $7$ and $p=53$ supplies $13$). Thus $B_0\varepsilon$ is a $91$st power in $\unit$. Its signs put it in $Z\unit^2$.

\textbf{The projective calculation.} At $p=21503$, $f$ has a single simple linear root and an irreducible quadratic factor. In particular the residual algebra is $\F_p\times\F_{p^2}$. Its invertible elements modulo common nonzero scalars form a cyclic group of order
\[
 S=p^2-1=3\cdot 2^{11}\cdot7\cdot13\cdot827 .
\]
We use brackets for the image in this group. Multiplying by a scalar preserves zero trace. The first cubing leaves the subgroup of order $N=S/3$. The orders of $[E^3]$ and $[T^3]$ are respectively $N$ and $N/2$. Thus $[(ETv^2)^3]$ has $2$-primary order exactly $2048$ for every $v\in\unit$: in cyclic coordinates its exponent is odd. For a $91$st power the $7$- and $13$-components vanish. The two possible orders of $[\beta^3]$ are therefore $2048$ and $1693696=2048\cdot827$.

With $g=[(\theta+2)^3]$, of order $N$, put $h_m=g^{N/m}$. Each entry in the right column supplies the exponent of one zero-trace representative $h_m^e$ in a different orbit of multiplication by $3$ on $(\Z/m\Z)^\times$.
\[
\begin{array}{r|r|l}\toprule
m&\text{period}&e\\\midrule
2048&512&231,\ 793\\
1693696&211456&36681,\ 106441,\ 193855,\ 598607\\\bottomrule
\end{array}
\]
The periods times the numbers of representatives equal $\varphi(m)$ in both rows. The verification in Appendix A checks the trace of each representative and the disjointness of the cycles. Hence every possible $[\beta^3]$ returns to a trace zero in a progression with one of the two periods. This proves Theorem A.



\sect{3. Restrictions outside the $91$st powers}
The same finite-field computation, without assuming global power divisibility, yields a second obstruction throughout $Z\unit^2$. At $p=34303$ the residual algebra again has type $(1,2)$, and $p^2-1=3\cdot1024\cdot67\cdot5717$. The orders of $[E^3]$ and $[T^3]$ there are $(p^2-1)/6$ and $(p^2-1)/3$. After cubing, every unit in the sign class has full $2$-primary order $1024$.

For each possible odd component $d$, the entries count the cycles containing a zero and the total cycles among elements of exact order $2^v d$. Cycles of a fixed order have the same length.
\[
\begin{array}{r|c@{\qquad}r|c}
\multicolumn{2}{c}{p=21503,\ v=11}&\multicolumn{2}{c}{p=34303,\ v=10}\\
d&\text{hits / cycles}&d&\text{hits / cycles}\\\midrule
1&2/2&1&2/2\\
7&0/4&67&0/12\\
13&2/8&5717&8/8\\
91&4/48&383039&48/48\\
827&4/4&&\\
5789&8/8&&\\
10751&16/16&&\\
75257&96/96&&
\end{array}
\]
In particular a survivor must have no $827$-component at $21503$, and must have \emph{exactly} the odd component $67$ at $34303$. With $\beta=E^aT^b$ in the fundamental units certified in Section 4, necessary conditions include
\[
\begin{aligned}
 a+170b&\equiv0\pmod{827},&
 (a+b\bmod7,\ a-b\bmod13)&\ne(0,0),\\
 a+2537b&\equiv0\pmod{5717},&
 a+34b&\not\equiv0\pmod{67}.
\end{aligned}
\]
The conditions come from projecting $[E],[T]$ to the odd prime-order factors. For $7,13,827$ at $21503$, the logarithms of $[T]$ relative to $[E]$ on these factors are $1,12,170$; at $34303$, for $67,5717$, they are $34,2537$.

For precision, the two local tests miss the following fraction of the sign class, counted in its finite quotient:
\[
  \frac{81}{75257}\frac{66}{383039}
        =\frac{5346}{28826366023}.
\]
The first count follows by weighting the three incompletely covered rows by the number of elements. At the second prime the missed elements have odd part of exact order $67$. The two local images are independent: in the product of their $2$-primary groups $E^3,T^3$ give independent images modulo squares, by the displayed orders, hence generate the product. The odd group orders are relatively prime and the indicated generators surject onto each. Prescribing the exceptional phase at $61$ imposes $a+b=0\bmod5$, independent of these groups. This is a fraction of residue classes; the tests do not determine the primality of the floors in the remaining classes.




\sect{4. An exponent-lattice bound on the remaining sign class}
Let $K=\mathbf Q(\theta)$, where
$\theta^3+19\theta^2-19\theta-19=0$, with the embeddings ordered by
increasing $\theta$. Write
\[
 E=2\theta^2+37\theta-56,\qquad T=2\theta^2+27\theta-46.
\]
Suppose $\beta$ is a norm-one cubic Pisot number in the remaining class,
positive in the middle embedding, with Saito's size inequalities and the
exceptional $21,38$ trace cycle modulo $61$. If its floors
$\lfloor\beta^{3^j}\rfloor$ have no infinite arithmetic progression of
composite terms, then
\[
                         \log_{10}\beta>2\cdot10^{186}.
\]
The bound is for the cubic power itself. It supplies no upper bound on the
level $t$ in a hypothetical identity $\beta=M^{3^t}$.

\paragraph{The unit coordinates.}
An integral basis is $1,s,t$, where
$s=(\theta^2+18\theta-19)/3$ and $t=(\theta^2+21\theta-1)/3$.
For $x+ys+zt$ its trace, second symmetric function and norm are
\begin{align*}
 A&=3x-z,\\
 B&=3x^2-2xz-76y^2-38yz-101z^2,\\
 C&=x^3-x^2z-76xy^2-38xyz-101xz^2\\
  &\hphantom{{}={}}-228y^3-532y^2z+646yz^2+387z^3.
\end{align*}
These integral forms give integrality and the lattice discriminant is
$88084=2^2\cdot19^2\cdot61$. For $p=2,19$, testing the $p^3$ residues
$0\le x,y,z<p$ gives only $(0,0,0)$ with
$p\mid A$, $p^2\mid B$ and $p^3\mid C$. These tests exclude an integral
overlattice of index divisible by $p$: an order-$p$ quotient element would
give an integral $(x+ys+zt)/p$ with nonzero residue triple. The squarefree
factor $61$ cannot divide the index.

The two unit logarithm vectors are independent, since in their first two
coordinates the diagonal and off-diagonal products have opposite signs.
For an arbitrary norm-one unit, subtracting integer multiples of these two
vectors places its logarithms in the centered fundamental parallelogram of
their span. In embedding $r$ its absolute value there is bounded by
$\sqrt{\max(|E_r|,1/|E_r|)\max(|T_r|,1/|T_r|)}$. The three upper bounds
are respectively less than $300,71,268$. Inverting the basis embedding matrix
then gives $|x|\le230$, $|y|\le26$, $|z|\le34$. All integer solutions
of $C=1$ in this box are
\[
 (1,0,0),\quad (-24,5,1),\quad(46,15,-9),\quad(-132,-22,19).
\]
They represent $1,E,T,E^{-1}$. Hence every norm-one unit is $E^aT^b$ with
unique integer exponents. Independence alone would not have proved generation.
The rational root-isolating intervals used for these bounds are
$(n,n+1)/10^{12}$ for
$n=-19906514237356,-623731016520,1530245253875$.

The signs $(-,+,-)$ require $a,b$ odd. At the repeated residual root
$\theta=3\bmod61$, the values of $E,T$ are $12=2^8$, $53=2^{33}$
respectively, where $2$ has order $60$. The exceptional repeated root has
order $4$ after powering by $3^j$. Consequently $5\mid 8a+33b$, or
$a+b\equiv0\bmod5$.


\paragraph{Independent congruences on the exponents.}
For each of the 51 rows $(p,k)$ recorded below define
$S=p^2+p+1$ and $m=S/\gcd(3,S)$. The number $m$ is prime and $f$ is
irreducible modulo $p$. In $\mathbf F_{p^3}^\times/\mathbf F_p^\times$
the norm-one residues have projective image of order $m$,
generated by $[E]$, and $[T]=[E]^k$. Both equalities are checked by
polynomial modular arithmetic. For the first 22 primes the preceding
$D=61$ study certified that every nonidentity class of this subgroup visits
a trace-zero class under cubing. An exhaustive recheck is supplied below.
For the last 29, $p\equiv2\bmod3$, $m=S$ and $3$ generates
$(\mathbf Z/m\mathbf Z)^\times$. Thus the orbit of every nonidentity
class visits all nonidentity classes. The trace-zero plane has $p+1$
projective points, none scalar. In either case a nonidentity projective
residue yields a composite progression. It follows that
\[
                 a+k b\equiv0\pmod m
                 \qquad\text{for every row }(p,k).
\]

\begin{center}\scriptsize
\begin{tabular}{rr@{\qquad}rr@{\qquad}rr}
\hline
$p$&$k$&$p$&$k$&$p$&$k$\\\hline
5&1&1583&255229&22277&350529041\\
41&787&1709&2745692&22937&297631469\\
97&2951&1723&654871&26489&31896576\\
103&1915&1811&205927&26921&404448787\\
127&4419&1993&630756&28901&65818922\\
131&14136&3557&8372873&31277&924860164\\
167&11077&4889&3215080&32693&1008873991\\
199&8942&8681&46418403&36497&695170719\\
271&5781&10589&44659163&36749&1317741455\\
293&50847&13469&31295481&37049&1219522376\\
773&552873&16229&77599160&37649&1390873606\\
827&419023&16301&66708745&38669&243832387\\
859&50622&16421&76985174&38993&1250966396\\
911&21647&16553&128280463&42797&987290451\\
1021&260342&19013&116496427&46229&1742607656\\
1163&549066&19661&95133794&47309&941740647\\
1483&349678&20849&270934284&49937&5343544\\
\hline
\end{tabular}
\end{center}
Read entries down each pair of columns. The 51 distinct $m$ are pairwise
coprime. Append the two further rows $(m,k)=(827,170),(5717,2537)$ from Section 3. These moduli are also coprime to each other and to all 51 previous moduli. Over all 53 rows, put $N=\prod m$ and let $0\le h<N$ solve $h\equiv k\pmod m$.
The exponent lattice has initial basis $(N,0),(-h,1)$. Apply binary
Gauss reduction: swap $u,v$ if $\lVert v\rVert^2<\lVert u\rVert^2$,
put $q=\lfloor\langle u,v\rangle/\lVert u\rVert^2+1/2\rfloor$,
and replace $v$ by $v-qu$ until $q=0$. This deterministic convention fixes
the signs of the reduced coordinates $(i,j)$ in $(a,b)=iu+jv$.

\paragraph{The finite check in the Saito cone.}
Put $X=10^{-183}\log|\beta_1|$ and $Y=10^{-183}\log\beta$,
where $\beta_1$ is the first embedding and $\beta$ is the middle one.
For norm one, Saito's inequalities imply
$17Y/40\le-X\le23Y/40$. Indeed both small conjugates have absolute value
at most $\beta^{-17/40}$ and the sum of their negative logarithms is
$\log\beta$. For $(X,Y)^{\mathsf T}=W(i,j)^{\mathsf T}$ the exact
interval calculation yields
\[
 W\in\begin{pmatrix}
  (82.7293306,82.7293312)&(185.2873499,185.2873510)\\
  (107.5045790,107.5045794)&(-89.1471288,-89.1471282)
 \end{pmatrix}.
\]
Here inclusion is entrywise. If $\log_{10}\beta\le 2\cdot10^{186}$, then
$0<Y\le2000\log 10$. Bounding $W^{-1}$ in the same interval
arithmetic gives $|i|\le39$, $|j|\le24$. Intersect this box with the
two parity conditions, $a+b\equiv0\bmod5$, and the indicated closed
logarithmic strip. Only the four pairs below can survive. Each has an
independent zero of the trace modulo the displayed prime at index $r$,
and its complete powered residue repeats after $L$ steps from there.
\[
 \begin{array}{rr|rrr}
 i&j&p&r&L\\\hline
 8&-7&11&2&4\\18&-17&23&10&20\\22&-23&11&0&4\\24&-21&11&1&4
 \end{array}
\]
The two small conjugates are negative, so trace and floor agree eventually.
The floors on $r+L\mathbf Z_{\ge0}$ grow past the indicated divisor.
This finishes the claimed exclusion.

\paragraph{Reproducibility.}
Besides modular arithmetic and integer enumeration, the only significant
real-valued inputs are root isolation and logarithms. For each rational
endpoint $1\le x\le2$, put $t=(x-1)/(x+1)$. Summing
$2\sum_{k=0}^{17}t^{2k+1}/(2k+1)$ gives a lower bound for $\log x$;
adding $2t^{37}/(37(1-t^2))$ gives an upper bound. Powers of two reduce
any positive rational endpoint to this range. The certificate uses these
bounds throughout, with outward rounding to rational multiples of $10^{-8}$.


\paragraph{Why a norm-one prime tail forces the stated signs.}
Saito's larger small conjugate is negative. Since the norm is positive, the other small conjugate is negative too. Hence the floor eventually equals the trace. Write $\gamma=E^aT^b$. At $3$ there is one simple linear root of $f$, and its discriminant squareclass $61$ is a square in $\Q_3$. The other two roots are therefore also in $\Q_3$, and the three reductions of $\theta$ are $1,-1,-1$. At them $E$ reduces to $1,-1,-1$ and $T$ to $1,1,1$. If $a$ is even, the trace is always zero modulo $3$.

If $a$ is odd and $b$ even, then $\gamma$ is $E$ times a square. At $7$ the projective residue group is cyclic of order $48$ (one linear and one irreducible quadratic factor). The element $[E]$ has order $16$, as direct powering gives $E^8=(0,5,6)$ and $E^{16}=(2,0,0)$ in the power basis. Thus $[\gamma^3]$ is an odd power of $[E]$. Multiplication by $3$ on odd exponents modulo $16$ has just two cycles, containing $3$ and $5$ respectively. Both $E^3$ and $E^5$ have zero trace modulo $7$. This excludes the second case on a progression. Hence $a,b$ must both be odd, giving $(-,+,-)$.

\paragraph{A limitation of inert-prime tests for order 67.}
Let $p$ be inert in this $K$. A projective class of exact order $67$ in $\F_{p^3}^\times/\F_p^\times$ cannot have trace zero. Indeed for any representative $x$ set $y=x^{p-1}$, of order $67$. Since $67\mid p^2+p+1$, the residue $r$ of $p$ modulo $67$ is $29$ or $37$. If $x$ has trace zero then $1+y+y^{r+1}=0$. For both possible $r$ the integer resultant is
\[
  \operatorname{Res}(\Phi_{67},1+X+X^{r+1})
  =1235376017=(29\cdot37)^3 .
\]
One exact reproduction in SymPy is
\texttt{resultant(cyclotomic\_poly(67,x),1+x+x**(r+1),x)}.
Therefore $p$ is $29$ or $37$. But $f$ has the simple linear factor $X+2$ at $29$, and $X+9$ at $37$; these primes are unramified and prime to the power-basis index. Neither is inert. This observation concerns exact order $67$ at the inert prime itself. The order measured at $34303$ need not persist at another prime.



\sect{5. A degree consequence for the original minimum}
By [1, Theorem 1.7], if the least Mills constant is algebraic, one of its powers $M^{3^m}$ is a totally real cubic Pisot number of the specified conjugate shape. Suppose such a power belongs to the field $K$ above and has norm $+1$. Descend by real cube roots in $K$ while the level is positive and the cube root belongs to $K$. The resulting $\gamma=M^{3^t}$ still generates $K$: its cube, or an iterated cube, generates $K$. It is integral, its norm is still $1$, and its sign and size inequalities persist. The sign calculation on the preceding page leaves only $Z\unit^2$ on a prime tail. The exceptional mod-$61$ phase also persists: on the order-$20$ group remaining after cubing, the cube map is bijective and preserves the order-four subgroup. For clarity, in the zero-shift case the mod-$61$ test underlying this necessity is recalled below.

If $t>0$, $\gamma$ is not a cube in $K$. Then $X^{3^t}-\gamma$ is irreducible over $K$. Start with $X^3-\gamma$ and take successive real cube roots. At each stage the norm to $K$ of the preceding root is $\gamma$ (its binomial minimal polynomial has odd degree). A cube at that stage would have cube norm to $K$, a contradiction. Thus $[\Q(M):\Q]=3^{t+1}$, also for $t=0$.

An explicit upper bound for $M$ suffices to convert the lower bound on $\gamma$ to a lower bound on $t$. By Cully-Hugill [4] there is a prime between $v^3$ and $(v+1)^3$ for all integers $v\ge L=\exp(\exp(32.537))$. Bertrand's postulate supplies a prime $p_1$ between $\lceil L\rceil$ and $2\lceil L\rceil$. Iterate the cube-interval theorem from $p_1$ to produce a Mills-representing number $A$ with $A^3<p_1+1$. Indeed the intervals $[p_j^{1/3^j},(p_j+1)^{1/3^j}]$ are nested and their upper endpoints strictly decrease: the possible uppermost integer $(p_j+1)^3-1$ is composite. Hence
\[
 \log_{10}M\le\log_{10}A<\tfrac13\log_{10}(2\lceil L\rceil+1)
 <2\cdot10^{13}.
\]
For example the last inequality follows from $e^{32.537}<1.36\cdot10^{14}$ and $\log 10>2.3$. By Section 4, $2\cdot10^{186}<\log_{10}\gamma=3^t\log_{10}M$. As $2\cdot10^{13}3^{362}<2\cdot10^{186}$, this forces $t\ge363$ and gives Theorem B. No assertion about a numerical decimal expansion of the minimum enters this argument.

\textbf{The bad phase.} In the remaining sign class $A_j\to-1$ in $\Z_3$: at $3$ the three unit residues are $1,-1,-1$. At $61$ the roots are $r,r,r^{-2}$, and for $j\ge1$ put $z=(r^3)^{3^j}$ in the order-$20$ subgroup. The quadratic character of the trace is that of $1+2z$. The cubing cycles, listed by their least representatives, are $1,3,8,9,11,23,60$. Only $11\leftrightarrow50$ gives exclusively nonsquares. These facts can be checked by iterating at most 20 elements. If another cycle occurs, one fixed phase has square trace. If there are arbitrarily late prime traces $p$ on it, quadratic reciprocity gives even Frobenius at $p$ on the cubic roots. They split over $\F_{p^3}$, in which cubing is invertible since $p\equiv-1\bmod3$. A prime seed returns on a progression and later terms exceed it. Otherwise that phase is eventually composite. This is the criterion from corpus [009]. The $61$ phase of every hypothetical prime tail must therefore be the one used in Section 4.




\sect{6. Scope and verification}
Theorem A handles every member of the earlier specified ray family that satisfies the Pisot conditions. The conclusion also holds for other $91$st powers in the sign class. The surviving relative fraction in Section 3 refers to the full remaining sign class and the two specified auxiliary primes. It is not a probability for the least constant. Section 4 bounds the cubic power within one fixed field and norm. We have no upper bound on its level in the original constant. Other norms and other fields remain possible.

\textbf{G1--G6.} G1 identifies Saito's original Problem 1.1. G2 retains the base and the $3^j$ exponents, with an explicit subgroup restriction. G3 follows the original problem and Saito's later conditional result. G4 remains open for the original problem. Theorems A and B prove the stated restricted conclusions. The specific advance counted under G5 links the old small-prime congruences to a global power condition and thence to all six local cycles; the height argument combines both mixed-prime congruences with the 51 inert-prime congruences in one exponent lattice. G6, historical priority, has not been established, and the manuscript has not been externally reviewed.

\textbf{Sources.}
\begin{enumerate}\raggedright
\item K. Saito, \emph{Transcendency of variants of Mills' constant}, Theorems 1.7 and 1.8. \href{https://arxiv.org/abs/2508.16068}{arXiv:2508.16068}.
\item K. Saito, \emph{Intervals without primes near an iterated linear recurrence sequence}, Problem 1.1. \href{https://arxiv.org/abs/2504.14968}{arXiv:2504.14968}.
\item \emph{Composite Mills floors in the exceptional discriminant class 61}, preceding research note, Sections 3--5. Related methods in QEEDDD's corpus [009], [042], [075], and audit terminology in [135]. \href{https://github.com/QEEDDD/Verifiable-Mathematical-Reasoning-Corpus/tree/38348353416f1f7fc50cecb93c3b554a00260c13}{Repository snapshot}.
\item M. Cully-Hugill, \emph{Explicit estimates for the distribution of primes}, Bull. Aust. Math. Soc. 108 (2023), 341--342, reporting the explicit cube-interval threshold. \href{https://doi.org/10.1017/S0004972723000588}{doi:10.1017/S0004972723000588}.
\end{enumerate}

The appendices contain two standalone listings. In Appendix A, compile as C++17 with optimization and without disabling assertions. The program tests the complete table by visiting the exact-order residue classes, as well as the shorter six witnesses and the discrete logarithms used in the proof. Appendix B uses Python 3 and its standard library. Its optional \texttt{--full} argument also repeats the original 22 orbit checks. Both listings can be copied directly from the single \LaTeX{} source.



\sect{Appendix A. Mixed-prime cycle certificate}
\begin{Verbatim}[fontsize=\footnotesize,breaklines,breakanywhere]
#include <array>
#include <vector>
#include <numeric>
#include <cassert>
#include <cstdint>
using I = std::int64_t;
using A = std::array<I,3>;
int p;
const A one={1,0,0}, E={-56,37,2}, T={-46,27,2};
I mod(I a){return (a%p+p)%p;}
A mul(A a,A b) {
    I w=a[2]*b[2], v=a[1]*b[2]+a[2]*b[1]-19*w;
    return {mod(a[0]*b[0]+19*v),
            mod(a[0]*b[1]+a[1]*b[0]+19*w+19*v),
            mod(a[0]*b[2]+a[1]*b[1]+a[2]*b[0]-19*v+19*w)};
}
A pw(A a,I n) {
    A r=one;
    while(n){if(n%2)r=mul(r,a); a=mul(a,a); n/=2;}
    return r;
}
bool scalar(A a){return a[0]!=0 && a[1]==0 && a[2]==0;}
bool same(A a,A b){
    for(int i=0;i<3;i++)for(int j=0;j<i;j++)
        if(mod(a[i]*b[j]-a[j]*b[i]))return false;
    return true;
}
bool zero(A a){return mod(3*a[0]-19*a[1]+399*a[2])==0;}
bool prime(I n) {
    if(n<2)return false;
    for(I i=2;i<=n/i;i++) if(n%i==0)return false;
    return true;
}
std::vector<I> factors(I n){
    std::vector<I> result;
    for(I d=2;d<=n/d;d++)if(n%d==0){
        result.push_back(d);while(n%d==0)n/=d;
    }
    if(n>1)result.push_back(n);
    return result;
}
void order(A a,I n){
    assert(scalar(pw(a,n)));
    for(I q:factors(n))assert(!scalar(pw(a,n/q)));
}
I powmod(I a,I n,I q){
    I r=1;
    while(n){if(n%2)r=r*a%q; a=a*a%q; n/=2;}
    return r;
}
void evaluations(int q,int l,int g,std::array<int,3> roots,
                 std::array<int,4> logs){
    assert(prime(q)&&prime(l)&&g!=1&&powmod(g,l,q)==1);
    for(I x:roots)assert((x*x*x+19*x*x-19*x-19)%q==0);
    for(int i=0;i<2;i++){
        I x=roots[i];
        assert(powmod((2*x*x+37*x-56)%q,(q-1)/l,q)==powmod(g,logs[2*i],q));
        assert(powmod((2*x*x+27*x-46)%q,(q-1)/l,q)==powmod(g,logs[2*i+1],q));
    }
    assert((logs[0]*logs[3]-logs[1]*logs[2])%l);
}
void six(A g,I N,I m,I period,std::vector<I> reps){
    A h=pw(g,N/m);std::vector<bool> marks(m,false);
    for(I e:reps){
        assert(std::gcd(e,m)==1 && !marks[e] && zero(pw(h,e)));
        I r=e, count=0;
        do {assert(!marks[r]); marks[r]=true; count++;r=r*3%m;}while(r!=e);
        assert(count==period);
    }
    for(I e=1;e<m;e++)if(std::gcd(e,m)==1)assert(marks[e]);
}
void whole_row(A g,I N,I m,int good_count,int cycle_count){
    std::vector<bool> seen(m,false);
    A h=pw(g,N/m);
    int good=0,total=0;
    for(I e=1;e<m;e+=2){
        if(seen[e]||std::gcd(e,m)!=1)continue;
        I j=e; A x=pw(h,e); bool hit=false;
        do{
            seen[j]=true;
            if(!hit){hit=zero(x); if(!hit)x=mul(mul(x,x),x);}
            j=3*j%m;
        }while(j!=e);
        good+=hit;total++;
    }
    assert(good==good_count && total==cycle_count);
}
int main(){
    evaluations(1933,7,1812,{71,112,1731},{1,4,3,4});
    evaluations(937,13,911,{501,666,688},{9,10,9,1});
    for(int val:{21503,34303}){
        p=val;assert(prime(p)); int roots=0;
        for(I x=0;x<p;x++)if(mod(x*x*x+19*x*x-19*x-19)==0)roots++;
        assert(roots==1 && (61*114*114)%p);
        I S=I(p)*p-1, N=S/3;
        A g=pw({p==21503?2:1,1,0},3);
        order(g,N); order(pw(E,3),p==21503?N:N/2);
        order(pw(T,3),p==21503?N/2:N);
        std::vector<std::array<int,2>> logs;
        std::vector<std::array<int,3>> rows;
        if(p==21503){
            six(g,N,2048,512,{231,793});
            six(g,N,1693696,211456,{36681,106441,193855,598607});
            logs={{7,1},{13,12},{827,170}};
            rows={{1,2,2},{7,0,4},{13,2,8},{91,4,48},
                  {827,4,4},{5789,8,8},{10751,16,16},{75257,96,96}};
        }else{
            logs={{67,34},{5717,2537}};
            rows={{1,2,2},{67,0,12},{5717,8,8},{383039,48,48}};
        }
        for(auto [q,k]:logs){
            A x=pw(E,S/q),y=pw(T,S/q);
            assert(!scalar(x) && same(y,pw(x,k)));
        }
        for(auto [d,hit,total]:rows)
            whole_row(g,N,(p==21503?2048LL:1024LL)*d,hit,total);
    }
}
\end{Verbatim}


\sect{Appendix B. Exponent-lattice certificate}
\begin{Verbatim}[fontsize=\footnotesize,breaklines,breakanywhere]
"""Verify the D=61 norm-one exponent-lattice bound. Python 3, standard library.

The optional --full flag rechecks the earlier twenty-two exhaustive finite-field
certificates. This part traverses several million cubic finite-field elements.
"""
from fractions import Fraction as F
from math import gcd, isqrt
import sys

DATA = [
    (5,1),(41,787),(97,2951),(103,1915),
    (127,4419),(131,14136),(167,11077),(199,8942),
    (271,5781),(293,50847),(773,552873),(827,419023),
    (859,50622),(911,21647),(1021,260342),(1163,549066),
    (1483,349678),(1583,255229),(1709,2745692),(1723,654871),
    (1811,205927),(1993,630756),(3557,8372873),(4889,3215080),
    (8681,46418403),(10589,44659163),(13469,31295481),(16229,77599160),
    (16301,66708745),(16421,76985174),(16553,128280463),(19013,116496427),
    (19661,95133794),(20849,270934284),(22277,350529041),(22937,297631469),
    (26489,31896576),(26921,404448787),(28901,65818922),(31277,924860164),
    (32693,1008873991),(36497,695170719),(36749,1317741455),(37049,1219522376),
    (37649,1390873606),(38669,243832387),(38993,1250966396),(42797,987290451),
    (46229,1742607656),(47309,941740647),(49937,5343544),
]
ONE=(1,0,0)
E=(-56,37,2)
T=(-46,27,2)

def mul(a,b,p):
    c=[0]*5
    for i in range(3):
        for j in range(3):
            c[i+j]+=a[i]*b[j]
    for k in (4,3):
        c[k-1]-=19*c[k]; c[k-2]+=19*c[k]; c[k-3]+=19*c[k]
    return tuple(x%p for x in c[:3])

def power(a,n,p):
    z=ONE
    while n:
        if n%2: z=mul(z,a,p)
        a=mul(a,a,p); n//=2
    return z

def project(a,p):
    c=pow(next(x%p for x in a if x%p),-1,p)
    return tuple(x*c%p for x in a)

def prime(n):
    return n>=2 and all(n%d for d in range(2,isqrt(n)+1))

def factors(n):
    d=2
    while d*d<=n:
        if n%d==0:
            yield d
            while n%d==0: n//=d
        d+=1
    if n>1: yield n

def trac(a,p):
    return (3*a[0]-19*a[1]+399*a[2])%p

def lattice():
    N,h=1,0
    for n,(p,k) in enumerate(DATA):
        S=p*p+p+1; m=S//gcd(S,3)
        assert prime(p) and prime(m) and gcd(N,m)==1
        assert all((t**3+19*t*t-19*t-19)%p for t in range(p))
        assert project(E,p)!=ONE and project(power(E,m,p),p)==ONE
        assert project(power(E,k,p),p)==project(T,p)
        if n>=22:
            assert p%3==2
            assert all(pow(3,(m-1)//d,m)!=1 for d in factors(m-1))
        elif '--full' in sys.argv:
            seen=bytearray(m); x=ONE
            for i in range(m):
                if not trac(x,p):
                    j=i
                    while not seen[j]:
                        seen[j]=1; j=3*j%m
                x=mul(E,x,p)
            assert sum(seen)==m-1 and not seen[0]
        h+=N*((k-h)*pow(N,-1,m)%m); N*=m
    for m,k in [(827,170),(5717,2537)]:
        # Projective congruences and cycle coverage at 21503,34303 in Appendix A.
        assert gcd(N,m)==1
        h+=N*((k-h)*pow(N,-1,m)%m); N*=m
    def dp(a,b): return sum(x*y for x,y in zip(a,b))
    u,v=(N,0),(-h,1)
    while True:
        if dp(v,v)<dp(u,u): u,v=v,u
        d=dp(u,u); q=(2*dp(u,v)+d)//(2*d)
        if not q: break
        v=tuple(x-q*y for x,y in zip(v,u))
    assert abs(u[0]*v[1]-u[1]*v[0])==N
    return u,v

def normform(x,y,z):
    return (x**3-x*x*z-76*x*y*y-38*x*y*z-101*x*z*z
            -228*y**3-532*y*y*z+646*y*z*z+387*z**3)

def units():
    for p in (2,19):
        for x in range(p):
            for y in range(p):
                for z in range(p):
                    b=3*x*x-2*x*z-76*y*y-38*y*z-101*z*z
                    if (3*x-z)%p==b%(p*p)==normform(x,y,z)%(p**3)==0:
                        assert (x,y,z)==(0,0,0)
    solutions=set()
    for y in range(-26,27):
        for z in range(-34,35):
            q=-76*y*y-38*y*z-101*z*z; r=normform(0,y,z)
            for x in range(-230,231):
                if ((x-z)*x+q)*x+r==1: solutions.add((x,y,z))
    assert solutions=={(1,0,0),(-24,5,1),(46,15,-9),(-132,-22,19)}

def interval(a,b=None):
    if b is None:b=a
    return (F(a),F(b))
def add(a,b):return a[0]+b[0],a[1]+b[1]
def neg(a):return -a[1],-a[0]
def prod(a,b):
    s=[x*y for x in a for y in b]
    return min(s),max(s)
def inv(a):
    assert a[0]*a[1]>0
    return 1/a[1],1/a[0]
def sub(a,b):return add(a,neg(b))
def ab(a):
    assert a[0]*a[1]>0
    return a if a[0]>0 else neg(a)
def dot(a,b):return add(prod(a[0],b[0]),prod(a[1],b[1]))
def inv2(M):
    a,b=M[0];c,d=M[1];D=inv(sub(prod(a,d),prod(b,c)))
    return [[prod(t,D) for t in R] for R in [[d,neg(b)],[neg(c),a]]]
def sqrt_ge(a,x):return a>=x*x
def natlog_one_two(x):
    t=(x-1)/(x+1)
    t=F(t); t2=t*t; q=t; s=0;n=18
    for k in range(n):s+=2*q/(2*k+1);q*=t2
    return (s,s+2*q/((2*n+1)*(1-t2)))
logtwo=natlog_one_two(F(2))
def logpoint(x):
    assert x>0;k=0
    while x<1:x*=2;k-=1
    while x>2:x/=2;k+=1
    return add(natlog_one_two(x),prod(interval(k),logtwo))
def logiv(a):
    lo,hi=logpoint(a[0])[0],logpoint(a[1])[1]
    f=10**8
    return interval(F((lo*f).__floor__(),f),F((hi*f).__ceil__(),f))
def box(n):return interval(F(n,10**12),F(n+1,10**12))
theta=[box(t) for t in [-19906514237356,-623731016520,1530245253875]]
def poly(a,coeffs):
    v=interval(0)
    for c in coeffs:v=add(prod(v,a),interval(c))
    return v
for a,b in theta:
    def f(x):return x**3+19*x*x-19*x-19
    assert f(a)*f(b)<0
es=[poly(t,[2,37,-56]) for t in theta]
ts=[poly(t,[2,27,-46]) for t in theta]
logs=[[logiv(ab(e)),logiv(ab(t))] for e,t in zip(es,ts)]
G=[[sum((prod(R[i],R[j])[0] for R in logs),F()),sum((prod(R[i],R[j])[1] for R in logs),F())] for i,j in [(0,0),(1,1),(0,1)]]
g,h,j=G
assert g[0]>47 and (g[0]-47)*(h[0]-47)>max(j[0]**2,j[1]**2)
def box_certificate():
    bounds=(300,71,268)
    for i in range(3):
        e,t=ab(es[i]),ab(ts[i])
        assert max(e[1],1/e[0])*max(t[1],1/t[0]) < bounds[i]**2
    R=[ [interval(1),poly(t,[F(1,3),6,F(-19,3)]),
          poly(t,[F(1,3),7,F(-1,3)])] for t in theta]
    def cofactor(i,j):
        u,v=[k for k in range(3) if k!=i]
        r,s=[k for k in range(3) if k!=j]
        det=sub(prod(R[u][r],R[v][s]),prod(R[u][s],R[v][r]))
        return neg(det) if (i+j)%2 else det
    D=interval(0)
    for j in range(3):D=add(D,prod(R[0][j],cofactor(0,j)))
    for i,maximum in enumerate((230,26,34)):
        assert sum(bounds[j]*ab(prod(inv(D),cofactor(j,i)))[1]
                   for j in range(3))<maximum
    units()

def cone_certificate():
    U,V=lattice(); scale=10**183; l10=logiv(interval(10))
    W=[[dot(r,[interval(F(x,scale)) for x in vec]) for vec in (U,V)]
       for r in logs[:2]]
    bounds=[int((ab(r[0])[1]*F(23,40)+ab(r[1])[1])*2000*l10[1])
            for r in inv2(W)]
    assert bounds==[39,24]
    candidates={}
    for i in range(-39,40):
        for j in range(-24,25):
            a,b=(i*U[t]+j*V[t] for t in (0,1))
            if not (a%2 and b%2 and (a+b)%5==0):continue
            X,Y=(dot(r,[interval(i),interval(j)]) for r in W)
            if Y[1]<=0 or Y[0]>2000*l10[1]:continue
            if sub(prod(interval(-40),X),prod(interval(17),Y))[1]<0:continue
            if add(prod(interval(23),Y),prod(interval(40),X))[1]<0:continue
            candidates[i,j]=(a,b)
    table={(8,-7):(11,2,4),(18,-17):(23,10,20),
           (22,-23):(11,0,4),(24,-21):(11,1,4)}
    assert candidates.keys()==table.keys()
    for ij,(a,b) in candidates.items():
        p,r,h=table[ij]
        # E and T have order dividing p^2-1: splitting type (1,2) or (1,1,1).
        d=p*p-1
        assert power(E,d,p)==power(T,d,p)==ONE
        z=mul(power(E,a%d,p),power(T,b%d,p),p)
        y=power(z,3**r,p)
        assert trac(y,p)==0 and power(y,3**h,p)==y

box_certificate()
cone_certificate()
print("Verified log10(beta) > 2*10^186.")
\end{Verbatim}


% Chapter early
\chapter{Early composite floors and recurring divisors}\label{ch:early}\markboth{Early composite floors}{}
Let $M$ be the least real number greater than one for which every $\lfloor M^{3^n}\rfloor$, $n\ge1$, is prime. Saito's theorem [1, Thm.~1.7] implies that if $M$ is algebraic, a power $M^{3^m}$ is a totally real cubic Pisot number. This note concerns the possibility that such a power belongs to
\[
 K=\Q(\theta),\qquad \theta^3+19\theta^2-19\theta-19=0,
\]
and has norm $+1$. Set $E=2\theta^2+37\theta-56$ and $T=2\theta^2+27\theta-46$. These are a basis of the norm-one units [3]. Order the embeddings by increasing $\theta$. The only surviving sign class is $(-,+,-)$, with $a,b$ odd for $\beta=E^aT^b$. Its middle embedding is the real number $\beta>1$. Write $A_j=\Tr(\beta^{3^j})$ and $F_j=\lfloor\beta^{3^j}\rfloor$.

\sect{1. Two primes test the first and the third floor}
Take the least $t\ge0$ for which the Saito power has descended to $\beta=M^{3^t}\in K$. Cubing preserves its signs and the inequalities on absolute values, so those inequalities hold after descent as well. The preceding bound $[\Q(M):\Q]\ge 3^{364}$ in this case implies $t>0$ [3]. Thus $\beta$ is not a cube in $K$, equivalently $a,b$ are not both divisible by three. Both small conjugates are bounded in absolute value by $\beta^{-17/40}$. Since the lower bound on $\beta$ in [3] exceeds $2^{40/17}$, their negative odd powers have sum strictly between $-1$ and $0$ for every $j\ge0$. Hence $A_j=F_j$ from $j=0$ on. In particular $F_0$ and $F_2$ must be primes exceeding $647$.

For $p=59,647$ write $[\beta]_p$ for the reduction modulo $p$, up to multiplication by a common nonzero scalar. Its order is computed in the invertible elements of the residue algebra modulo $\F_p^\times$. The prime-to-three component is trivial precisely when this order is a power of three, allowing order one.

\para{Theorem 1.} If $3\nmid b$ and $[\beta]_{59}$ has order a power of three, then $59\mid F_0$. If $3\mid b$ and $[\beta]_{647}$ has order a power of three, then $647\mid F_2$. In particular, no canonical Mills candidate has order a power of three at both primes.

\para{Proof.} Both primes have splitting type $(1,2)$, as the discriminant has nonsquare class $61$ modulo either prime. The projective group is cyclic of order $s=p^2-1$. Raise the units to $c$, where $c$ is $1$ modulo the three-part of $s$ and $0$ modulo its coprime part. For $(p,c)=(59,2320)$ the projected units are $(1,g)$, with $g$ of order $3$ and $g,g^2$ both of trace zero. For $(p,c)=(647,25840)$ the projected units are $(h,h^4)$, with $h$ of order $27$ and $h^9,h^{18}$ both of trace zero. These identities are polynomial multiplications modulo $p$. In the first case the exponent is $b$. In the second, $3\mid b$ forces $3\nmid a+4b$. The ninth power therefore has trace zero. Multiplying by a scalar preserves zero trace. In exponent coordinates the two necessary escape conditions read $a\not\equiv903b\pmod{1160}$ if $3\nmid b$, and $b\not\equiv551a\pmod{2584}$ if $3\mid b$. The certificate verifies the discrete logarithms. \hfill$\square$


\sect{2. An infinite family surviving 554 primes}
The following construction specifies an infinite family for which a sieve using only eventual prime divisors would encounter a genuine obstruction. Theorem 1 excludes every normalized member of this family as a power of the least Mills constant. Section 3 supplies a recurring composite progression even without normalizing the power.

Let $\mathcal P$ be the union of all primes at most $3000$, the 51 inert primes in [3], the 94 inert primes of Section 4 and $21503,34303,524287$. The union has 554 distinct members. The certificate specifies their values. Put $P=\prod_{p\in\mathcal P}p$. We construct an odd integer $L$ with $v_3(L)=13$. The unit
\[
                 X=(ET^{23})^L
\]
has trace coprime to $P$ after every iterated cubing, including the initial state.

Here is the construction of $L$. Let $\mathcal P_0$ consist of the small primes, the 51 earlier inert primes, the three mixed primes, the 35 rows with small factor $7$ or $13$ and the five $67$-rows with slope $34$. There are 502 primes in the union. For $B_p=p(p^2-1)(p^3-1)$ put $H=\operatorname{lcm}_{p\in\mathcal P_0} B_p$. Strip all factors $2,7,67$ from $H$ to obtain $L_0$ and write $R=H/L_0$. For each of the remaining 52 primes let $q_p=(p^2+p+1)/(67\gcd(3,p^2+p+1))$ and set $F=\prod q_p$. Define $L=L_0 F(F^{-1}\bmod R)$ using the least positive inverse. The certificate checks that the unit $ET^{23}$ has order dividing $B_p$ at every prime used here. It checks a complete projective cubing period starting at $X$ at each prime. There are $796543$ states in total, and the longest period is $786432$. Every trace is nonzero.

\para{The normalized family.} Put $u=L/3^{13}$, $X_0=(ET^{23})^u$ and let $\epsilon$ be any totally positive unit equal to $1$ modulo $P\OO$. Then
\[
                         \beta=X_0\epsilon^3
\]
is a noncube in $K$. Indeed its unit exponents are $u+3v,23u+3w$, with $3\nmid u$. At each of the 554 primes its $3^{13}$th power agrees with $X$. Thus every $\Tr(\beta^{3^j})$, $j\ge13$, avoids all 554 primes. The exponents at the two primes of Theorem 1 also have scalar prime-to-three component, as verified directly for $X$; taking a $3^{13}$th root leaves this property intact. In this displayed family $3\nmid b$, so $59$ already divides $F_0$ for its large Pisot members. The $647$ alternative in Theorem 1 also handles noncubes with $3\mid b$ in the more general local class.

\para{Why the two early tests work.} At any mixed prime $p\equiv2\bmod3$, a projective class of order $3$ can be represented by $(1,\omega)$ in $\F_p\times\F_{p^2}$, with $\omega$ a primitive cube root of unity. Both its trace and the trace of its inverse vanish, since $1+\omega+\omega^p=0$. Thus a projective order $3^e>1$ gives a zero after $e-1$ cubings. In the present Saito range the second coefficient $B_j=\Tr(\beta^{-3^j})$ is negative and $0<|B_j|<F_j$ for sufficiently large $\beta$ at every $j$. Indeed for $x=\beta^{3^j}$, the two small conjugate moduli have product $1/x$ and are at most $x^{-17/40}$; hence $B_j=x^{-1}-x(|\beta_2|^{3^j}+|\beta_3|^{3^j})$ and $|B_j|<2x^{23/40}$. Any prime dividing both coefficients is strictly smaller than $F_j$. This also handles auxiliary primes greater than $\beta$.

\para{Existence and size.} Totally positive units equal to $1$ modulo $P\OO$ have finite index in the rank-two unit group. Their log vectors, multiplied by three, still form a full lattice in the plane with coordinate sum zero. After translation by the log vector of $X_0$, bounded-distance approximations to $t(-0.46,1,-0.54)$ give arbitrarily large Pisot numbers in the strict Saito range. Their signs are those of $ET$. Consequently this construction excludes infinitely many genuine, arbitrarily large noncube Pisot candidates while witnessing avoidance of the 554 primes on the tail.


\sect{3. A recurring divisor outside the 554 primes}
\para{Theorem 2.} Let $\beta$ be a Pisot unit of signature $(-,+,-)$. Suppose that for some $h\ge0$ the residues of $\beta^{3^h}/X$ have projective orders which are powers of three at each of the six detecting primes in the table below. Every sufficiently late block of $16464$ indices contains a $j$ with $4801\mid F_j$. In particular this holds throughout the family of Section 2 with $h=13$.

\para{Proof.} Write $\beta^{3^h}/X=E^vT^w$. Each table entry gives the order $o$ of a generator of the indicated primary part of the subgroup generated by the two units, followed by the logarithms $(e,t)$ of $E,T$ after raising them to the prime-to-primary part of the projective group exponent. For the respective primes this forces $ev+tw\equiv0\bmod o$.
\[
\begin{array}{r|r|r|rr}
\text{required modulus}&p&o&e&t\\\midrule
128&21503&2048&1&1598\\
&34303&1024&438&1\\
25&1249&625&140&1\\
&2999&125&1&105\\
343&1733&343&1&331\\
&1373&343&1&292
\end{array}
\]
Each pair of logarithm vectors is independent modulo the prime below the required modulus. Thus $v,w$ vanish modulo $128,25,343$. Now $4801$ is prime, with splitting type $(1,2)$ and projective group order $128\cdot3\cdot25\cdot2401$. Both units have trivial $3$-component there. The exponents of $\beta^{3^h}$ equal $(25,63)$ modulo $128$, $(0,0)$ modulo $25$, and are of the form $(2364+343r,1550+343s)$ modulo $2401$, $0\le r,s<7$.

\[
\begin{array}{c|rrrrrrr}
r+s\bmod7&0&1&2&3&4&5&6\\\midrule
j_0&802&3154&5506&1159&1373&771&3123
\end{array}
\]
For all 49 lifts, powering gives zero trace at $j_0$ and projective return after $16464$ cubings. Here $j_0$ is counted from $\beta^{3^h}$, so the index for $\beta$ is $h+j_0$ modulo $16464$. The eventual equality $F_j=A_j$ proves the claim. \hfill$\square$

The two mechanisms have different consequences. Theorem 1 tests specified early floors of a minimum-level power. Theorem 2 yields infinitely many composite floors for the family without using minimality. Both apply beyond every height cutoff.


\sect{4. Exact local conditions at 94 inert primes}
Write $S_p=p^2+p+1$ and $N_p=S_p/\gcd(3,S_p)$. At an inert prime the projective image of the norm-one residues is cyclic of order $N_p$. This follows by dividing the kernel of the finite-field norm by its scalar subgroup. The following rows all have $N_p=\ell q_p$ with distinct primes $\ell,q_p$, and $[E]$ generates this group.

\para{Theorem 3.} At each of the 94 certified inert primes, a unit $\beta=E^aT^b$ avoids a recurring trace divisor $p$ exactly when $[\beta]_p$ has order dividing the indicated $\ell$. Equivalently $a+k_pb\equiv0\bmod q_p$, where $k_p$ is specified in the certificate.
\[
\begin{array}{c|rrr}
\ell&7&13&67\\\midrule
\text{number of primes }p&29&6&59
\end{array}
\]
For example, among the $\ell=67$ rows, the following five have the same linear condition modulo $67$ as the exact order test at $34303$ in [3].
\[
\begin{array}{r|r|r}
p&q_p&k_p\\\midrule
720413&7746203149&3475884807\\
989761&4873770283&1806255930\\
1118267&18664510471&13699147176\\
1253071&7811881513&1721835122\\
1672751&41762650459&16677457130
\end{array}
\]
\para{Proof.} The finite certificate tests inertness, the stated orders and $[T^\ell]=[E^{\ell k_p}]$. Its witnesses cover every cubing orbit of exact order $q_p$ or $\ell q_p$. For order $q_p$, the orbit label of the exponent $u$ is $u^{\operatorname{ord}_{q_p}(3)}\bmod q_p$. For order $\ell q_p$ write $u\bmod\ell=c3^j$ with one fixed $c$ per orbit and $0\le j<d_\ell=\operatorname{ord}_\ell(3)$. Set $d=\operatorname{ord}_{q_p}(3)/\gcd(\operatorname{ord}_{q_p}(3),d_\ell)$. The two labels $(c,(u3^{-j})^d\bmod q_p)$ then classify the orbits. The checker verifies each witness by modular powering and compares distinct label counts with the exact group cardinalities. For orders dividing $\ell$ it tests every element and finds none of trace zero. All these orders are coprime to three, so every hit recurs. \hfill$\square$

\para{At least 56 exact orders.} The 59 slopes modulo $67$ occupy 34 distinct lines. One slope, $34$, occurs five times. Every other slope occurs at most three times. A surviving candidate must have $a+34b\not\equiv0\bmod67$ by [3]. Hence $(a,b)$ is nonzero modulo $67$ and at most one distinct slope can vanish. Thus at least 56 of the 59 projective orders are exactly $67$. At all 59, $\beta^{67}$ reduces to a scalar.


\sect{5. A lower bound for the degree of the original constant}
\para{Theorem 4.} If the least Mills constant is algebraic and has a Saito power in $K$ of norm $+1$, then its descended cubic power and its degree satisfy
\[
       \log_{10}\beta>10^{563},\qquad [\Q(M):\Q]\ge3^{1154}.
\]

The preceding 51 inert-prime congruences, two mixed-prime congruences [3], and the 94 congruences of Theorem 3 have pairwise coprime moduli. Let their product be $N$. In each row $a+kb\equiv0\bmod m$. Chinese remaindering gives a single $h\bmod N$, and the lattice has basis $(N,0),(-h,1)$. Two-dimensional Gauss reduction gives a shortest squared length $S>2073\cdot10^{1122}$. The certificate checks the determinant and both reduction inequalities. Indeed if $U,V$ are reduced, the length of the component of $V$ perpendicular to $U$ is at least $\sqrt3\,\lVert U\rVert/2$. This handles coefficients of $V$ of magnitude at least two; for $0,\pm1$ the reduced inner product gives the bound directly.

The Gram matrix $G$ of the two unit log vectors satisfies $G>47I$, verified with rational root isolation and logarithm intervals. Thus the squared length of the log vector of a nonidentity unit in the lattice exceeds $47S$. In the Saito interval the two small logs are $-x\log\beta$ and $-(1-x)\log\beta$, with $17/40\le x\le23/40$. Hence
\[
 47S < \left(1+(17/40)^2+(23/40)^2\right)(\log\beta)^2 .
\]
The first inequality in Theorem 4 follows by exact interval arithmetic.

For the second, the explicit bound for primes between consecutive cubes [4] and Bertrand's postulate give $\log_{10}M<2\cdot10^{13}$, as in [3]. In detail, take $H=\exp(\exp(32.537))$, start at a prime between $\lceil H\rceil$ and $2\lceil H\rceil$, and successively choose primes strictly between $p^3$ and $(p+1)^3$. Nested intervals give a Mills-representing number at most $(2\lceil H\rceil+1)^{1/3}$. Since $2\cdot10^{13}3^{1152}<10^{563}$, it follows that $t\ge1153$. By minimal descent $\beta$ is not a cube in $K$. The binomial $X^{3^t}-\beta$ is irreducible over $K$: at each stage a cube root of the previous root cannot already belong to its current field, since its norm to $K$ would make $\beta$ a cube in $K$. Inductively the norm of the previous root is $\beta$ by its odd-degree binomial minimal polynomial. Since $\Q(\beta)=K$, this proves Theorem 4.

\para{Precisely what an infinite extension would give.} If infinitely many inert primes force $[\beta]_p$ to have order dividing $67$ for every surviving candidate, then the entire present field and norm case is excluded. For any fixed candidate $\beta^{67}$ would be a scalar at infinitely many primes. In an integral basis containing $1$, its other two integer coordinates would then both be zero. A positive power of a cubic Pisot number cannot be rational, since its other two conjugates have magnitude less than one. We have certified 59 such primes. Neither their infinitude nor an obstruction covering the unrestricted cases follows from the finite certificates.


\sect{6. Reproducibility and scope}
All prime-indexed tests concern exact modular residues. No primality experiment on a giant decimal approximation or estimate of a random survival probability is used. The two-prime early obstruction requires a normalized noncube and equality of the trace and floor at the indicated early indices. The recurrent-divisor theorem is specific to the displayed ray or its six-prime enlargement. The degree theorem assumes the specified field and norm. Other fields, other norms and the unrestricted transcendence of the least Mills constant remain open in this analysis.

\para{How to reproduce.} The certificate for this chapter is bundled with the other listings. After extracting the combined driver as described in the verification guide, run \texttt{python3 certificates.py run 06 --thorough}.
Python 3.9 or later is required. Run with assertions enabled. The full option rechecks the 22 inherited exhaustive finite-field rows. Both runs verify the 94 new rows, the log Gram bound, the 554-prime construction, the 49 lifts at $4801$ and the two early trace zeros. The previous report [3] supplies the proofs and certificates for the two mixed-prime restrictions and the norm-one sign classification. The source is a continuation of that report.

\para{G1--G6.} The original target is identifiable in Saito's Problem 1.1, with a sufficient reduction for the minimum supplied by Saito's Theorem 1.7 (G1--G3). The unconditional results above close the specified subfamilies. A complete proof of the unrestricted target required by G4 is still missing. The G5 argument in this note uses the noncube at the first cubic level to expose transient trace divisors, together with matching residue characters and a shortest-vector bound. Historical priority under G6 has not been established, and the note has not received external review.

\medskip\textbf{References}\par
{\small\raggedright
[1] K. Saito, \emph{Transcendency of variants of Mills' constant}, Theorems 1.7--1.8. \href{https://arxiv.org/abs/2508.16068}{arXiv:2508.16068}.\par
[2] K. Saito, \emph{Intervals without primes near an iterated linear recurrence sequence}, Problem 1.1. \href{https://arxiv.org/abs/2504.14968}{arXiv:2504.14968}.\par
[3] \emph{A recurring divisor beyond the small-prime barrier for cubic Mills floors}, preceding research note, 13 November 2026. Unit basis, 51 inert primes, and mixed-prime certificates. Related corpus arguments in [009], [042], [075], with G1--G6 in [135]. \href{https://github.com/QEEDDD/Verifiable-Mathematical-Reasoning-Corpus}{Verifiable Mathematical Reasoning Corpus}.\par
[4] M. Cully-Hugill, \emph{Explicit estimates for the distribution of primes}, Bull. Aust. Math. Soc. 108 (2023), 341--342. \href{https://doi.org/10.1017/S0004972723000588}{doi:10.1017/S0004972723000588}.\par}


% Chapter frob
\chapter{Frobenius orbits and the minimum shift order}\label{ch:frob}\markboth{Frobenius orbits and shift order}{}
Let $M$ denote the least Mills constant. Saito's Theorem 1.7 [1] forces a totally real cubic Pisot power $\beta=M^{3^m}$ if $M$ is algebraic. It has conjugates $\beta>1>|\beta_2|>|\beta_3|$, with $\beta_2<0$. Consequently its integer floors $F_j=\lfloor\beta^{3^j}\rfloor$ agree with $\Tr(\beta^{3^j})$ for all sufficiently large $j$. Each candidate must have prime floors at every index. We derive a local criterion allowing all norms, sharpen an existing fixed-field bound, and classify the order of polynomial identities among the floors.

\head{1. A trace-zero criterion for arbitrary norms}
Let $K=\Q(\beta)$, $p>3$ inert in $K$ and prime to the norm of $\beta$. Put
\[
 S=p^2+p+1,\qquad h=\gcd(3,S),\qquad n=S/h.
\]
Write $G=\F_{p^3}^{\times}/\F_p^\times$, and $H=G^3$, of index $h\le3$ in $G$. A scalar is the identity class of $G$. Suppose that
\[
              n\text{ is prime},\qquad
          \langle3,p\rangle=(\Z/n\Z)^\times.                      \tag{1}
\]
Then either $\beta^3$ is scalar modulo $p$, or $p\mid F_j$ on an infinite arithmetic progression of sufficiently late indices. The assertion does not require $\beta$ to be a unit.

\textbf{Proof.} The trace-zero classes form a set $D$ of $p+1$ projective points. If $g\ne1$ in $G$, the planes defining $D$ and $gD$ are distinct and intersect in a single projective point. Equality would give an invariant plane for multiplication by a nonscalar element of $\F_{p^3}$, whose minimal polynomial over $\F_p$ has degree three. It follows for every nontrivial multiplicative character $\chi$ of this cyclic group that
\[
 \left|\sum_{z\in D}\chi(z)\right|^2
 =(p+1)+\sum_{g\ne1}\chi(g)=p.
\]
Averaging the $h$ characters trivial on $H$ gives
\[
              |D\cap H|\ \ge\ (p+1-(h-1)\sqrt p)/h>0.
\]
No scalar has zero trace, since $p>3$. All nonidentity points of $H$ have order $n$. Cubing and Frobenius on $H$ multiply exponents by $3$ and $p$, respectively. By (1), from any nonidentity $[\beta^3]$ a mixed power $3^j p^k$ takes its class into $D$. Frobenius preserves $D$ in both directions, so cubing alone hits $D$. A period gives infinitely many returns. The floors grow past $p$. \hfill$\square$


\head{2. From local obstructions to a global criterion}
For a fixed $p$ with $n$ prime call the prime \emph{covering} if every orbit of cubing and Frobenius on $H\setminus\{1\}$ meets $D$. This depends only on $p$. Condition (1) suffices, but separate witnesses for each orbit also suffice. If every candidate field has infinitely many inert covering primes, the least Mills constant is transcendental. For any fixed candidate $\beta$, its norm has finitely many prime divisors, and $\beta^3$ can be scalar at only finitely many primes: in an integral basis containing $1$, some nonconstant coordinate is a nonzero integer, since the distinct conjugate sizes imply $\beta^3\notin\Q$. Thus even a single inert covering prime outside those exceptions rules out this candidate. Infinitude of suitable primes is open in this argument. The stronger version restricted to (1) encounters a quadratic obstruction, demonstrated in Section 5.

\textbf{Actual compression in the discriminant-61 example.} In the previous note [2] we used $K=\Q(\theta)$ with $\theta^3+19\theta^2-19\theta-19=0$, and the norm-one basis
\[
          E=2\theta^2+37\theta-56,\qquad T=2\theta^2+27\theta-46 .
\]
At $p=11299$, the prime $n=42559567$ satisfies
\[
 \ord_n(3)=14186522=(n-1)/3,\qquad \ord_n(p)=3,\qquad
                    3\nmid 14186522 .
\]
Thus (1) holds. For the norm-one projective image $H$, the certificate gives $[T]=[E]^{25827896}$ with $[E]$ a generator. Direct powering even supplies the point $E^{7098}=(4682,6455,1235)$ in the power basis modulo $p$; its trace in these representatives is $384166=34p$. Its Frobenius images meet all three cubing cycles. Every surviving norm-one candidate $E^aT^b$ therefore satisfies
\[
                     a+25827896 b\equiv0\pmod{42559567}.          \tag{2}
\]
For the ten primes below we test the cosets of $\langle3,p\rangle$ individually when that subgroup is smaller. Here its order is $3o$ with $o=\ord_n(3)$, and $w$ is the number of verified zero witnesses. Distinct $e^{3o}\bmod n$ label distinct cosets. All witnesses and discrete logarithms are supplied in the source.
\[
\begin{array}{r|r|r|r@{\quad}r|r|r|r}
p& n& k & w & p & n& k&w\\\midrule
3491&12190573&5409730&6&16267&88210519&49352284&11\\
6389&40825711&3297004&3&18427&113190919&3023671&1\\
10781&116240743&7713735&3&18773&352444303&347944980&3\\
11299&42559567&25827896&1&19237&123360469&17606603&2\\
14593&70990081&62632233&4&22643&512728093&273688730&6
\end{array}
\]
Each row gives $a+kb\equiv0\bmod n$. Their moduli are pairwise coprime and coprime to the preceding 147. Two-dimensional Gauss reduction in the resulting lattice gives a shortest squared length $S_0$. Using the previously certified log Gram bound $G_{\log}>47I$ and Saito's small logarithm interval [2] gives
\[
 \log_{10}\beta >
 \frac{1}{\log 10}\sqrt{\frac{47S_0}{1+(17/40)^2+(23/40)^2}}
 >5.14\cdot10^{602}.
\]
The earlier bounds $\log_{10}M<2\cdot10^{13}$ and $[\Q(M):\Q]=3^{t+1}$ at the least level $t$ then give $[\Q(M):\Q]\ge3^{1237}$ in this field and norm case.


\head{3. The precise shift order on the Mills schedule}
For any Saito candidate let $c=\mathrm N(\beta)$ and $A_j=\Tr(\beta^{3^j})$. Define the order of a polynomial identity as the least $r$ for which a nonzero $P\in\C(j)[X_0,\ldots,X_r]$ satisfies $P(A_j,\ldots,A_{j+r})=0$ eventually  (coefficients are evaluated at $j$). Then the minimum order is exactly two for $|c|=1$, and exactly three for $|c|>1$. For linear identities no finite order works in either case. These assertions also hold for $F_j$ by eventual equality.

\textbf{Proof.} Write $x_i=\beta_i^{3^j}$, $\beta_1=\beta$, and $s_n=\sum x_i^n$. If $|c|=1$, the product of the three $x_i$ is $c$. The pair of elementary symmetric functions $(s_1,b)$ is algebraically independent on the product-$c$ torus: for any prescribed complex sum and pair sum, the cubic with those coefficients and nonzero product $c$ has three nonzero complex roots. The pair $(s_1,s_3)$ is also independent since $s_3=s_1^3-3s_1 b+3c$ recovers $b$ away from $s_1=0$. Three functions on a two-dimensional torus are algebraically dependent.

For the nonunit case work on the full three-dimensional torus. The polynomials $s_1,s_3,s_9$ are independent since their Jacobian is nonzero at $(1,2,3)$, where the determinant after dividing rows by $1,3,9$ is $17640$. Four functions on this torus are dependent. These arguments give identities of the claimed upper orders.

To transfer independence to a single orbit, normalize a multiplicative relation among the absolute values to integer exponents. Such relations form a rational subspace invariant under a cyclic permutation of the three embeddings, since the square of a product equal to $\pm1$ equals $1$ and can be conjugated. The rational permutation space splits into the invariant line $(1,1,1)$ and the irreducible plane of coordinate sum zero, with different minimal polynomials for the cycle. The plane cannot be in the relation space since the absolute values are distinct. The invariant line lies in it precisely in the unit case. Consequently Laurent monomials have distinct growth rates on the orbit after eliminating $x_3=c/(x_1x_2)$ in the unit case, or directly in three variables otherwise. A nonzero Laurent polynomial with coefficients polynomial in $j$ has a unique dominant monomial for large $j$; the exponential scale is $3^j$. It cannot vanish eventually. This proves the lower orders even with rational-function coefficients. For a linear recurrence the largest shifted $A_{j+r}\sim\beta^{3^{j+r}}$ dominates every preceding shift times any rational function of $j$. \hfill$\square$

For $c=1$, setting $x=A_j$, $y=A_{j+1}$ gives the explicit order-two relation
\[
 b=(x^3+3-y)/(3x),\qquad A_{j+2}=y^3+3-3y(b^3-3xb+3).
\]
Its validity for prime and composite terms alike leaves the primality question open.


\head{4. A residue family with an early composite floor}
In the fixed field and norm-one case take the least level $t$ with $\beta=M^{3^t}\in K$. The preceding bound implies $t>0$, so $\beta=E^aT^b$ is not a cube. At least one exponent is prime to three. The large height and Saito inequalities in [2] give $F_j=A_j$ already at $j=0$, and $F_0>2789$.

If $a\equiv1585643b\bmod3457120$, then $17\mid F_1$ when $3\nmid b$, and $2789\mid F_0$ when $3\mid b$. Indeed the stated congruence implies $a\equiv11b\bmod32$ and $a\equiv289223b\bmod432140$. Both primes are mixed. At $17$, $g=[\theta+4]$ has order $288$ and $[E],[T]$ have exponents $45,145$ relative to $g$. The unit exponent $45a+145b$ is zero modulo $32$ and equals $b$ modulo $9$. If $3\nmid b$, one cubing gives a class of order three. At $2789$, $g=[(\theta+4)^6]$ has order $1296420=3\cdot432140$. The exponents of $[E],[T]$ are $318371,175467$. If $3\mid b$, then $3\nmid a$, so the combined exponent is divisible by $432140$ and nonzero modulo $3$. In each case a projective element of order three in $\F_p^\times\times\F_{p^2}^\times$ modulo common scalars has trace zero: its normalized coordinates are $(1,\omega)$ and $\omega^p=\omega^2$. This tests the original floors without an upper cutoff. Both primes themselves were included in the previous small-prime set; the conclusion here uses transient indices in an explicitly described congruence family.


\head{5. A missing squareclass and a two-orbit repair}
Assume $n=S/h$ prime. Then $n\equiv1\bmod3$ and $p$ has order three modulo $n$ (since $n>p+1$ for $p>3$). In particular $p$ is a square modulo $n$. Condition (1) requires $3$ to be a nonsquare modulo $n$. For $p\equiv1,5,7,11\bmod12$, respectively, $n\equiv1,7,7,1\bmod12$. Hence (1) forces $p\equiv 5,7\bmod12$, or $\left(\frac3p\right)=-1$. In a non-Galois cubic field with discriminant squareclass $3$, every unramified inert prime has even Frobenius and $\left(\frac3p\right)=+1$. Thus (1) is impossible at every inert $p>3$ in such a field. Counting more primes subject to (1) cannot address this obstruction.

The weaker covering condition can hold even here. Let $\theta$ satisfy $f=X^3-6X^2-3X+2$, irreducible over $\Q$ with discriminant $2700=3\cdot30^2$. Modulo $59$ the cubic has no root, so $59$ is inert. The projective group has prime order $n=3541$, and $g=[\theta^3]\ne1$ generates it. The order of $3$ modulo $3541$ is $1770$, so $\langle3,59\rangle$ consists exactly of the squares. In the basis $(1,\theta,\theta^2)$ two powers of the representative $\theta^3$ reduce to
\[
 g^9:\ (26,4,6),\qquad g^{315}:\ (11,21,51).
\]
The actual representatives have traces $354=6\cdot59$ and $2301=39\cdot59$. Euler's criterion gives $9^{1770}=1$ and $315^{1770}=-1$ modulo $3541$. There is one zero witness in each orbit; this proves coverage of all $3540$ nonidentity classes. It applies to any cubic field in which $59$ is inert, while candidates scalar there still require another prime or argument.

\textbf{Source distinctions.} Corpus [055] averages odd-order symmetries of a linear Okada system to preserve an odd integral constraint. Our orbit calculation uses a different, verified invariance, namely the invariance of trace zero under field Frobenius. The Singer difference set and Frobenius enlargement are present in [075], R15--R18. Here the character sum is over the cyclic group of projective points when counting its index-three subgroup, not a character on the multiplicative group of exponent labels. The two character estimates are distinct. Corpus [006] concerns linear minimum order and an integral solution lattice for Franel numbers. Section 3 concerns nonlinear polynomial identities on a different schedule. The Newton state and the conjugate growth argument are related to corpus [010]. No historical priority or external peer review is claimed.

\textbf{References and verification.} [1] K. Saito, \emph{Transcendency of variants of Mills' constant}, Theorem 1.7. \href{https://arxiv.org/abs/2508.16068}{arXiv:2508.16068}.
[2] \emph{Two early composite floors and recurring divisors for a cubic Mills candidate}, preceding working note. [3] QEEDDD, \href{https://github.com/QEEDDD/Verifiable-Mathematical-Reasoning-Corpus}{Verifiable Mathematical Reasoning Corpus}, snapshot \texttt{3834835}, [006], [010], [055], [075].
The Chapter 7 listing in the verification guide verifies the new discrete logarithms, trace witnesses, orders, early congruences and lattice bound. The two fixed-field numerical bounds use the earlier 147 proved local rows and the earlier upper bound on $M$. Run \texttt{python3 certificates.py run 07} after extracting the combined driver.


% Chapter seven
\chapter{The divisor seven, signed norms and a gap law}\label{ch:seven}\markboth{The divisor seven and signed norms}{}
Let $\theta$ be a root of $f=X^3-X^2-27X+3$, let $K=\Q(\theta)$, and write $F_j(\beta)=\lfloor\beta^{3^j}\rfloor$. This totally real cubic field has discriminant $20172=3\cdot82^2$. Its simple root $1$ modulo $3$ gives an embedding into $\Q_3$. Thus it passes the necessary field test of corpus [042].

\head{1. A field and two full unit branches}
\textbf{Theorem.} For every cubic Pisot unit $\beta$ in any real embedding of $K$, the sequence $F_j(\beta)$ is composite along an infinite arithmetic progression. Moreover, if the larger of its two contracting conjugates in absolute value is negative, there is an $r\bmod4$ such that
\[
                     7\mid F_j(\beta)\qquad (j\gg1,\quad j\equiv r\bmod4).
\]
Both signed norms are included. Section 3 extends the same statements to an explicit multiplicative set of absolute norms containing infinitely many primes.

Let $M$ be the least Mills constant. Saito's Theorem 1.7 [1] implies that an algebraic $M$ has a power $\beta=M^{3^m}$ that is a totally real cubic Pisot number with the indicated negative conjugate. In this case $F_j(\beta)$ is prime for every $j$. The theorem rules out $K$ as the field of that candidate whenever its norm belongs to the set proved here.

\textbf{The earlier field at $59$.} For $f_0=X^3-6X^2-3X+2$, one has $f_0(Y+1)=Y^3-3Y^2-12Y-6$, Eisenstein at $3$. Its cubic field is therefore totally ramified at $3$. By corpus [042], no nonzero rational power of $M$ belongs to this field, for any norm. The finite-field coverage at $59$ was already known in corpus [075], R16. The previous note provided a two-orbit certificate for that coverage. The argument for $K$ below does not require a search for inert primes.

\textbf{Proof ingredients.} Write $U=-796-24\theta+29\theta^2$ and $V=-65+82\theta+20\theta^2$. The exact unit certificate in Section 4 proves
\[
 \mathcal O_K^\times=\{sU^aV^b\colon s\in\{1,-1\},\ (a,b)\in\Z^2\}.
\]
In increasing order of the roots of $f$, both $U$ and $V$ have signs $(-,-,+)$, both have norm $+1$, and $-1$ has norm $-1$.



\head{2. The proof at seven and at three}
For a Pisot element write $A_j=\Tr(\beta^{3^j})$. The two smaller roots have distinct absolute values: otherwise they would be an opposite pair and the remaining root would equal the rational trace. For sufficiently large $j$, $F_j=A_j-h$ with $h=0$ or $1$ according as the larger small root is negative or positive. When $h=0$ the three real signs are mixed. Any positive real conjugate chosen as the Pisot root gives the same trace.

Modulo $7$, $f=(X+1)(X^2-2X+3)$ with the quadratic irreducible. The order $\Z[\theta]$ has index $2$, so this is the residue algebra of $\mathcal O_K$. Denote a class modulo nonzero common scalars by brackets. In the power basis $1,\theta,\theta^2$, set $g=[U^3]$. The certificate gives
\[
 g=[(1,2,5)],\qquad [V^3]=g^5=[(1,6,3)],\qquad
 |g|=16 .
\]
Among $0\le e<16$, the trace of $g^e$ vanishes modulo $7$ precisely at $e=1,7$. Vanishing is independent of the choice of representative. On odd exponents the two cycles under multiplication by $3$ modulo $16$ are
\[
 1,3,9,11\qquad\hbox{and}\qquad 5,15,13,7 .
\]
For $\beta=sU^aV^b$, mixed signs force $a+b$ odd. Then $[\beta^3]=g^{a+5b}$ has odd exponent. Each cycle has one zero position, so for $h=0$ every fourth sufficiently late floor is divisible by $7$ and exceeds $7$. This proves the claimed progression and lower density $1/4$ for the composite indices in this case.

For $h=1$ the following argument is the prime-return criterion of [009]. The residues of all three conjugates of a unit at $3$ lie in $\{1,-1\}$, since $f\bmod3=X^2(X-1)$, the index is prime to $3$, and a unit reduces to a unit. For $sU^aV^b$ these signs are constant under cubing. A local unit of sign $\delta$ satisfies $x^{3^j}\to\delta$ in a finite $3$-adic extension. Hence $A_j$ converges $3$-adically to a sum $\lambda\in\{-3,-1,1,3\}$. As $\lambda-1\ne\pm1$, the valuations $v_3(F_j\pm1)$ stay bounded for large $j$.

If only finitely many $F_j$ are prime, all late floors are composite. Otherwise choose an arbitrarily late prime $p=F_j$ prime to $6\Nm(\beta)$ and the discriminant of its minimal polynomial. The conjugates reduce into $\F_{p^d}^\times$, $d\in\{1,2,3\}$. The bounds on $v_3(p\pm1)$ bound $v_3(p^d-1)$ uniformly (the extra factor $p^2+p+1$ for $d=3$ has valuation at most one). For $j$ sufficiently late, cubing permutes the prime-to-three subgroup containing all three powered conjugates. Their return to the original state gives $p\mid F_{j+kL}$ for some $L\ge1$, all $k\ge0$. All subsequent terms of the progression eventually exceed $p$.



\head{3. Infinitely many norms, including arbitrary powers}
Let $\mathcal P$ be the rational primes $q\nmid 2\cdot3\cdot7\cdot41$ admitting a degree-one prime ideal $P=(\pi)$ in $K$ for which $\pi$ is totally positive and $\pi\equiv5\bmod84\mathcal O_K$. Every cubic Pisot element of $K$ such that all prime divisors of $|\Nm(\beta)|$ belong to $\mathcal P$ satisfies the theorem, with the same period-four assertion when $h=0$. An empty product gives the two unit norms.

Indeed $P$ belongs to the narrow ray class of $(5)$ modulo $84\mathcal O_K$ and all real places. The prime ideal theorem in ray classes [2] supplies a positive density of prime ideals in this class; discarding the degree $>1$ prime ideals leaves infinitely many degree-one ones, since the discarded ones of norm at most $X$ number $O(\sqrt X)$. Conversely the class condition supplies precisely the positive generator with $\pi\equiv5$ at the modulus. Its norm gives $q\equiv5^3\equiv41\bmod84$, hence $q\equiv5\bmod12$ and $(3/q)=-1$. The Frobenius permutation is odd by the discriminant squareclass, so the splitting type in $K$ is $(1,2)$. Writing $(q)=PQ$, the unique other prime $Q$ has totally positive generator $q/\pi$, of nonzero scalar residue modulo $84$ as well.

Every integral ideal of norm supported on $\mathcal P$ is therefore generated by a totally positive element $\gamma$ that reduces to a scalar unit modulo $84$. In particular, every possible $\beta$ in the theorem is $s\gamma U^aV^b$. Since $\gamma$ is positive at all embeddings and scalar at $7$, the proof using mixed signs and the two cycles applies verbatim. The three residues at $3$ are again signs, so the complementary $h=1$ proof also applies. This treats every generator of every possible ideal with these norms, allowing arbitrary powers and either sign of $\Nm(\beta)$.

\textbf{Two certified prime norms.}
\[
\begin{array}{c|r}
\pi&\Nm(\pi)\\ \midrule
425+84\theta&9\,192\,413\\
509-84\theta&14\,916\,929
\end{array}
\]
Both generators are totally positive, and trial division proves both norms prime. Thus all signed norms $\pm9\,192\,413^a14\,916\,929^b$, $a,b\ge0$, are covered, independently of the infinitude theorem.

\textbf{Existence of candidate shapes.} For each permitted positive norm $C$, choose a totally positive generator $\gamma$ of norm $C$, using a product of degree-one generators. The logarithms of squares of units form a rank-two lattice in the sum-zero plane. Multiplying $\gamma U$ or $-\gamma U$ by such squares allows bounded-distance approximations to a translate of $t(1,-0.46,-0.54)$, in any chosen ordering of embeddings. Align the expanding place with a positive sign and the middle place with a negative sign. Both signed norms admit such an ordering. The bounds in [1] then hold strictly for $t$ sufficiently large, since $0.46>17/40$ and $0.46>0.54\cdot17/23$. There are arbitrarily large Pisot elements in the very conjugate ranges to which the exclusion applies.



\head{4. Exact certificate for the unit group}
The cubic $f$ has no rational root among $\pm1,\pm3$. Intervals of width $10^{-18}$ in the attached certificate isolate its three real roots. The element $(1+\theta^2)/2$ has monic polynomial $X^3-29X^2+212X-100$. The lattice with basis $1,\theta,(1+\theta^2)/2$ is integral and has discriminant $20172$. If it were smaller than $\mathcal O_K$, a prime $p$ dividing the index would satisfy $p^2\mid20172$. For $p=2,41$ the certificate checks all $p^3$ classes for an integral vector with denominator $p$ by testing its trace, second coefficient and norm. Only the zero class succeeds; the lattice is the full ring of integers.

The logs of $|U|,|V|$ are independent: at the first place their signs as logarithms are opposite and at the second both are positive. For any unit of norm $+1$, multiply by powers of $U,V$ to place its log vector in their centered fundamental parallelogram. In this parallelogram the three absolute embedding values are less than $1100,220,4100$, respectively. This follows by bounding each value by the square root of the product of $\max(|u|,|u|^{-1})$ for $u=U,V$ at the given place. If the translated unit is $(i+j\theta+k\theta^2)/2$, rational interval inversion of the embedding matrix yields
\[
 |i|<600,\qquad |j|<1000,\qquad |k|<240,\qquad
 j\equiv0\bmod2,\quad i\equiv k\bmod2 .
\]
The exhaustive norm equation $\Nm(i+j\theta+k\theta^2)=8$ in this slightly enlarged box gives exactly three triples:
\[
 (2,0,0),\quad(-130,164,40),\quad(76,-700,122).
\]
They represent $1,V,(UV)^{-1}$. This proves generation of all norm-one units. Negating a unit changes its norm sign, proving the full unit formula. The certificate uses integer modular sieving only to accelerate this exhaustive box check.

\textbf{Scope.} The prime $7$ acts here through mixed reduction, the real signatures of a complete unit basis, and the scalar and sign constraints on generators of ideals. Field-independent coverage by inert primes is unnecessary for this result. Other cubic fields and unrestricted norms in the displayed field remain open in this argument.

\textbf{Sources.} [1] K. Saito, \emph{Transcendency of variants of Mills' constant}, Theorem 1.7, \href{https://arxiv.org/abs/2508.16068}{arXiv:2508.16068}. [2] J. S. Milne, \href{https://www.jmilne.org/math/CourseNotes/CFT.pdf}{\emph{Class Field Theory}}, reciprocity and density in ray classes. [3] QEEDDD, \href{https://github.com/QEEDDD/Verifiable-Mathematical-Reasoning-Corpus/tree/38348353416f1f7fc50cecb93c3b554a00260c13}{VMRC snapshot 3834835}, [009] on prime returns, [042] on the field and sign conditions, [075] R16 on the earlier local certificates. The earlier working note on discriminant squareclass $61$ also used ray classes for norm obstructions. This note concerns a different field and a single recurring divisor for both unit norms.

\textbf{Reproduction.} Run \texttt{python3 certificates.py run 08} using the driver supplied in the verification guide.


\head{5. A necessary gap law independent of the field}
Let $p_j=\lfloor\beta^{3^j}\rfloor$ on a hypothetical Mills tail, let $C=\Nm(\beta)$, and suppose $3\nmid C$. Put
\[
 g_j=p_{j+1}-p_j^3,\qquad
 e_j=\operatorname{cres}_{p_j}(g_j/3),\qquad
 R_j=\frac{g_{j+1}}{g_j e_j^2},
\]
where the centered residue lies in $(-p_j/2,p_j/2)$. At sufficiently large indices [010] gives $e_j=C^{3^j}$ and $g_{j+1}/g_j=\prod_{i<k}(X_i^2-X_iX_k+X_k^2)$ with $X_i=\beta_i^{3^j}$. Thus three consecutive primes determine $R_j$. For a unit $e_j^2=1$.

\textbf{Theorem.} Exactly one of the following two laws holds for all large $j$, with a constant integer $c$ depending on the candidate and indexing:
\[
 v_3(R_j+3)=j+c,\qquad\hbox{or}\qquad v_3(R_j-9)=2j+c.
\]
Here $v_3(3)=1$. These are limits and exact orders in the $3$-adic topology. Failure of both eventual laws would exclude all candidate norms prime to $3$. The required failure has not been shown for the least Mills sequence.

\textbf{Proof.} Write the cubic as $X^3-aX^2+bX-C$. The prime-tail criterion [009] gives $a\equiv\varepsilon=\pm1$, $C\equiv ab\bmod3$, hence the residual cubic is $(X-\varepsilon)(X^2+b)$. In a $3$-adic splitting field write $\beta_i=\zeta_i u_i$ with $\zeta_i$ the root of unity lifting its residue and $u_i$ a principal unit. The three $\zeta_i$ have product $\pm1$, so the product of all $(u_i u_k)^{3^j}$ over pairs is $e_j^2$. Set $m=3^j$, $r_{ik}=(u_i/u_k)^m$, $w_{ik}=r_{ik}+r_{ik}^{-1}-2$, and $H(x)=x-x^{-1}$.

If $b=1\bmod3$, the lifted residues are $(\varepsilon,i,-i)$ with $i^2=-1$. Here $\zeta_1\zeta_2\zeta_3=\varepsilon$. The roots live in an unramified extension, so $u_k\in1+3\OO$. Put $\delta=(i/\varepsilon)^m$. After dividing each pair factor by $u_i^m u_k^m$, its three constant terms are $-1,-1,-3$ at $m=0$ on either parity of $j$. Direct multiplication gives
\[
 R_j=-(3+w_{23})\{1+\delta[H(r_{12})-H(r_{13})]+H(r_{12})H(r_{13})\}.
\]
Write $s=v_3(u_2/u_3-1)\ge1$, finite because the original roots have distinct absolute values. Since $H(x)-H(y)=(x-y)(1+(xy)^{-1})$, the valuation of the linear contribution to $R_j+3$ is $1+j+s$ by the lifting-the-exponent identity for principal units. The quadratic contribution has valuation at least $2j+3$, and the term containing $w_{23}$ has valuation $2j+2s$. Thus for $j\ge s$ the first law holds.



\head{6. The other residue pattern}
If $b=-1\bmod3$, choose the ordering so that $\zeta_1$ is the unique sign and $\zeta_2=\zeta_3=-\zeta_1$. The normalized product is exactly
\[
       R_j=(3+w_{12})(3+w_{13})(1+w_{23}).
\]
Set $A=\log_3(u_1/u_2)$ and $B=\log_3(u_1/u_3)$, where $\log_3$ denotes the $3$-adic logarithm. For sufficiently large $j$, $r_{12}=\exp(mA)$ and $r_{13}=\exp(mB)$; raising principal units to large powers of three removes any $3$-power torsion. Expanding the convergent series at $m=0$ gives
\[
 R_j/9=1+m^2 Q(A,B)+O(m^4),\qquad
 Q(A,B)=\frac{A^2+B^2}{3}+(A-B)^2 .
\]
The bound here is in the $3$-adic absolute value. The quadratic coefficient does not vanish. Indeed a rational linear relation between $A$ and $B$ would make a nonzero integral power of
$(\beta_1/\beta_2)^r(\beta_1/\beta_3)^s$ equal to $1$ for some integers $(r,s)\ne(0,0)$, since the logarithm kernel on local units is torsion. Its exponent vector $(r+s,-r,-s)$ and its cyclic permutations span the rational plane of coordinate sum zero: the Galois group contains a three-cycle, whose action on that plane is irreducible over $\Q$. Conjugating the multiplicative relation would force all three real absolute values to agree, which is impossible. This proves rational linear independence. The Baker--Brumer theorem [4], applied after raising to a common power within its convergence neighborhood, gives linear independence over the algebraic numbers. Since $Q$ factors into linear forms over the algebraic numbers, $Q(A,B)\ne0$. The quadratic term now gives $v_3(R_j-9)=2j+2+v_3(Q(A,B))$ for all large $j$. This proves the second law and completes the dichotomy.

\textbf{Verification of both patterns.} The attached script checks finite instances using only integer Newton identities. These finite examples check the two expansions; the theorem for a hypothetical Mills tail follows from the proof above. The logical input from [010] is its gap and norm reconstruction, from [031] its local multiplier viewpoint, and from [075], R17 its use of logarithmic independence. The result here measures the distance of the normalized multiplier to $-3$ or $9$.

[4] A. Brumer, \href{https://www.cambridge.org/core/journals/mathematika/article/abs/on-the-units-of-algebraic-number-fields/8390DA0413F8229C2961A35604480784}{\emph{On the units of algebraic number fields}}, Mathematika 14 (1967), 121--124.

\backmatter
\chapter{Verification guide}\markboth{Verification guide}{}
The source is a standalone LaTeX document. Run pdfLaTeX twice on \texttt{Mills\_Collected.tex} to resolve the table of contents and cross-references. The numerical certificates follow the final document marker as an executable archive. The archive and its ten listings are plain text. To extract the driver on a system with \texttt{sed}, use
\begin{Verbatim}[fontsize=\small,breaklines]
sed '1,/^% BEGIN MILLS CERTIFICATE ARCHIVE$/d' Mills_Collected.tex > certificates.py
python3 certificates.py list
python3 certificates.py all --thorough
\end{Verbatim}
Python 3.9 or newer and a C++17 compiler named \texttt{c++} on the command path suffice for the collection. The driver runs each certificate in a fresh temporary directory and checks its SHA-256 digest. Select a single chapter with, for example, \texttt{python3 certificates.py run 04}. Save all ten original listings as individual files with \texttt{python3 certificates.py extract all programs}.

\begin{center}\small\begin{tabular}{@{}clp{79mm}@{}}\toprule
Chapter & Listing & Computation\\\midrule
1 & \texttt{01-pisot.py} & Finite search through the coefficient bounds of the canonical cubic.\\[2pt]
2 & \texttt{02-gap.py} & Cube-residue failures at three of nine published transitions.\\[2pt]
3 & \texttt{03-resolvent.py} & Complete enumeration of the stated local character cycles.\\[2pt]
4 & \texttt{04-orbits.cpp} & Projective orbits for 22 primes.\\[2pt]
4 & \texttt{04-barrier.py} & The 430-prime avoidance construction.\\[2pt]
5 & \texttt{05-mixed.cpp} & Mixed-prime orbits and trace-zero witnesses.\\[2pt]
5 & \texttt{05-lattice.py} & Unit coordinates, log intervals and exponent lattice.\\[2pt]
6 & \texttt{06-early.py} & Early floors, 554-prime construction and enlarged lattice.\\[2pt]
7 & \texttt{07-franel.py} & Ten more congruences, the degree bound and the two cycles at 59.\\[2pt]
8 & \texttt{08-seven.py} & Unit generation, trace-zero cycles, signed norms and both gap examples.\\\bottomrule
\end{tabular}\end{center}
The finite checks can be rerun independently of the analytic arguments. The proof of the infinite ray-prime family uses class field theory. The density-one reciprocal-sum theorem uses the Subspace Theorem. The necessary gap law at the center $+9$ uses the Baker--Brumer theorem. Their hypotheses and applications appear in the corresponding chapters.
\end{document}

% BEGIN MILLS CERTIFICATE ARCHIVE
#!/usr/bin/env python3
"""Extract, inspect, or run the original finite certificates in this volume.

Requires Python 3.9 or later. The two C++ listings require a C++17 compiler
named c++ on PATH. Each certificate runs in its own temporary directory, with
assertions enabled. No network access or third-party packages are needed.
"""
import argparse
import hashlib
from pathlib import Path
import subprocess
import sys
import tempfile

PROGRAMS = {}
DESCRIPTIONS = {}
DIGESTS = {}

# Source: verify_mills_pisot.py
DESCRIPTIONS['01-pisot.py'] = 'Coefficient catalogue and recurrent divisors through 10,000.'
DIGESTS['01-pisot.py'] = '2ff3ff658c77719597702265700690a5bc4ac0fef1c0eded5d8df486d558a756'
PROGRAMS['01-pisot.py'] = r'''"""Certify that a cubic Pisot shadow of the algebraic least Mills constant
must exceed 10,000. Run with Python 3.9+ and no external requirements.

For a quick check of the same method at 1,000, pass --bound 1000.
The proof and assumptions are given in the accompanying manuscript.
This program tests all triples in a superset of the necessary catalogue.
A position on a nonrecurring part of a finite orbit never certifies an
exclusion. A zero trace on its eventual cycle always does.
"""
import argparse
from array import array


SMALL_PRIMES = (2, 3, 5, 7, 11, 13, 17, 19, 23, 29, 31, 37, 41, 43, 47,
                53, 59, 61, 67, 71, 73, 79, 83, 89, 97, 101, 103, 107)

# Each record contains the two coefficient bounds, the number of starting
# triples, and the successive survivor counts after each prime.
CASES = {
    1000: (2, 108, 307930,
           (117653, 78411, 34586, 13362, 6747, 3416, 1991, 1303, 616,
            343, 180, 85, 40, 15, 4, 3, 0)),
    10000: (3, 401, 18064626,
            (7533373, 5022141, 2208469, 857106, 431440, 217659, 126995,
             82780, 39078, 22311, 12109, 5049, 2627, 937, 257, 185, 69,
             35, 17, 14, 9, 6, 3, 2, 1, 1, 1, 0)),
}


def candidate_coefficients(bound, max_d, max_c):
    B = bound
    for a in range(1, B + 1):
        a3 = a**3
        for d in range(1, max_d + 1):
            d3 = d**3
            lead = a*a*d*d + 4*d3
            p_at_B = B*B*(B-a) - d*B
            for c in range(1, max_c + 1):
                if d <= 2*a+2 and d <= a+c and p_at_B > c:
                    discr = lead - 4*a3*c - 27*c*c - 18*a*d*c
                    if discr > 0:
                        yield a, -d, c
                if d+c <= a:
                    yield a, -d, -c


def step(state, p):
    t, b, e = state
    return ((t**3 - 3*t*b + 3*e) % p,
            (b**3 - 3*t*b*e + 3*e*e) % p,
            pow(e, 3, p))


def survival_test(p):
    known = {}

    def survives(a, b, c):
        state = a % p, b % p, c % p
        if state in known:
            return known[state]
        seen = {}
        while state not in known and state not in seen:
            seen[state] = len(seen)
            state = step(state, p)
        if state in known:
            verdict = known[state]
        else:
            first = seen[state]
            verdict = all(s[0] != 0 for s, j in seen.items() if j >= first)
        known.update(dict.fromkeys(seen, verdict))
        return verdict

    return survives


def pack(a, b, c):
    # 0 < a <= 10000, 1 <= -b <= 401, and -3 <= c <= 3.
    return (a << 12) | ((-b) << 3) | (c+3)


def unpack(code):
    return code >> 12, -((code >> 3) & 511), (code & 7)-3


def verify(bound):
    max_c, max_d, total, expected = CASES[bound]
    assert bound**3 < (max_c+1)**20
    assert (2**40)*bound**23 < (max_d-1)**40
    assert pack(bound, -max_d, max_c) < 2**32
    assert array('I').itemsize >= 4

    stages = zip(SMALL_PRIMES[:len(expected)], expected)
    prime, survivors = next(stages)
    assert prime == 2
    test = survival_test(prime)
    pending = array('I')
    n = 0
    for a, b, c in candidate_coefficients(bound, max_d, max_c):
        n += 1
        if test(a, b, c):
            pending.append(pack(a, b, c))
    assert n == total
    assert len(pending) == survivors
    print(f'{n:,} starting triples. After prime 2: {len(pending):,}.', flush=True)
    for prime, survivors in stages:
        assert all(prime % d for d in range(2,prime) if d*d <= prime)
        test = survival_test(prime)
        pending = array('I', (code for code in pending if test(*unpack(code))))
        assert len(pending) == survivors
        print(f'After prime {prime}: {len(pending):,}.', flush=True)
    assert not pending
    print(f'PASS: canonical cubic Pisot root beta > {bound:,}.')


if __name__ == '__main__':
    parser = argparse.ArgumentParser(description=__doc__)
    parser.add_argument('--bound', type=int, choices=sorted(CASES), default=10000)
    verify(parser.parse_args().bound)
'''

# Source: Mills_Gap_Residues.tex
DESCRIPTIONS['02-gap.py'] = 'Local congruences for the nine published offsets.'
DIGESTS['02-gap.py'] = '82c460f214ed5a4c93a69d90ebb322ded46b6f5eba0816cbbd0690cb5211ec1d'
PROGRAMS['02-gap.py'] = r'''gaps = [3, 30, 6, 80, 12, 450, 894, 3636, 70756]
b = 2
failures = []
for n, gap in enumerate(gaps, start=1):
    if not any(pow(x, 3, gap) == b % gap for x in range(gap)):
        failures.append(n)
    if n == 9:
        assert (b % 7, b % 19) == (2, 6)
    b = b**3 + gap
assert failures == [6, 8, 9]
'''

# Source: Mills_Resolvent_Exclusions.tex
DESCRIPTIONS['03-resolvent.py'] = 'Discriminant-character cycle enumeration.'
DIGESTS['03-resolvent.py'] = 'd127d0a35ca41a373eebee9a269ba30fe0b0adffc4fd589e9c74e960b32395bb'
PROGRAMS['03-resolvent.py'] = r'''from math import isqrt

def prime(p):
    return p >= 2 and all(p % d for d in range(2, isqrt(p) + 1))

def all_cycles_good(p):
    h = p - 1
    while h % 3 == 0:
        h //= 3
    seen = set()
    for x in range(1, p):
        if x in seen or pow(x, h, p) != 1:
            continue
        z, good = x, False
        while z not in seen:
            seen.add(z)
            good |= pow(1 + 2*z, (p-1)//2, p) != p-1
            z = pow(z, 3, p)
        if not good:
            return False
    return True

primes = [p for p in range(13, 2000, 12) if prime(p)]
print(len(primes), [p for p in primes if all_cycles_good(p)])
assert all(prime(2*3**s+1) for s in (1,2,4,5,6,9,16,17))
'''

# Source: Mills_Discriminant_61.tex
DESCRIPTIONS['04-orbits.cpp'] = 'All projective orbits for the 22 primes in the squareclass-61 field.'
DIGESTS['04-orbits.cpp'] = '1badc81b8b10a7fcd58fe8b055ba1b577c44bfb09e32267e63aef8bd8fe04d42'
PROGRAMS['04-orbits.cpp'] = r'''#include <array>
#include <cassert>
#include <vector>
using A = std::array<int,3>;
int p;
int mod(long long v) { return (v%p+p)%p; }
bool prime(int n) {
  if (n<2) return false;
  for(int d=2; d<=n/d; ++d) if(n%d==0) return false;
  return true;
}
A mul(A a,A b) {
  long long c[5]={};
  for(int i=0;i<3;++i) for(int j=0;j<3;++j) c[i+j]+=1LL*a[i]*b[j];
  for(int k=4;k>=3;--k) {
    c[k-1]-=19*c[k]; c[k-2]+=19*c[k]; c[k-3]+=19*c[k];
  }
  return {mod(c[0]),mod(c[1]),mod(c[2])};
}
int main() {
  int primes[]={5,41,97,103,127,131,167,199,271,293,773,827,
                859,911,1021,1163,1483,1583,1709,1723,1811,1993};
  for(int v:primes) {
    p=v; assert(prime(p));
    for(long long x=0;x<p;++x) assert(mod(x*x*x+19*x*x-19*x-19));
    int S=p*p+p+1, m=S%3?S:S/3;
    assert(prime(m) && m%3);
    A e={mod(-56),mod(37),2}, g={1,0,0}, a={1,0,0};
    for(int i=0;i<S/m;++i) g=mul(g,e);
    assert(g[1] || g[2]);
    std::vector<bool> hit(m,false);
    for(int i=0;i<m;++i) {
      if(mod(3*a[0]-19*a[1]+399*a[2])==0)
        for(int j=i;!hit[j];j=(3*j)%m) hit[j]=true;
      a=mul(a,g);
    }
    assert(a[0] && !a[1] && !a[2] && !hit[0]);
    for(int i=1;i<m;++i) assert(hit[i]);
  }
}
'''

# Source: Mills_Discriminant_61.tex
DESCRIPTIONS['04-barrier.py'] = 'Small-prime avoidance on the 430-prime ray.'
DIGESTS['04-barrier.py'] = '4a4b2b7ff5cf14330d0db4930157da390cf4b11503679ac063f9e4b50aec2e30'
PROGRAMS['04-barrier.py'] = r'''from math import isqrt, lcm
def prime(p):
    return p >= 2 and all(p % d for d in range(2, isqrt(p)+1))
def mul(a, b, p):
    c = [0]*5
    for i in range(3):
        for j in range(3):
            c[i+j] += a[i]*b[j]
    for k in (4,3):
        c[k-1] -= 19*c[k]; c[k-2] += 19*c[k]; c[k-3] += 19*c[k]
    return tuple(x % p for x in c[:3])
def power(a, n, p):
    z = (1,0,0)
    while n:
        if n % 2: z = mul(z,a,p)
        a = mul(a,a,p)
        n //= 2
    return z
primes = [p for p in range(2,3001) if prime(p)]
n = 1
for p in primes:
    m = p*(p*p-1)*(p**3-1)
    while m % 2 == 0: m //= 2
    n = lcm(n,m)
states = 0
for p in primes:
    z = power((3564,-2150,-117),n,p)
    seen = set()
    while z not in seen:
        seen.add(z)
        assert (3*z[0]-19*z[1]+399*z[2]) % p
        z = mul(mul(z,z,p),z,p)
    states += len(seen)
assert (len(primes),len(str(n)),states) == (430,2874,1485)
'''

# Source: Mills_Unit_Ray_Return.tex
DESCRIPTIONS['05-mixed.cpp'] = 'Exhaustive mixed-prime orbit rows and returning witnesses.'
DIGESTS['05-mixed.cpp'] = '33b40fe0af280a543b969638cec93200a4870ea9036992c6947b6039cf1ce848'
PROGRAMS['05-mixed.cpp'] = r'''#include <array>
#include <vector>
#include <numeric>
#include <cassert>
#include <cstdint>
using I = std::int64_t;
using A = std::array<I,3>;
int p;
const A one={1,0,0}, E={-56,37,2}, T={-46,27,2};
I mod(I a){return (a%p+p)%p;}
A mul(A a,A b) {
    I w=a[2]*b[2], v=a[1]*b[2]+a[2]*b[1]-19*w;
    return {mod(a[0]*b[0]+19*v),
            mod(a[0]*b[1]+a[1]*b[0]+19*w+19*v),
            mod(a[0]*b[2]+a[1]*b[1]+a[2]*b[0]-19*v+19*w)};
}
A pw(A a,I n) {
    A r=one;
    while(n){if(n%2)r=mul(r,a); a=mul(a,a); n/=2;}
    return r;
}
bool scalar(A a){return a[0]!=0 && a[1]==0 && a[2]==0;}
bool same(A a,A b){
    for(int i=0;i<3;i++)for(int j=0;j<i;j++)
        if(mod(a[i]*b[j]-a[j]*b[i]))return false;
    return true;
}
bool zero(A a){return mod(3*a[0]-19*a[1]+399*a[2])==0;}
bool prime(I n) {
    if(n<2)return false;
    for(I i=2;i<=n/i;i++) if(n%i==0)return false;
    return true;
}
std::vector<I> factors(I n){
    std::vector<I> result;
    for(I d=2;d<=n/d;d++)if(n%d==0){
        result.push_back(d);while(n%d==0)n/=d;
    }
    if(n>1)result.push_back(n);
    return result;
}
void order(A a,I n){
    assert(scalar(pw(a,n)));
    for(I q:factors(n))assert(!scalar(pw(a,n/q)));
}
I powmod(I a,I n,I q){
    I r=1;
    while(n){if(n%2)r=r*a%q; a=a*a%q; n/=2;}
    return r;
}
void evaluations(int q,int l,int g,std::array<int,3> roots,
                 std::array<int,4> logs){
    assert(prime(q)&&prime(l)&&g!=1&&powmod(g,l,q)==1);
    for(I x:roots)assert((x*x*x+19*x*x-19*x-19)%q==0);
    for(int i=0;i<2;i++){
        I x=roots[i];
        assert(powmod((2*x*x+37*x-56)%q,(q-1)/l,q)==powmod(g,logs[2*i],q));
        assert(powmod((2*x*x+27*x-46)%q,(q-1)/l,q)==powmod(g,logs[2*i+1],q));
    }
    assert((logs[0]*logs[3]-logs[1]*logs[2])%l);
}
void six(A g,I N,I m,I period,std::vector<I> reps){
    A h=pw(g,N/m);std::vector<bool> marks(m,false);
    for(I e:reps){
        assert(std::gcd(e,m)==1 && !marks[e] && zero(pw(h,e)));
        I r=e, count=0;
        do {assert(!marks[r]); marks[r]=true; count++;r=r*3%m;}while(r!=e);
        assert(count==period);
    }
    for(I e=1;e<m;e++)if(std::gcd(e,m)==1)assert(marks[e]);
}
void whole_row(A g,I N,I m,int good_count,int cycle_count){
    std::vector<bool> seen(m,false);
    A h=pw(g,N/m);
    int good=0,total=0;
    for(I e=1;e<m;e+=2){
        if(seen[e]||std::gcd(e,m)!=1)continue;
        I j=e; A x=pw(h,e); bool hit=false;
        do{
            seen[j]=true;
            if(!hit){hit=zero(x); if(!hit)x=mul(mul(x,x),x);}
            j=3*j%m;
        }while(j!=e);
        good+=hit;total++;
    }
    assert(good==good_count && total==cycle_count);
}
int main(){
    evaluations(1933,7,1812,{71,112,1731},{1,4,3,4});
    evaluations(937,13,911,{501,666,688},{9,10,9,1});
    for(int val:{21503,34303}){
        p=val;assert(prime(p)); int roots=0;
        for(I x=0;x<p;x++)if(mod(x*x*x+19*x*x-19*x-19)==0)roots++;
        assert(roots==1 && (61*114*114)%p);
        I S=I(p)*p-1, N=S/3;
        A g=pw({p==21503?2:1,1,0},3);
        order(g,N); order(pw(E,3),p==21503?N:N/2);
        order(pw(T,3),p==21503?N/2:N);
        std::vector<std::array<int,2>> logs;
        std::vector<std::array<int,3>> rows;
        if(p==21503){
            six(g,N,2048,512,{231,793});
            six(g,N,1693696,211456,{36681,106441,193855,598607});
            logs={{7,1},{13,12},{827,170}};
            rows={{1,2,2},{7,0,4},{13,2,8},{91,4,48},
                  {827,4,4},{5789,8,8},{10751,16,16},{75257,96,96}};
        }else{
            logs={{67,34},{5717,2537}};
            rows={{1,2,2},{67,0,12},{5717,8,8},{383039,48,48}};
        }
        for(auto [q,k]:logs){
            A x=pw(E,S/q),y=pw(T,S/q);
            assert(!scalar(x) && same(y,pw(x,k)));
        }
        for(auto [d,hit,total]:rows)
            whole_row(g,N,(p==21503?2048LL:1024LL)*d,hit,total);
    }
}
'''

# Source: Mills_Unit_Ray_Return.tex
DESCRIPTIONS['05-lattice.py'] = 'Unit generation and the exponent-lattice bound.'
DIGESTS['05-lattice.py'] = 'df4fe0a5f08157c3b78c00e7b43a8f0d18f3529df25fce37b587a68540467a0b'
PROGRAMS['05-lattice.py'] = r'''"""Verify the D=61 norm-one exponent-lattice bound. Python 3, standard library.

The optional --full flag rechecks the earlier twenty-two exhaustive finite-field
certificates. This part traverses several million cubic finite-field elements.
"""
from fractions import Fraction as F
from math import gcd, isqrt
import sys

DATA = [
    (5,1),(41,787),(97,2951),(103,1915),
    (127,4419),(131,14136),(167,11077),(199,8942),
    (271,5781),(293,50847),(773,552873),(827,419023),
    (859,50622),(911,21647),(1021,260342),(1163,549066),
    (1483,349678),(1583,255229),(1709,2745692),(1723,654871),
    (1811,205927),(1993,630756),(3557,8372873),(4889,3215080),
    (8681,46418403),(10589,44659163),(13469,31295481),(16229,77599160),
    (16301,66708745),(16421,76985174),(16553,128280463),(19013,116496427),
    (19661,95133794),(20849,270934284),(22277,350529041),(22937,297631469),
    (26489,31896576),(26921,404448787),(28901,65818922),(31277,924860164),
    (32693,1008873991),(36497,695170719),(36749,1317741455),(37049,1219522376),
    (37649,1390873606),(38669,243832387),(38993,1250966396),(42797,987290451),
    (46229,1742607656),(47309,941740647),(49937,5343544),
]
ONE=(1,0,0)
E=(-56,37,2)
T=(-46,27,2)

def mul(a,b,p):
    c=[0]*5
    for i in range(3):
        for j in range(3):
            c[i+j]+=a[i]*b[j]
    for k in (4,3):
        c[k-1]-=19*c[k]; c[k-2]+=19*c[k]; c[k-3]+=19*c[k]
    return tuple(x%p for x in c[:3])

def power(a,n,p):
    z=ONE
    while n:
        if n%2: z=mul(z,a,p)
        a=mul(a,a,p); n//=2
    return z

def project(a,p):
    c=pow(next(x%p for x in a if x%p),-1,p)
    return tuple(x*c%p for x in a)

def prime(n):
    return n>=2 and all(n%d for d in range(2,isqrt(n)+1))

def factors(n):
    d=2
    while d*d<=n:
        if n%d==0:
            yield d
            while n%d==0: n//=d
        d+=1
    if n>1: yield n

def trac(a,p):
    return (3*a[0]-19*a[1]+399*a[2])%p

def lattice():
    N,h=1,0
    for n,(p,k) in enumerate(DATA):
        S=p*p+p+1; m=S//gcd(S,3)
        assert prime(p) and prime(m) and gcd(N,m)==1
        assert all((t**3+19*t*t-19*t-19)%p for t in range(p))
        assert project(E,p)!=ONE and project(power(E,m,p),p)==ONE
        assert project(power(E,k,p),p)==project(T,p)
        if n>=22:
            assert p%3==2
            assert all(pow(3,(m-1)//d,m)!=1 for d in factors(m-1))
        elif '--full' in sys.argv:
            seen=bytearray(m); x=ONE
            for i in range(m):
                if not trac(x,p):
                    j=i
                    while not seen[j]:
                        seen[j]=1; j=3*j%m
                x=mul(E,x,p)
            assert sum(seen)==m-1 and not seen[0]
        h+=N*((k-h)*pow(N,-1,m)%m); N*=m
    for m,k in [(827,170),(5717,2537)]:
        # Projective congruences and cycle coverage at 21503,34303 in Appendix A.
        assert gcd(N,m)==1
        h+=N*((k-h)*pow(N,-1,m)%m); N*=m
    def dp(a,b): return sum(x*y for x,y in zip(a,b))
    u,v=(N,0),(-h,1)
    while True:
        if dp(v,v)<dp(u,u): u,v=v,u
        d=dp(u,u); q=(2*dp(u,v)+d)//(2*d)
        if not q: break
        v=tuple(x-q*y for x,y in zip(v,u))
    assert abs(u[0]*v[1]-u[1]*v[0])==N
    return u,v

def normform(x,y,z):
    return (x**3-x*x*z-76*x*y*y-38*x*y*z-101*x*z*z
            -228*y**3-532*y*y*z+646*y*z*z+387*z**3)

def units():
    for p in (2,19):
        for x in range(p):
            for y in range(p):
                for z in range(p):
                    b=3*x*x-2*x*z-76*y*y-38*y*z-101*z*z
                    if (3*x-z)%p==b%(p*p)==normform(x,y,z)%(p**3)==0:
                        assert (x,y,z)==(0,0,0)
    solutions=set()
    for y in range(-26,27):
        for z in range(-34,35):
            q=-76*y*y-38*y*z-101*z*z; r=normform(0,y,z)
            for x in range(-230,231):
                if ((x-z)*x+q)*x+r==1: solutions.add((x,y,z))
    assert solutions=={(1,0,0),(-24,5,1),(46,15,-9),(-132,-22,19)}

def interval(a,b=None):
    if b is None:b=a
    return (F(a),F(b))
def add(a,b):return a[0]+b[0],a[1]+b[1]
def neg(a):return -a[1],-a[0]
def prod(a,b):
    s=[x*y for x in a for y in b]
    return min(s),max(s)
def inv(a):
    assert a[0]*a[1]>0
    return 1/a[1],1/a[0]
def sub(a,b):return add(a,neg(b))
def ab(a):
    assert a[0]*a[1]>0
    return a if a[0]>0 else neg(a)
def dot(a,b):return add(prod(a[0],b[0]),prod(a[1],b[1]))
def inv2(M):
    a,b=M[0];c,d=M[1];D=inv(sub(prod(a,d),prod(b,c)))
    return [[prod(t,D) for t in R] for R in [[d,neg(b)],[neg(c),a]]]
def sqrt_ge(a,x):return a>=x*x
def natlog_one_two(x):
    t=(x-1)/(x+1)
    t=F(t); t2=t*t; q=t; s=0;n=18
    for k in range(n):s+=2*q/(2*k+1);q*=t2
    return (s,s+2*q/((2*n+1)*(1-t2)))
logtwo=natlog_one_two(F(2))
def logpoint(x):
    assert x>0;k=0
    while x<1:x*=2;k-=1
    while x>2:x/=2;k+=1
    return add(natlog_one_two(x),prod(interval(k),logtwo))
def logiv(a):
    lo,hi=logpoint(a[0])[0],logpoint(a[1])[1]
    f=10**8
    return interval(F((lo*f).__floor__(),f),F((hi*f).__ceil__(),f))
def box(n):return interval(F(n,10**12),F(n+1,10**12))
theta=[box(t) for t in [-19906514237356,-623731016520,1530245253875]]
def poly(a,coeffs):
    v=interval(0)
    for c in coeffs:v=add(prod(v,a),interval(c))
    return v
for a,b in theta:
    def f(x):return x**3+19*x*x-19*x-19
    assert f(a)*f(b)<0
es=[poly(t,[2,37,-56]) for t in theta]
ts=[poly(t,[2,27,-46]) for t in theta]
logs=[[logiv(ab(e)),logiv(ab(t))] for e,t in zip(es,ts)]
G=[[sum((prod(R[i],R[j])[0] for R in logs),F()),sum((prod(R[i],R[j])[1] for R in logs),F())] for i,j in [(0,0),(1,1),(0,1)]]
g,h,j=G
assert g[0]>47 and (g[0]-47)*(h[0]-47)>max(j[0]**2,j[1]**2)
def box_certificate():
    bounds=(300,71,268)
    for i in range(3):
        e,t=ab(es[i]),ab(ts[i])
        assert max(e[1],1/e[0])*max(t[1],1/t[0]) < bounds[i]**2
    R=[ [interval(1),poly(t,[F(1,3),6,F(-19,3)]),
          poly(t,[F(1,3),7,F(-1,3)])] for t in theta]
    def cofactor(i,j):
        u,v=[k for k in range(3) if k!=i]
        r,s=[k for k in range(3) if k!=j]
        det=sub(prod(R[u][r],R[v][s]),prod(R[u][s],R[v][r]))
        return neg(det) if (i+j)%2 else det
    D=interval(0)
    for j in range(3):D=add(D,prod(R[0][j],cofactor(0,j)))
    for i,maximum in enumerate((230,26,34)):
        assert sum(bounds[j]*ab(prod(inv(D),cofactor(j,i)))[1]
                   for j in range(3))<maximum
    units()

def cone_certificate():
    U,V=lattice(); scale=10**183; l10=logiv(interval(10))
    W=[[dot(r,[interval(F(x,scale)) for x in vec]) for vec in (U,V)]
       for r in logs[:2]]
    bounds=[int((ab(r[0])[1]*F(23,40)+ab(r[1])[1])*2000*l10[1])
            for r in inv2(W)]
    assert bounds==[39,24]
    candidates={}
    for i in range(-39,40):
        for j in range(-24,25):
            a,b=(i*U[t]+j*V[t] for t in (0,1))
            if not (a%2 and b%2 and (a+b)%5==0):continue
            X,Y=(dot(r,[interval(i),interval(j)]) for r in W)
            if Y[1]<=0 or Y[0]>2000*l10[1]:continue
            if sub(prod(interval(-40),X),prod(interval(17),Y))[1]<0:continue
            if add(prod(interval(23),Y),prod(interval(40),X))[1]<0:continue
            candidates[i,j]=(a,b)
    table={(8,-7):(11,2,4),(18,-17):(23,10,20),
           (22,-23):(11,0,4),(24,-21):(11,1,4)}
    assert candidates.keys()==table.keys()
    for ij,(a,b) in candidates.items():
        p,r,h=table[ij]
        # E and T have order dividing p^2-1: splitting type (1,2) or (1,1,1).
        d=p*p-1
        assert power(E,d,p)==power(T,d,p)==ONE
        z=mul(power(E,a%d,p),power(T,b%d,p),p)
        y=power(z,3**r,p)
        assert trac(y,p)==0 and power(y,3**h,p)==y

box_certificate()
cone_certificate()
print("Verified log10(beta) > 2*10^186.")
'''

# Source: Mills_Early_Floors.tex
DESCRIPTIONS['06-early.py'] = 'Early and periodic divisors, the 554-prime construction and lattice.'
DIGESTS['06-early.py'] = 'ddd22366478fd7452b41af2518e45d97cf30a7e06298c8e75d63facef3eeb6be'
PROGRAMS['06-early.py'] = r'''"""Run with Python 3.9+; --full repeats the inherited 22 exhaustive orbit checks.
This file is embedded verbatim in the TeX source. No third-party dependencies.
"""
from math import gcd, isqrt
ONE=(1,0,0); E=(-56,37,2); T=(-46,27,2)
def fac(n):
    d=2
    while d*d<=n:
        if n%d==0:
            yield d
            while n%d==0:n//=d
        d+=1
    if n>1:yield n
def prime(n):return n>1 and next(fac(n))==n
def mul(a,b,p):
    w=a[2]*b[2];v=a[1]*b[2]+a[2]*b[1]-19*w
    return ((a[0]*b[0]+19*v)%p,(a[0]*b[1]+a[1]*b[0]+19*w+19*v)%p,
            (a[0]*b[2]+a[1]*b[1]+a[2]*b[0]-19*v+19*w)%p)
def pw(a,n,p):
    x=ONE
    while n:
        if n%2:x=mul(x,a,p)
        a=mul(a,a,p);n//=2
    return x
def proj(a,p):
    s=pow(next(x%p for x in a if x%p),-1,p)
    return tuple(x*s%p for x in a)

def order3(n):
    u=n-1
    for r in fac(u):
        while u%r==0 and pow(3,u//r,n)==1:u//=r
    return u

def verify(p,l,q,k,wq,wm,slope):
    S=p*p+p+1; N=S//gcd(3,S)
    assert N==l*q and prime(p) and prime(l) and prime(q) and p not in (2,3,19,61)
    assert pow(61,(p-1)//2,p)==1 and pw((0,1,0),p,p)!=(0,1,0)
    assert proj(pw(E,N,p),p)==proj(pw(T,N,p),p)==ONE
    assert proj(pw(E,q,p),p)!=ONE and proj(pw(E,l,p),p)!=ONE
    assert proj(pw(pw(E,l,p),k,p),p)==proj(pw(T,l,p),p)
    def zero(x):
        a,b,c=x
        return (3*a-19*b+399*c)%p==0
    small=pw(E,q,p)
    assert all(not zero(pw(small,j,p)) for j in range(l))
    if slope is not None:
        assert proj(pw(T,q,p),p)==proj(pw(small,slope,p),p)
    h=pw(E,l,p);oq=order3(q);ol=order3(l)
    assert len({pow(e,oq,q) for e in wq})==len(wq)==(q-1)//oq
    assert all(gcd(e,q)==1 and zero(pw(h,e,p)) for e in wq)
    tab={}
    for a in range(1,l):
        if a in tab:continue
        u=a;j=0
        while u not in tab:tab[u]=(a,j);u=u*3%l;j+=1
        assert j==ol
    d=oq//gcd(oq,ol);labels=set()
    for e in wm:
        assert gcd(e,N)==1 and zero(pw(E,e,p))
        c,j=tab[e%l];labels.add((c,pow(e*pow(3,-j,q),d,q)))
    assert len(labels)==len(wm)==(l-1)//ol*((q-1)//d)


records = [
    (967, 67, 4657, 4170, [1128, 2054],
     [184212, 47210, 280008, 55928, 197712, 217700, 173599, 279556, 290734, 124768, 187701, 96767], 57),
    (2441, 67, 88969, 64229, [27148, 20053],
     [1771578, 5697903, 3460021, 5840972, 5948750, 4267817, 4798003, 3262280, 5053056, 1005385, 732753,
      5745398, 5616026, 2090047, 5857728, 5448435, 4204172, 4878747, 660902, 4226262, 4096449, 3545532,
      4871228, 1656792, 686576, 3115271, 5434673, 5208495, 1806768, 4254624, 2659159, 884265, 1186419,
      411722, 1922343, 2686775, 3032479, 5900973, 3794842, 543514, 1474900, 5190206, 922834, 790510,
      3260068, 4856212, 2388794, 2856041, 3461136, 5500666, 862425, 922397, 1151284, 5931991, 4640371,
      2320536, 4791269, 393263, 1765552, 2751953, 1464892, 3493385, 3026716, 4588397, 5152174, 1177152,
      4477207, 3057801, 5159409, 3918543, 3047431, 3269630, 1266546, 5956599, 1107583, 881662, 1572261,
      4217038, 5317719, 4547228, 2711208, 5771401, 79117, 2917236, 5373806, 1595521, 3932490, 1351314,
      4544514, 5245759, 2878484, 1049458, 377037, 5224522, 1075218, 1968678, 881348, 4902043, 3962459,
      5707263, 4954211, 5398461, 1673312, 4519913, 1828023, 3311552, 5438892, 5781449, 3205420, 1394781,
      3072552, 2122060, 3504055, 4847051, 762604, 1723100, 258569, 176393, 5393275, 4676194, 2910484,
      2098029, 4332677, 5895873, 5801571, 4050889, 5349958, 4163838, 4463264, 5416908, 2862737,
      565404],
     43),
    (3049, 7, 442831, 150449, [32414], [694774, 2076125, 868414, 2018475, 805946, 2071809], None),
    (4987, 67, 123757, 48500, [73635, 19869, 48169, 51372],
     [4442991, 8144060, 4877263, 380549, 6950975, 6158320, 1974666, 6176397, 5253585, 1252949, 386566,
      4743300],
     20),
    (5531, 67, 456679, 354949, [66263], [13492994, 10236271, 29687795, 8121829, 998005, 6796141], 48),
    (7793, 7, 8676949, 53311, [2440220, 7120207],
     [47468441, 36208810, 40965210, 10992255, 9376840, 4186197, 16117699, 37062360, 50964727, 51690888,
      55784170, 14475955],
     None),
    (8803, 7, 3690553, 1948595, [2996947, 406668],
     [8815803, 5883377, 22165657, 16841644, 18115922, 5558487, 21143425, 24222692, 15670997, 17810208,
      17876235, 22812918],
     None),
    (10079, 67, 1516363, 1218236, [1502877, 501887, 170590],
     [22963847, 5401033, 17745207, 95986099, 34477424, 4955046, 58616053, 987807, 44127334, 63407475,
      45755429, 19263232, 65980708, 56282645, 25846970, 98846133, 48276060, 92153957],
     4),
    (10487, 13, 8460589, 1096523, [3323769, 1838188, 5031352, 4417339],
     [97354128, 47985733, 47869103, 108928309, 103479442, 88921433, 103473384, 17143998, 1654402,
      34464283, 9064779, 62407775, 31849914, 93380852, 69547640, 31893087, 70574769, 60893084, 95210403,
      16679175, 47678867, 7187510, 64906727, 101305885, 5431958, 64807469, 36581752, 12544360, 32287696,
      14121527, 91021837, 18681199, 6519278, 76873666, 84567144, 14113241, 77875792, 13333250, 81054566,
      20681096, 1073106, 71783577, 18910162, 1414447, 41250022, 64173389, 108958051, 108822088],
     None),
    (10883, 67, 1767919, 832718, [1505733],
     [28975844, 51678554, 27679652, 21963560, 89453282, 45531508], 40),
    (11821, 67, 695263, 105769, [313529], [20036661, 30179799, 44687740, 20034978, 30826029, 46042563],
     12),
    (11953, 7, 6804103, 1928470, [5524694, 3560144, 2110832],
     [39235633, 40893523, 19000771, 1804328, 41117268, 37664203], None),
    (16657, 7, 13212967, 3028876, [2276809], [78019061, 43731173, 81663872, 6615200, 2948748, 42137839],
     None),
    (16699, 7, 13279681, 21911, [463381, 2210724],
     [85279988, 11954066, 23060725, 82285979, 55385898, 33487863, 72481295, 10209203, 74840919,
      15286977, 24166360, 25406249],
     None),
    (16823, 7, 40432879, 15212592, [31620714, 18642031, 37419102],
     [198156557, 32286341, 6128196, 197293482, 61895264, 116077943, 102827343, 280585577, 34775037,
      110937146, 218544987, 265324979, 213979106, 65754705, 172594027, 50964973, 10198287, 14512979],
     None),
    (17959, 7, 15359221, 8616,
     [11206319, 4082994, 14196990, 10244480, 189997, 8950843, 2116933, 11133249, 2193736, 10339726,
      2821778, 2662011],
     [89526772, 73544424, 45271239, 14217324, 32917907, 40036414, 10996388, 20071377, 35743150,
      46325919, 1831812, 31964339, 104423443, 30761243, 106562751, 102853549, 23762778, 99950349,
      43372573, 52533389, 33602473, 35832849, 64836708, 54706660, 37483151, 11460314, 57563937,
      39965967, 52721005, 23114180, 47740138, 11609601, 29904186, 99799618, 34563121, 66095812],
     None),
    (21347, 7, 65102251, 22090485, [19137059],
     [173319954, 235432167, 261858420, 134030473, 382239356, 179239408], None),
    (31393, 67, 4903243, 4499034, [4628771],
     [207107778, 299886402, 6916597, 238525795, 53809789, 129820672], 58),
    (32939, 7, 155001523, 137389042, [92704767],
     [1075075511, 201623552, 78703028, 265115187, 924872337, 323086465], None),
    (33749, 7, 162718393, 61179805,
     [48384652, 159668474, 148190180, 151431778, 82184141, 18395180, 86042971, 109735479],
     [260690568, 973276214, 212469905, 807478345, 886815534, 226144166, 1080971032, 384094934,
      990087143, 218301617, 111307386, 482043033, 1114578212, 629377932, 785023686, 1062518489,
      178473593, 131161474, 895719236, 611002423, 889460210, 874736251, 819856082, 346793520],
     None),
    (34267, 7, 55917217, 28006376, [48209706, 14791617, 47461801, 25601532],
     [140981746, 328902867, 26115561, 86601820, 349391478, 34321388, 270571004, 100318175, 171611907,
      383643312, 132525963, 237269484, 327157192, 174741737, 379022732, 307199594, 6777678, 350346900,
      361062837, 284297075, 139270616, 288964337, 284796674, 16448042],
     None),
    (35933, 7, 184459489, 144883764, [140394749, 46110848, 179319446, 173959134],
     [538593224, 21662545, 1085383818, 673612330, 380751530, 1212635222, 650102646, 727603335,
      130280562, 976461219, 924102849, 1130925997, 687147087, 909662165, 874786488, 231031891,
      508478282, 1050963042, 137320830, 487054527, 1105983820, 353108296, 280330000, 70269317],
     None),
    (36343, 67, 6571393, 2719704, [3320947, 5784393, 720591, 6509952, 5177743, 2155046],
     [44861607, 141902594, 114962575, 93707417, 402350112, 166623613, 367777665, 67584519, 18384907,
      264065863, 809944, 82548938, 231166791, 370326749, 429255199, 168818724, 409413619, 173481948,
      760587, 352294811, 53623439, 239054059, 119013285, 393801584, 349665656, 317987642, 188859262,
      77357693, 61868687, 393936598, 196094637, 18002660, 79337077, 229203144, 333151293, 146486129],
     39),
    (36353, 7, 188796709, 102736586,
     [105205497, 35506479, 102841834, 63646069, 162169863, 161177477, 167851581, 183577402, 150711100,
      100434003, 178137332, 39812511, 21783139, 12664035, 77077090, 116238135, 125306126, 164055765],
     [630060085, 914995222, 775196512, 39667695, 269580340, 703165702, 1179824378, 861450851,
      1028668061, 1214081039, 1146682637, 72889652, 442375302, 333875826, 430144671, 702082762,
      1130183297, 363211691, 811987767, 344076073, 741833452, 941534145, 168236799, 393134288,
      446363004, 456353472, 144916117, 536857022, 579373736, 966402333, 1254455512, 950892254,
      1028956931, 774430059, 679960657, 971280446],
     None),
    (37549, 67, 7014751, 595805, [1638954],
     [415330154, 458172434, 405538956, 385363495, 239357620, 39679775], 65),
    (39157, 67, 7628407, 39416, [6330966],
     [431475188, 85423205, 37198095, 451042192, 510723768, 287042297], 56),
    (40699, 13, 42473059, 36263348, [40917577],
     [118688904, 160372228, 90505457, 95825905, 221101899, 270946886, 378755064, 78581949, 110883952,
      415923866, 391148623, 327872656],
     None),
    (43801, 7, 91360543, 22263101, [50807059, 41913772, 8055527],
     [608635123, 154131000, 342181428, 72730030, 191684594, 109646308], None),
    (46351, 7, 102307693, 56953496,
     [35243188, 87246657, 54426909, 46446442, 62343957, 14152173, 55997029, 83925580, 72539152,
      35808959, 74825653, 35301373],
     [683299976, 580202155, 182043004, 612420337, 434437858, 392897450, 33786450, 587206271, 640644895,
      125537504, 349750062, 174769942],
     None),
    (47339, 67, 33448183, 5617033, [7545491, 21213455, 16285611],
     [2077938581, 290440320, 705227566, 25796773, 1301451016, 1143737910, 2041860988, 252480454,
      147198831, 1732904019, 439405583, 1068160152, 230731975, 1828494363, 491261244, 813597635,
      1359705802, 1337874929],
     21),
    (47741, 67, 34018669, 20984962, [3544686, 5820533],
     [1055132187, 404974278, 1810117834, 1497286037, 57022267, 1080465700, 2240811140, 2152074642,
      1221675500, 1978021654, 1460104886, 1477121819],
     20),
    (47819, 7, 326672083, 183970904,
     [195514778, 117927219, 62211443, 21984097, 314720950, 304575684, 119619762, 196514283, 200682832],
     [1782576181, 1290092562, 425713074, 2129161444, 1911534195, 1514507910, 27171068, 1945310250,
      756535651, 412018902, 96296111, 1926607766, 1219970581, 915572524, 395252785, 899185953,
      1301624490, 2156358091, 408699170, 1886580942, 1025940081, 2114587744, 1120609844, 2265114276,
      1918244644, 99913368, 2160907971, 998388122, 1208684859, 35572421, 1895190105, 856540802,
      718643182, 216054082, 204757172, 722821124, 880015746, 934752200, 1842524699, 832952284,
      439305660, 1587880169, 1432982479, 1758032559, 275597131, 1859886498, 2228185039, 1914112898,
      401240244, 363189046, 1677501075, 2274938493, 1460251448, 1244173493],
     None),
    (48947, 67, 35759071, 5654718, [1125438],
     [2373245876, 2017145018, 2212093455, 299966749, 972440187, 1869305014], 56),
    (52967, 67, 41873971, 12929849, [3871718],
     [1012340965, 2043587385, 1754792984, 159201193, 681953988, 273599963], 7),
    (57389, 67, 49157533, 21484646, [15775414, 13010601],
     [1272675698, 2058194404, 2662498656, 1655839189, 712633046, 28999519, 1642437225, 717025945,
      585621923, 1207693482, 3037312579, 2683283288],
     40),
    (57731, 7, 476132299, 433328269, [171336550],
     [2851431825, 927146728, 1979426556, 455976190, 1670183682, 3036900002], None),
    (58111, 7, 160806973, 118130478, [121601835, 97002456, 119362506, 10823616],
     [662699498, 801997222, 227083706, 619941811, 867663628, 64849216, 853014556, 1097455521, 494918756,
      1105965750, 360818678, 64299216],
     None),
    (58477, 13, 87682513, 64852616, [50390575, 60467197, 42141956, 21377785, 76916155, 67253172],
     [233724806, 27385192, 376787180, 181799790, 450559994, 163725, 783956084, 205208659, 25889502,
      842136665, 421419507, 528370806, 248452034, 478622845, 736910808, 116498730, 1078139373,
      913366464, 181446016, 982265317, 1055837765, 816873714, 197642686, 81063276],
     None),
    (60029, 7, 514791553, 366613010, [216450881, 149378958],
     [2417217956, 1830682185, 1332331547, 3202277994, 3057980704, 2367844719, 2779700223, 1059747261,
      2127464301, 2877579871, 2445980219, 2205773559],
     None),
    (61057, 13, 95590213, 25560133, [14741062, 17117334, 30210776, 88444163],
     [501919291, 339338671, 222643500, 473953424, 854920457, 67724114, 1041312866, 673293492, 272442628,
      940702480, 514880583, 205007455, 58548934, 165463534, 310341687, 282577453, 168725611, 888723422,
      453246367, 901197667, 966536404, 1210544354, 401319313, 517781838, 732406453, 1192321049,
      899243498, 595469406, 1155856504, 626644659, 117134255, 780581775, 778516125, 1171307212,
      1090387087, 970240151, 352140826, 385939705, 792166825, 33466206, 1064950739, 728382519,
      936404802, 375385632, 232162756, 621976180, 1189724115, 756919639],
     None),
    (62347, 67, 19339357, 11157766, [15954377, 2178894],
     [346119461, 229944570, 899531159, 984668428, 391621337, 415642810, 113668724, 274253319, 773234852,
      646510260, 349286754, 1251850875],
     41),
    (70379, 67, 73929463, 70851137, [38154359, 37716440, 63244401],
     [756820262, 2828298701, 2353236356, 2291039268, 3722063567, 4821591784, 3237671008, 1631236294,
      2861621525, 3374772011, 2443376834, 1407730595, 4759171204, 295064556, 4103586609, 1446379922,
      4023356007, 4063671995],
     16),
    (71861, 67, 77075749, 72648020, [3076188, 41139491],
     [1206004139, 539607675, 961613527, 4053993162, 1576003928, 1936314157, 4808871655, 1876858087,
      1535321014, 299771550, 3147697116, 3819092446],
     31),
    (73771, 13, 139544467, 133833515, [105390346],
     [927545089, 962739846, 1372818671, 289243788, 168997094, 1310883576, 665427455, 459301595,
      916316782, 957080811, 1270139193, 1624780087],
     None),
    (73859, 7, 779317963, 271732298, [604097653],
     [458211616, 3274252872, 3299968836, 874478651, 1349795009, 3466260515], None),
    (74131, 67, 27340693, 25463160, [24058621, 25332920, 1597359, 7975871],
     [1764716792, 1093282045, 928147921, 738096903, 1369405909, 219217944, 1323010201, 1095528939,
      1320811549, 15845074, 1762255392, 32238286],
     38),
    (76493, 7, 835893649, 451425189, [833050494, 664821364],
     [3605536856, 2066092405, 2948767084, 880219668, 826669548, 867547285, 1363160002, 5386969498,
      701804519, 4446497388, 5453522476, 5208050665],
     None),
    (77893, 7, 288923683, 119137600, [180356652],
     [750967190, 1413482172, 424526178, 1069249605, 1852974117, 902221041], None),
    (78427, 67, 30601357, 7501815, [9805672, 4754104],
     [1254544297, 1402487605, 1753311469, 211815648, 1539537570, 1246858327, 1165474134, 25647786,
      1036775915, 1257410939, 1105335845, 191778851],
     52),
    (79193, 7, 895944349, 382819642, [342067659, 408707201, 595011751, 28366340, 241298698, 741399838],
     [1932867791, 4543323605, 4371200283, 3175932544, 5339658860, 3851123754, 5859283602, 4063535961,
      1317275080, 1041607489, 2919231691, 2070469746, 5596461412, 5654644307, 3832381706, 5138438831,
      4950769331, 4718153824, 5741660705, 4140819115, 521992361, 3054338011, 5676272865, 2711329510,
      2071004022, 6146182992, 1110913395, 4768177404, 2288431051, 4106038897, 5712189134, 2060091358,
      2614704222, 2951432473, 2604999732, 4407080674],
     None),
    (79829, 13, 490211467, 376237680, [257810225],
     [4652786765, 2965582946, 5719671057, 2551607805, 3610946176, 893431038, 4377088273, 4407690570,
      2042303592, 3218469638, 1248785741, 867662101],
     None),
    (82457, 7, 971319901, 512010047, [256566418, 498380174, 492502671, 808620270],
     [5454963476, 2653358049, 1577769003, 2105297252, 67293866, 922465431, 695705566, 5886973516,
      237173067, 4662453659, 313333190, 4782858833],
     None),
    (83243, 67, 103425079, 25012967, [39055127],
     [3813793950, 5934568295, 3902939538, 5452354334, 1630346860, 5302482410], 32),
    (87719, 7, 1099244383, 402987394, [1024310669],
     [5368626222, 4950912634, 238860989, 7079862522, 4176522866, 6705944012], None),
    (89849, 7, 1153276093, 942784979, [73051314, 779310381, 433499004, 850636303],
     [3878323153, 6981058246, 7362933010, 3267268251, 24556125, 6292781557, 2984946326, 1375570120,
      4108827894, 1965721557, 3243817382, 1657662290],
     None),
    (90703, 7, 391767853, 74651683, [300372243, 132279354, 48201269, 28359552],
     [106439272, 560471816, 2465366002, 2551746119, 51615412, 2102354477, 2160709686, 1856936351,
      1125570504, 2041182116, 2567432398, 1162527041],
     None),
    (95461, 7, 433947523, 169250685, [277927111],
     [1093897911, 98367834, 2752543579, 995148285, 528101614, 2393053032], None),
    (97849, 67, 47634451, 5014751, [6897225, 26966346, 5834349],
     [184653881, 2933719674, 2160756961, 898703989, 1657382415, 1944940130, 191360718, 148929776,
      3150891440, 2472039756, 2216024893, 2257743482, 2398920426, 1520839707, 547214786, 2111564075,
      2656569745, 710738789],
     3),
    (98251, 67, 48026653, 40890907,
     [5847308, 21471492, 44031507, 36515053, 30272219, 9944176, 17114209, 30771030, 17772757, 31555103,
      32283886, 17623633],
     [1221635494, 1070890461, 1444284690, 341121497, 1796398153, 3079286472, 1371078002, 992442233,
      1954874235, 2608278782, 2476860102, 570398353, 2802122025, 2986444002, 2236055651, 494094177,
      755183566, 1261267073, 1595820602, 84639011, 1826040936, 738640788, 1758869709, 353635124,
      2202376705, 1796665096, 1327204893, 2268100015, 2484237212, 986519552, 1942162984, 2360390915,
      2958951390, 2155932306, 29571727, 2938347102],
     4),
    (104147, 67, 161891071, 95875488, [17138075, 88319996, 602876],
     [545441744, 6367041891, 2540762612, 4816080875, 1350729078, 1088361559, 312092221, 4579807699,
      4510833569, 1657703013, 1300317222, 1382548756, 8499831813, 4070390664, 8120982831, 2676720256,
      3492601450, 4049522445, 7574376234, 2475674328, 3585530274, 3080829440, 7808846775, 10731211277,
      7615296139, 4441480126, 10748347125, 5919441017, 3575988886, 1594300654, 6701724652, 9693332155,
      7934387496, 3545713151, 7778046231, 9861586700, 2008207339, 8889964228, 4626050357, 8244199126,
      2232671823, 2579326275, 7780629255, 8000611355, 4950133361, 607074924, 1586000651, 1842689808,
      760572087, 6018434941, 9383501057, 1232609600, 10624017610, 6149257323, 8282320318, 10686075613,
      10639131107, 5337993662, 7807133160, 4063982149, 1917010201, 10124747392, 4086102067, 2637761593,
      2063667313, 456196723, 3522545443, 5574345631, 7347696984, 5200893922, 10591095221, 4949446348,
      5742215823, 10641186036, 9420836142, 264793820, 3006053156, 8104148790, 7528518451, 8200415675,
      336329151, 8487240667, 1615521320, 933941813, 6306477540, 618460623, 8366710007, 4477925277,
      2400041911, 2032426450, 10740897849, 259907954, 10264895524, 3725081491, 3198583861, 7225431719,
      7050838504, 4750261460, 4715686516, 5637199022, 4444717839, 2772851992, 8044606575, 1797877667,
      10695629009, 1626107796, 2488986008, 504341187, 7617436187, 4890905423, 9298657537, 2174055786,
      992376184, 1510606508, 3125352054, 7385969240, 8976210885, 6907534697, 754705902, 5551128635,
      1717372182, 6064369397, 3322720385, 8446673935, 1509305575, 337939542, 2567971840, 8750528958,
      8330963359, 8453003900, 5345742494, 8261972376, 10643400811, 10066128074, 5102029355, 3304237001,
      4856067210, 5180772708, 5608625680, 2056395725, 6042146015, 1504894677, 6902486252, 7806060858,
      3829397003, 1054088563, 4304863591, 2915833399, 4451069562, 153502758, 8957615072, 3602096578,
      9582287134, 7441964467, 5533668026, 6096045199, 3951425201, 3044161029, 6331272456, 738034042,
      4778922643, 9459577631, 627231168, 2905392278, 8996120910, 6133051709, 2453744769, 3901826230,
      1102279974, 10664349186, 7504500342, 2157289141, 5869595646, 5516508436, 1138293297, 9329382936,
      973353776, 4685744532, 4731546439, 7813989660, 4248819575, 6808213272, 7825862433, 2426308449,
      1299598595, 9638451240, 4610216431, 5368747486, 4342091214, 6955111718, 10267161989, 1489030409,
      9630656831, 1614197438, 2925967805, 10604318374, 4956136500, 6543757186],
     9),
    (131759, 67, 259112923, 5888300, [163196198, 51817638, 27009030],
     [11255717965, 1663827800, 6771338641, 7950405076, 16103227230, 10815232871, 13708592326,
      15215776470, 4346897272, 374609642, 13174903523, 16934233321, 6419134748, 14887739419, 7781800482,
      3047543792, 12168223440, 15739689313],
     33),
    (132697, 67, 87605107, 78958763, [69438997],
     [3839584903, 5022353898, 1723262313, 160170690, 5614503411, 3671858436], 6),
    (133769, 67, 267078793, 33096160, [27968379, 78303857, 23564243, 159175464],
     [12525057026, 7851264942, 13175833696, 1146888727, 2569275466, 16051598241, 12368827301,
      5899449252, 4307908941, 8909928564, 9965454883, 9233377067, 17011540372, 8465243614, 15920705406,
      5582032284, 5747402042, 13998999653, 14212033403, 15800360019, 6800907448, 13006468936,
      1583445695, 11172159066],
     44),
    (134707, 67, 90279157, 55648125, [20556765, 64815276],
     [4381267331, 4720012965, 90170821, 449552755, 1632221452, 2717740214, 3319176257, 3152526331,
      3834017347, 5276490043, 2088858116, 2945901067, 5750556110, 319668285, 5907965107, 5040655311,
      937518782, 2560131030, 3124439199, 4721214143, 1574385695, 3860736749, 2876229336, 4882781,
      444377369, 3134420995, 3186770420, 4785614221, 4843401920, 2457007129, 3285392880, 5005126334,
      4983579134, 3296847137, 2325651608, 126211874, 1984791235, 169370318, 1465948286, 1983726135,
      2323851232, 1889145612, 2272013887, 2400984723, 1165255373, 412733115, 2131466482, 444196567,
      4052635949, 791279550, 97253196, 3888526155, 257938580, 1296866316, 3684796596, 4137435928,
      789678474, 3540346195, 4846973681, 1196776004, 1914114146, 240991686, 4611755763, 5134911953,
      4855363126, 1423116529, 4192240056, 2605456600, 1460352449, 5977287211, 3188318275, 1148516563,
      3239239921, 3674944565, 987529906, 3175194625, 4579959594, 1514038357, 2674326502, 3342508868,
      1388017904, 349128255, 2441334443, 2169262801, 3726865571, 809826249, 4154438961, 1236896154,
      3347898394, 4843450195, 190237222, 545366595, 4133490793, 1629084658, 5125952416, 3015366135,
      2904645389, 3841269913, 2315007239, 5324385012, 277336711, 1885496353, 4770917060, 4944064480,
      79165558, 5799149845, 5178227828, 5054520594, 3570157119, 4899793121, 2580950910, 1233966291,
      2567625812, 5343636314, 953939639, 5673328714, 3734107399, 3396347047, 2314708479, 3958342011,
      1992432929, 3830348317, 3112784460, 3506411147, 490369960, 4602687416, 5587791873, 4986412760,
      1157141443, 2314059623, 1971733947, 3968743423],
     44),
    (142211, 67, 301852399, 209743788, [151089094],
     [13816589610, 13125460582, 9154318460, 17832888045, 5083372823, 13660142579], 31),
    (150779, 67, 339320263, 101114423, [216264724],
     [16006500059, 18701258832, 8461805402, 5751044968, 20769761283, 16637688085], 29),
    (154799, 67, 357655003, 15516254, [204101760, 25551901, 155532165],
     [10948834164, 16608790485, 22770963084, 2201432634, 12239082083, 8199107888, 2945697743,
      9676316064, 14320322160, 4769178009, 619171891, 5620321988, 22280300438, 9635007184, 13473543594,
      23063688795, 22828490517, 5798778677],
     13),
    (155209, 67, 119850691, 72464659, [40522579, 916334, 48103140],
     [852777920, 6679515576, 990322763, 3511747375, 3430144260, 3026487563, 5496503262, 5855071363,
      6101593742, 609049200, 2992094056, 231847214, 4572714137, 3639872371, 1803878514, 6914464378,
      7707836258, 797417892],
     38),
    (176641, 67, 155234923, 100789324, [58462401],
     [9409743734, 2782409348, 5706925423, 8136813675, 220268711, 1037726941], 37),
    (198491, 67, 588042919, 581459294, [291162641],
     [6211299311, 20281461582, 35312711070, 29053480047, 17536269016, 29850229078], 62),
    (222067, 67, 245343157, 18621495, [220839636, 68316256, 213522200, 195449063],
     [13385572469, 5854212671, 9903722375, 6651323036, 1798965693, 3132138738, 12503501697, 9826425977,
      14776183141, 15574599309, 9785310397, 3950587232],
     16),
    (226631, 67, 766594579, 319269706, [31375108],
     [11232566890, 8077038505, 2030649578, 32113448890, 14405150680, 13284625322], 21),
    (259453, 67, 334906063, 64167256, [330587347],
     [2311302437, 2201767617, 17781247833, 17363677062, 11062633460, 20276337035], 27),
    (275399, 67, 1132013203, 225150291, [1020180894, 535024697, 1096686870],
     [37011748157, 31344374939, 2986469781, 361402333, 36655194468, 36684454196, 55771642549,
      65269431042, 58247829392, 1442805617, 66184871761, 71437762564, 24424876811, 71101630876,
      14162940761, 4806728529, 36255379071, 7976463610],
     10),
    (278347, 67, 385459357, 294595561, [295525321, 276800574, 96192723, 281354937],
     [19469941031, 15750607951, 2953523736, 14688826323, 16782382380, 7524466621, 14921652207,
      6617424046, 21645665404, 13927810674, 1067765047, 22679745520],
     21),
    (294829, 67, 432459871, 430762393, [74572775, 54757131, 125841026, 93459145, 191483926],
     [12953385269, 8368844699, 26440259614, 3617369863, 27902809204, 17754312552, 22586669568,
      18442806727, 12010963858, 14506985463, 2973081597, 19475492790, 849957373, 8713348639, 9536439933,
      277367942, 11788346980, 20866062215, 14229513451, 3545749735, 17674701937, 2938471828,
      22111793845, 26721154868, 17076148041, 18645848973, 10806225254, 14166277397, 15252952499,
      27876287979],
     20),
    (298187, 67, 1327101271, 973297854, [737928903, 623647403, 499907998],
     [27822103787, 79302831054, 4161929709, 53594161005, 59495972659, 49118183532, 54164557014,
      56769749642, 38388761237, 81014599690, 64708881702, 60332666733, 13382631310, 50796797500,
      88349501629, 83635493372, 52624405743, 66254532061, 78592047680, 81202784910, 39690248065,
      37650749236, 81227354711, 63214970556, 35157022592, 67237120229, 48288787349, 81832812350,
      53065913852, 45840524519, 32846597740, 76086854479, 73571623856, 20439099579, 60266622907,
      24369723847, 50357970413, 27988486253, 45733371968, 59256382842, 79141975652, 3400885958,
      79554922713, 37930495575, 29177040622, 22787004489, 47693470272, 83252946293, 7374203189,
      43273568826, 16536317512, 29425988337, 35578944388, 50180853958, 40107762828, 82470625796,
      43180446397, 19813509785, 82131240423, 78660941661, 11215935146, 13832264191, 45024021597,
      16965470207, 25448473675, 9987410627, 39288916306, 36887806667, 20352255390, 57226936083,
      25065050644, 65330390501, 85216676711, 42849570602, 10705730552, 27246272869, 32223558782,
      45218815596, 43484097384, 60438934520, 50540668735, 7787494177, 52422405825, 28250611024,
      88097667241, 8139684939, 78307855446, 19664754960, 22359291485, 2381797743, 59954073688,
      45930392049, 56424650391, 1562361305, 79369816390, 25870432157, 21680232411, 74023809847,
      15164731500, 85055064337, 55378017857, 26216918936, 9112813491, 32486534375, 41224402911,
      69201944478, 74646626842, 61802013631, 58584875861, 83881142170, 60834906886, 49949180230,
      83449089241, 13487764755, 84615768070, 11505407638, 81558545649, 68994047696, 45549091908,
      59733990180, 15879295853, 64576694929, 77671882204, 26625156482, 42375888783, 59417824460,
      8435813502, 66680230655, 38795793484, 59812440319, 39908990472, 61672041864, 53280984503,
      57206465314, 3627899747, 49571452482, 72903578749, 77427609146, 9791011436, 616132240, 9864117820,
      47901082185, 40293465787, 11409100008, 18810757961, 3693650323, 44334520385, 74060186155,
      86941198462, 43821336958, 43890976504, 13130061454, 3044009891, 21732885751, 66810969984,
      17026987385, 60753087429, 7940995980, 17022060020, 29011427081, 35472719367, 60446968508,
      12715779234, 1345912398, 80663686543, 31999467456, 87431019190, 10246662289, 46560471256,
      55138713702, 8187569268, 72014870589, 47856405580, 56125820150, 16161881242, 2596007287,
      33462285897, 77018768668, 21161969689, 46376011398, 9846652261, 50193209873, 2582586866,
      33015826937, 710285381, 19300770134, 70643556555, 62075700835, 27087053884, 20428278001,
      60556925161, 1910050265, 49456728503, 56427452098, 66859812869, 46100492110, 48200882618,
      59593656828],
     58),
    (332617, 67, 550419907, 503474519, [45214974],
     [19354405224, 30399107453, 34447474579, 7081165415, 18403567947, 22343141510], 42),
    (336503, 67, 1690068739, 1394597284, [502978979],
     [81525555859, 33433406208, 66730305152, 105066795835, 68849212308, 102755909161], 62),
    (344543, 67, 1771794379, 1758632887, [1205049381],
     [87419672046, 34196194898, 97168819597, 102362578780, 110076423724, 101935278904], 37),
    (356737, 67, 633142507, 587255922, [600349744, 329171336, 31281846],
     [34588095080, 2983715078, 42239017193, 7189011190, 41183946818, 15601873240, 38547957333,
      3190229628, 22020494743, 22115280331, 9338512052, 39961340920, 36373480388, 23654246386,
      34845295812, 33233122658, 37298273938, 8903439817],
     11),
    (357281, 67, 1905224929, 1883433119,
     [1274332726, 1073043281, 237079317, 875831147, 1539402295, 476291212, 1155626863, 1788086850,
      426952768, 1312776602, 1533917871, 1480949105],
     [84688965187, 112747433900, 121048934715, 78739675942, 57474190576, 99927111060, 8369556155,
      23552870668, 61309024042, 15214208946, 64979391274, 28105073642, 99545125723, 29794916898,
      1635454072, 49742532801, 89175950901, 28619319835, 104698994080, 94271245586, 88026517995,
      46207854719, 72649870316, 25432444496, 125744389258, 107823046722, 118246288669, 123901998977,
      112507363146, 8824003659, 125824028842, 48281969711, 7680741636, 2690285229, 116665306329,
      90774630147, 27887583634, 25781520548, 64630328705, 93823454182, 92676944816, 87140942328,
      38059025451, 53732213888, 114525147522, 77928828274, 1707914079, 111817832261, 76991403950,
      90649638255, 115483398871, 58233010377, 99501277942, 76877836752, 47954267210, 74039594432,
      122808134622, 21039601949, 8356361200, 102675191488, 64880766946, 44091784229, 22968842468,
      28721230791, 89483337403, 91157682927, 88841133426, 19674967829, 46042379879, 50491096809,
      39226461536, 9115007377],
     64),
    (373621, 67, 694492663, 585960269, [194513373],
     [28890428145, 25541604372, 28839532618, 41124069530, 25307367660, 8745422870], 52),
    (378713, 67, 2140655449, 1596047346, [1194113807, 270640838],
     [95668890469, 118309174063, 32110040356, 59861372091, 56994796720, 36435205543, 87424686509,
      54755606849, 34912153992, 82487586571, 107174781630, 72307031644],
     38),
    (381937, 67, 725752507, 617042177, [687942622],
     [30005769077, 2248247516, 17983027879, 28697383732, 12535614925, 18035804456], 58),
    (385153, 67, 738025963, 707538938, [363590660, 452274233, 308303833],
     [8353247126, 32992006749, 22282867191, 2047212955, 35330466469, 15746621561, 25135002365,
      31214118164, 21623371510, 15773778710, 35592495911, 14010713600, 13596801642, 33529388507,
      31905430526, 6565385495, 36015782814, 44688183585],
     4),
    (388897, 67, 752444107, 577226010, [299734287],
     [47719615682, 22306291887, 36642860015, 32122581503, 19706691511, 11789384031], 43),
    (393587, 67, 2312106271, 392796347, [1027457999, 1293741307, 986726537],
     [43847885591, 84126958905, 7227172533, 56754061456, 82531770268, 33788168759, 86918164835,
      10164104193, 64968459514, 89528623015, 46115251681, 32873103169, 109183852135, 54766031089,
      13196950538, 135842082845, 75717231316, 132040951367],
     64),
    (395597, 67, 2335781821, 851909424, [1093028079, 1411360450, 1359056915, 864454297],
     [16261219668, 97786086619, 104153171260, 5800938056, 38593432153, 70023688938, 107020049874,
      108699414061, 128338739499, 31975577179, 61019205390, 147220964086],
     32),
    (720413, 67, 7746203149, 3475884807,
     [4196781196, 515388264, 1822604618, 5294902932, 2217084810, 7615241370, 1406771527, 4957123836,
      6315684921, 5506166111, 6186024087, 4513368547],
     [281518282442, 150924243795, 372883274698, 456040186283, 185023015423, 264307571721, 9894013048,
      174105409096, 438268514799, 141309647520, 127794811294, 103362888210, 91614920063, 236148274137,
      216155326973, 416461232584, 272933242874, 133856496723, 234999303588, 433111178200, 64127293959,
      371688067102, 260572728343, 10863742572, 45943771637, 318004609953, 49513011454, 83555032528,
      338938106644, 386513187519, 319458223852, 26017187315, 34729365339, 485833008308, 160214257619,
      84149832627],
     34),
    (989761, 67, 4873770283, 1806255930, [3021816622],
     [247982562148, 62422771496, 218225628396, 95440274468, 303428801206, 23914441181], 34),
    (1118267, 67, 18664510471, 13699147176, [15075541790, 15075383787, 1795963401],
     [318329569833, 1161302768705, 284324706219, 1192740927550, 294986908233, 324196317740,
      794797330881, 852765350014, 903933132829, 398660867209, 441464161945, 782372617981, 475351883122,
      297446528574, 117888209497, 1010756722741, 28632413296, 27895322940],
     34),
    (1253071, 67, 7811881513, 1721835122,
     [2007906336, 859378333, 2963205532, 197851444, 1414397587, 3645350444, 1962938644, 2579981686,
      5021903661, 6256733830, 4913829747, 5599765042, 2958711942, 4377305399, 7490020533, 793416521,
      7582294555, 5978938025, 4286592917, 5566616105, 7603774960, 5287954141, 4984585754, 6218552250],
     [296012638913, 405741783004, 432496635891, 401547808417, 77746901181, 9487166722, 404643551480,
      394802239081, 316162983369, 310291600129, 317859784406, 353756800908, 457551478420, 253056203362,
      87781813542, 360218194885, 227798973164, 284519682664, 508465536518, 28444323596, 395011872763,
      368774522959, 256469118269, 376996640783, 462272242706, 322471398972, 102172772113, 349673361486,
      310102635483, 444685048472, 364008618025, 421895539900, 210164486729, 473016758315, 273990947662,
      119838186162, 308872262516, 456502941613, 493117419242, 127196179091, 133836430791, 344655869322,
      499283709372, 185778963120, 486763594247, 39625718056, 453547056101, 264307490525, 320564008105,
      102143439962, 282271453315, 500086880872, 365063255979, 13126613918, 205740192233, 305305477518,
      390584046532, 375109443473, 226389939584, 442559602952, 403664486666, 389884224465, 212678305108,
      417985114119, 456308564360, 371870894446, 244842696080, 311551275853, 26086612279, 414313250068,
      381028570361, 125829887016],
     34),
    (1672751, 67, 41762650459, 16677457130, [26358052608],
     [932644942803, 727498919603, 1607331262730, 2255154426540, 1015058584816, 19167961535], 34),
]
for record in records:
    verify(*record)
from collections import Counter
slopes=Counter(x[-1] for x in records if x[1]==67)
assert slopes[34]==5 and max(v for k,v in slopes.items() if k!=34)==3
assert len(slopes)==34 and all(x[2] not in (2,3,5,7,13,67,827,5717) for x in records)
print("Verified all 94 local equivalences.")

"""Verify the D=61 norm-one exponent-lattice bound. Python 3, standard library.

The optional --full flag rechecks the earlier twenty-two exhaustive finite-field
certificates. This part traverses several million cubic finite-field elements.
"""
from fractions import Fraction as F
from math import gcd, isqrt
import sys

DATA = [
    (5,1),(41,787),(97,2951),(103,1915),
    (127,4419),(131,14136),(167,11077),(199,8942),
    (271,5781),(293,50847),(773,552873),(827,419023),
    (859,50622),(911,21647),(1021,260342),(1163,549066),
    (1483,349678),(1583,255229),(1709,2745692),(1723,654871),
    (1811,205927),(1993,630756),(3557,8372873),(4889,3215080),
    (8681,46418403),(10589,44659163),(13469,31295481),(16229,77599160),
    (16301,66708745),(16421,76985174),(16553,128280463),(19013,116496427),
    (19661,95133794),(20849,270934284),(22277,350529041),(22937,297631469),
    (26489,31896576),(26921,404448787),(28901,65818922),(31277,924860164),
    (32693,1008873991),(36497,695170719),(36749,1317741455),(37049,1219522376),
    (37649,1390873606),(38669,243832387),(38993,1250966396),(42797,987290451),
    (46229,1742607656),(47309,941740647),(49937,5343544),
]
ONE=(1,0,0)
E=(-56,37,2)
T=(-46,27,2)

def mul(a,b,p):
    c=[0]*5
    for i in range(3):
        for j in range(3):
            c[i+j]+=a[i]*b[j]
    for k in (4,3):
        c[k-1]-=19*c[k]; c[k-2]+=19*c[k]; c[k-3]+=19*c[k]
    return tuple(x%p for x in c[:3])

def power(a,n,p):
    z=ONE
    while n:
        if n%2: z=mul(z,a,p)
        a=mul(a,a,p); n//=2
    return z

def project(a,p):
    c=pow(next(x%p for x in a if x%p),-1,p)
    return tuple(x*c%p for x in a)

def prime(n):
    return n>=2 and all(n%d for d in range(2,isqrt(n)+1))

def factors(n):
    d=2
    while d*d<=n:
        if n%d==0:
            yield d
            while n%d==0: n//=d
        d+=1
    if n>1: yield n

def trac(a,p):
    return (3*a[0]-19*a[1]+399*a[2])%p

def lattice():
    N,h=1,0
    for n,(p,k) in enumerate(DATA):
        S=p*p+p+1; m=S//gcd(S,3)
        assert prime(p) and prime(m) and gcd(N,m)==1
        assert all((t**3+19*t*t-19*t-19)%p for t in range(p))
        assert project(E,p)!=ONE and project(power(E,m,p),p)==ONE
        assert project(power(E,k,p),p)==project(T,p)
        if n>=22:
            assert p%3==2
            assert all(pow(3,(m-1)//d,m)!=1 for d in factors(m-1))
        elif '--full' in sys.argv:
            seen=bytearray(m); x=ONE
            for i in range(m):
                if not trac(x,p):
                    j=i
                    while not seen[j]:
                        seen[j]=1; j=3*j%m
                x=mul(E,x,p)
            assert sum(seen)==m-1 and not seen[0]
        h+=N*((k-h)*pow(N,-1,m)%m); N*=m
    # Necessary mixed-prime rows certified by the mixed-prime verifier.
    for m,k in [(827,170),(5717,2537)]+[(x[2],x[3]) for x in records]:
        assert gcd(N,m)==1
        h+=N*((k-h)*pow(N,-1,m)%m); N*=m
    def dp(a,b): return sum(x*y for x,y in zip(a,b))
    u,v=(N,0),(-h,1)
    while True:
        if dp(v,v)<dp(u,u): u,v=v,u
        d=dp(u,u); q=(2*dp(u,v)+d)//(2*d)
        if not q: break
        v=tuple(x-q*y for x,y in zip(v,u))
    assert abs(u[0]*v[1]-u[1]*v[0])==N
    return u,v

def normform(x,y,z):
    return (x**3-x*x*z-76*x*y*y-38*x*y*z-101*x*z*z
            -228*y**3-532*y*y*z+646*y*z*z+387*z**3)

def units():
    for p in (2,19):
        for x in range(p):
            for y in range(p):
                for z in range(p):
                    b=3*x*x-2*x*z-76*y*y-38*y*z-101*z*z
                    if (3*x-z)%p==b%(p*p)==normform(x,y,z)%(p**3)==0:
                        assert (x,y,z)==(0,0,0)
    solutions=set()
    for y in range(-26,27):
        for z in range(-34,35):
            q=-76*y*y-38*y*z-101*z*z; r=normform(0,y,z)
            for x in range(-230,231):
                if ((x-z)*x+q)*x+r==1: solutions.add((x,y,z))
    assert solutions=={(1,0,0),(-24,5,1),(46,15,-9),(-132,-22,19)}

def interval(a,b=None):
    if b is None:b=a
    return (F(a),F(b))
def add(a,b):return a[0]+b[0],a[1]+b[1]
def neg(a):return -a[1],-a[0]
def prod(a,b):
    s=[x*y for x in a for y in b]
    return min(s),max(s)
def inv(a):
    assert a[0]*a[1]>0
    return 1/a[1],1/a[0]
def sub(a,b):return add(a,neg(b))
def ab(a):
    assert a[0]*a[1]>0
    return a if a[0]>0 else neg(a)
def dot(a,b):return add(prod(a[0],b[0]),prod(a[1],b[1]))
def inv2(M):
    a,b=M[0];c,d=M[1];D=inv(sub(prod(a,d),prod(b,c)))
    return [[prod(t,D) for t in R] for R in [[d,neg(b)],[neg(c),a]]]
def sqrt_ge(a,x):return a>=x*x
def natlog_one_two(x):
    t=(x-1)/(x+1)
    t=F(t); t2=t*t; q=t; s=0;n=18
    for k in range(n):s+=2*q/(2*k+1);q*=t2
    return (s,s+2*q/((2*n+1)*(1-t2)))
logtwo=natlog_one_two(F(2))
def logpoint(x):
    assert x>0;k=0
    while x<1:x*=2;k-=1
    while x>2:x/=2;k+=1
    return add(natlog_one_two(x),prod(interval(k),logtwo))
def logiv(a):
    lo,hi=logpoint(a[0])[0],logpoint(a[1])[1]
    f=10**8
    return interval(F((lo*f).__floor__(),f),F((hi*f).__ceil__(),f))
def box(n):return interval(F(n,10**12),F(n+1,10**12))
theta=[box(t) for t in [-19906514237356,-623731016520,1530245253875]]
def poly(a,coeffs):
    v=interval(0)
    for c in coeffs:v=add(prod(v,a),interval(c))
    return v
for a,b in theta:
    def f(x):return x**3+19*x*x-19*x-19
    assert f(a)*f(b)<0
es=[poly(t,[2,37,-56]) for t in theta]
ts=[poly(t,[2,27,-46]) for t in theta]
logs=[[logiv(ab(e)),logiv(ab(t))] for e,t in zip(es,ts)]
G=[[sum((prod(R[i],R[j])[0] for R in logs),F()),sum((prod(R[i],R[j])[1] for R in logs),F())] for i,j in [(0,0),(1,1),(0,1)]]
g,h,j=G
assert g[0]>47 and (g[0]-47)*(h[0]-47)>max(j[0]**2,j[1]**2)
def box_certificate():
    bounds=(300,71,268)
    for i in range(3):
        e,t=ab(es[i]),ab(ts[i])
        assert max(e[1],1/e[0])*max(t[1],1/t[0]) < bounds[i]**2
    R=[ [interval(1),poly(t,[F(1,3),6,F(-19,3)]),
          poly(t,[F(1,3),7,F(-1,3)])] for t in theta]
    def cofactor(i,j):
        u,v=[k for k in range(3) if k!=i]
        r,s=[k for k in range(3) if k!=j]
        det=sub(prod(R[u][r],R[v][s]),prod(R[u][s],R[v][r]))
        return neg(det) if (i+j)%2 else det
    D=interval(0)
    for j in range(3):D=add(D,prod(R[0][j],cofactor(0,j)))
    for i,maximum in enumerate((230,26,34)):
        assert sum(bounds[j]*ab(prod(inv(D),cofactor(j,i)))[1]
                   for j in range(3))<maximum
    units()

box_certificate()
U,V=lattice()
S=U[0]**2+U[1]**2;D=abs(U[0]*V[1]-U[1]*V[0])
# Binary Gauss reduction: U is a shortest nonzero vector.
assert S<=sum(v*v for v in V)
assert 2*abs(sum(u*v for u,v in zip(U,V)))<=S
assert len(records)==94 and S>2073*10**1122
assert len(str(S))==1126
# ||log vector||² ≤ (1+(17/40)²+(23/40)²)*log(beta)²
# Gram matrix certificate above gives ||log vector||²>47*(a²+b²).
C=1+F(17,40)**2+F(23,40)**2
assert 47*S>C*(10**563*logiv(interval(10))[1])**2
assert 2*10**13*3**1152<10**563
print("Verified log10(beta) > 10^563.")

"""Exact check of a simultaneous Pisot-prime-sieve survivor. Python 3."""
from math import gcd, isqrt, lcm

EXTRA=[3049, 3557, 4889, 7793, 8681, 8803, 10487, 10589, 11953, 13469, 16229, 16301, 16421, 16553, 16657, 16699, 16823, 17959, 19013, 19661, 20849, 21347, 21503, 22277, 22937, 26489, 26921, 28901, 31277, 32693, 32939, 33749, 34267, 34303, 35933, 36353, 36497, 36749, 37049, 37649, 38669, 38993, 40699, 42797, 43801, 46229, 46351, 47309, 47819, 49937, 57731, 58111, 58477, 60029, 61057, 73771, 73859, 76493, 77893, 79193, 79829, 82457, 87719, 89849, 90703, 95461, 524287, 720413, 989761, 1118267, 1253071, 1672751]
def prime(n):
    return n>=2 and all(n%i for i in range(2,isqrt(n)+1))
Q=[p for p in range(2,3001) if prime(p)]+EXTRA

def mul(a,b,p):
    w=a[2]*b[2]
    v=a[1]*b[2]+a[2]*b[1]-19*w
    return ((a[0]*b[0]+19*v)%p,(a[0]*b[1]+a[1]*b[0]+19*w+19*v)%p,
            (a[0]*b[2]+a[1]*b[1]+a[2]*b[0]+19*w-19*v)%p)

def power(a,n,p):
    z=(1,0,0)
    while n:
        if n%2:z=mul(z,a,p)
        a=mul(a,a,p)
        n//=2
    return z

def clock(p,L):
    E=(-56,37,2); T=(-46,27,2)
    # Compare in the common projective unit group; trace 0 is scalar invariant.
    D=p*(p*p-1)*(p**3-1)
    base=mul(E,power(T,23,p),p)
    assert power(base,D,p)==(1,0,0)
    start=power(base,L%D,p); x=start
    k=next(i for i in range(3) if start[i])
    for n in range(1,786434):
        assert (3*x[0]-19*x[1]+399*x[2])%p
        x=mul(mul(x,x,p),x,p)
        if x[k] and all((x[i]*start[k]-start[i]*x[k])%p==0 for i in range(3)):
            return n
    raise ValueError(p)

assert len(Q)==len(set(Q))==502 and all(prime(p) for p in Q)
L=lcm(*(p*(p*p-1)*(p**3-1) for p in Q))
for r in (2,7,67):
    while L%r==0:L//=r
assert len(str(L))==4024
R=lcm(*(p*(p*p-1)*(p**3-1) for p in Q))//L
new=[4987, 5531, 10079, 10883, 11821, 31393, 36343, 37549, 39157, 47339, 47741, 48947, 52967, 57389, 62347, 70379, 71861, 74131, 78427, 83243, 97849, 98251, 104147, 131759, 132697, 133769, 134707, 142211, 150779, 154799, 155209, 176641, 198491, 222067, 226631, 259453, 275399, 278347, 294829, 298187, 332617, 336503, 344543, 356737, 357281, 373621, 378713, 381937, 385153, 388897, 393587, 395597]
assert len(new)==52 and all(prime(p) for p in new) and len(set(new)|set(Q))==554
F=1
for p in new:
    S=p*p+p+1
    q=S//gcd(3,S)//67
    assert S//gcd(3,S)==67*q and prime(q) and gcd(q,R)==1
    F*=q
L*=F*pow(F,-1,R)
clocks=[clock(p,L) for p in Q]
extra_clocks=[clock(p,L) for p in new]
assert (sum(clocks),max(clocks))==(795441,786432)
print(len(Q)+len(new),sum(clocks+extra_clocks),max(clocks+extra_clocks))
print("extra clocks",sum(extra_clocks),max(extra_clocks))
print("3-adic valuation",next(j for j in range(80) if L%(3**(j+1))!=0))

assert L%3**13==0 and L%3**14!=0
assert {r[0] for r in records}|{p for p,k in DATA}|{21503,34303,524287}|{p for p in range(2,3001) if prime(p)} == set(Q+new)
assert pow(61,2400,4801)==4800 and prime(4801) and 4801 not in Q+new
assert all(power(power(mul((-56,37,2),power((-46,27,2),23,p),p),L%(p*p-1),p),27,p)[1:]==(0,0) for p in (59,647))
assert (L%128,(23*L)%128,L%25,L%2401,(23*L)%2401)==(25,63,0,2364,1550)
"""Certificate that p=4801 hits every member of the 554-prime ray.
Append to the construction of L if its numeric residues also need checking.
"""
from math import gcd
E=(-56,37,2);T=(-46,27,2);ONE=(1,0,0)
def mul(a,b,p):
    w=a[2]*b[2]; v=a[1]*b[2]+a[2]*b[1]-19*w
    return ((a[0]*b[0]+19*v)%p,
            (a[0]*b[1]+a[1]*b[0]+19*w+19*v)%p,
            (a[0]*b[2]+a[1]*b[1]+a[2]*b[0]+19*w-19*v)%p)
def pw(a,n,p):
    z=ONE
    while n:
        if n%2:z=mul(z,a,p)
        a=mul(a,a,p);n//=2
    return z
def pr(a,p):
    t=pow(next(x%p for x in a if x%p),-1,p)
    return tuple(x*t%p for x in a)
def check_row(p,l,o,e,t):
    v=pw((0,1,0),p,p)
    if pow(61,(p-1)//2,p)==p-1:s=p*p-1
    elif v==(0,1,0):s=p-1
    else:s=p*p+p+1
    h=s
    while h%l==0:h//=l
    x=pw(E,h,p);y=pw(T,h,p)
    gen=x if e==1 else y
    assert pr(pw(gen,o,p),p)==ONE and pr(pw(gen,o//l,p),p)!=ONE
    assert pr(pw(gen,e,p),p)==pr(x,p) and pr(pw(gen,t,p),p)==pr(y,p)
def prove():
    pairs=[(128,2,(21503,2048,1,1598),(34303,1024,438,1)),
           (25,5,(1249,625,140,1),(2999,125,1,105)),
           (343,7,(1733,343,1,331),(1373,343,1,292))]
    for d,l,a,b in pairs:
        for p,o,e,t in (a,b):
            assert o%d==0
            check_row(p,l,o,e,t)
        assert (a[2]*b[3]-a[3]*b[2])%l
    p=4801;s=p*p-1; assert s==128*3*25*2401
    assert pr(pw(E,s//3,p),p)==pr(pw(T,s//3,p),p)==ONE
    def crt(x,d):
        k=s//d;return (k*pow(k,-1,d)*x)%s
    hits=[802,3154,5506,1159,1373,771,3123]
    H=16464
    for i in range(7):
        for j in range(7):
            a=crt(25,128)+crt(2364+343*i,2401)
            b=crt(63,128)+crt(1550+343*j,2401)
            z=mul(pw(E,a,p),pw(T,b,p),p)
            assert pr(pw(z,s,p),p)==ONE
            h=hits[(i+j)%7]
            y=pw(z,pow(3,h,s),p)
            assert (3*y[0]-19*y[1]+399*y[2])%p==0
            assert pr(pw(z,pow(3,H,s),p),p)==pr(z,p)
    print("Verified 4801 in every 16464 late indices on the 554-prime ray.")
if __name__=='__main__':prove()

"""Unbounded two-prime obstruction for noncube cubic Mills shadows."""
E=(-56,37,2);T=(-46,27,2); ONE=(1,0,0)
def mul(a,b,p):
    w=a[2]*b[2]
    v=a[1]*b[2]+a[2]*b[1]-19*w
    return ((a[0]*b[0]+19*v)%p,(a[0]*b[1]+a[1]*b[0]+19*w+19*v)%p,
            (a[0]*b[2]+a[1]*b[1]+a[2]*b[0]+19*w-19*v)%p)
def power(a,n,p):
    z=ONE
    while n:
        if n%2:z=mul(z,a,p)
        a=mul(a,a,p);n//=2
    return z
for p,c,ord3 in [(59,2320,3),(647,25840,27)]:
    assert pow(61,(p-1)//2,p)==p-1
    n=p*p-1; h=1
    while n%3==0:n//=3;h*=3
    assert c%h==1 and c%n==0
    e,t=power(E,c,p),power(T,c,p)
    g=t if p==59 else e
    if p==59:assert e==ONE
    else:assert t==power(g,4,p)
    assert power(g,ord3,p)==ONE
    for i in (1,2):
        x,y,z=power(g,ord3*i//3,p)
        assert (3*x-19*y+399*z)%p==0 and (y,z)!=(0,0)
print("Two-prime obstruction at floors j=0 and j=2 verified.")

# Full projective generators give the prime-to-three escape conditions.
for p,h,u,v,A,B,m in [(59,3,2589,1853,1,-903,1160),
                       (647,4,198246,394950,-551,1,2584)]:
    g=(h,1,0); n=p*p-1
    assert pr(pw(g,n,p),p)==ONE
    assert all(pr(pw(g,n//r,p),p)!=ONE for r in fac(n))
    assert pr(pw(g,u,p),p)==pr(E,p) and pr(pw(g,v,p),p)==pr(T,p)
    while n%3==0:n//=3
    D=gcd(gcd(u,v),n)
    assert n//D==m
    cu,cv=u//D,v//D
    assert gcd(gcd(cu,cv),m)==1 and (cu*B-cv*A)%m==0
print("Verified the two escape congruences.")
'''

# Source: Mills_Frobenius_Franel.tex
DESCRIPTIONS['07-franel.py'] = 'Ten further rows, the new degree bound and the two orbits at 59.'
DIGESTS['07-franel.py'] = '9c8747729503ca39bd9757bdcdec229c9be7cf45c2998345fc6bee7e7efa6471'
PROGRAMS['07-franel.py'] = r'''from math import gcd,isqrt
ONE=(1,0,0);E=(-56,37,2);T=(-46,27,2)
def mul(a,b,p):
    q=a[2]*b[2];r=a[1]*b[2]+a[2]*b[1]-19*q
    return ((a[0]*b[0]+19*r)%p,
            (a[0]*b[1]+a[1]*b[0]+19*(r+q))%p,
            (a[0]*b[2]+a[1]*b[1]+a[2]*b[0]+19*(q-r))%p)
def power(a,n,p):
    x=ONE
    while n:
        if n&1:x=mul(x,a,p)
        a=mul(a,a,p);n//=2
    return x
def project(a,p):
    k=pow(next(v%p for v in a if v%p),-1,p)
    return tuple(v*k%p for v in a)
def zero(a,p):return (3*a[0]-19*a[1]+399*a[2])%p==0
def prime(n):return n>=2 and all(n%d for d in range(2,isqrt(n)+1))
def factors(n):
    t=2
    while t*t<=n:
        if n%t==0:
            yield t
            while n%t==0:n//=t
        t+=1
    if n>1:yield n
def verify(p,n,k,o,wits):
    s=p*p+p+1
    assert p not in (2,3,19,61) and prime(p) and prime(n) and n==s//gcd(3,s)
    assert pow(61,(p-1)//2,p)==1 and power((0,1,0),p,p)!=(0,1,0)
    assert project(power(E,n,p),p)==ONE and project(E,p)!=ONE
    assert project(power(E,k,p),p)==project(T,p)
    assert pow(3,o,n)==1 and all(pow(3,o//r,n)!=1 for r in factors(o))
    assert pow(p,3,n)==1 and p%n!=1 and o%3
    d=3*o
    assert all(0<u<n and zero(power(E,u,p),p) for u in wits)
    assert len({pow(u,d,n) for u in wits})==len(wits)==(n-1)//d
ROWS=[
(3491, 12190573, 5409730, 677254, [4133, 8773, 9347, 44143, 65387, 113742]),
(6389, 40825711, 3297004, 4536190, [1650, 9576, 14582]),
(10781, 116240743, 7713735, 12915638, [42399, 43084, 68303]),
(11299, 42559567, 25827896, 14186522, [7098]),
(14593, 70990081, 62632233, 5915840, [22172, 32722, 112723, 129249]),
(16267, 88210519, 49352284, 2673046, [3089, 44056, 57109, 75992, 98257, 144353, 272230, 291807, 323847, 351828, 388768]),
(18427, 113190919, 3023671, 37730306, [15834]),
(18773, 352444303, 347944980, 39160478, [1556, 35454, 40208]),
(19237, 123360469, 17606603, 20560078, [22652, 46607]),
(22643, 512728093, 273688730, 28484894, [19566, 56343, 68298, 74708, 83709, 757213]),
]
if __name__=="__main__":
    for row in ROWS:verify(*row)
    print("Verified",len(ROWS),"Frobenius-compressed scalar obstructions.")

"""Verify the D=61 norm-one exponent-lattice bound. Python 3, standard library.

The optional --full flag rechecks the earlier twenty-two exhaustive finite-field
certificates. This part traverses several million cubic finite-field elements.
"""
from fractions import Fraction as F
from math import gcd, isqrt
import sys

DATA = [
    (5,1),(41,787),(97,2951),(103,1915),
    (127,4419),(131,14136),(167,11077),(199,8942),
    (271,5781),(293,50847),(773,552873),(827,419023),
    (859,50622),(911,21647),(1021,260342),(1163,549066),
    (1483,349678),(1583,255229),(1709,2745692),(1723,654871),
    (1811,205927),(1993,630756),(3557,8372873),(4889,3215080),
    (8681,46418403),(10589,44659163),(13469,31295481),(16229,77599160),
    (16301,66708745),(16421,76985174),(16553,128280463),(19013,116496427),
    (19661,95133794),(20849,270934284),(22277,350529041),(22937,297631469),
    (26489,31896576),(26921,404448787),(28901,65818922),(31277,924860164),
    (32693,1008873991),(36497,695170719),(36749,1317741455),(37049,1219522376),
    (37649,1390873606),(38669,243832387),(38993,1250966396),(42797,987290451),
    (46229,1742607656),(47309,941740647),(49937,5343544),
]
ONE=(1,0,0)
E=(-56,37,2)
T=(-46,27,2)

def mul(a,b,p):
    c=[0]*5
    for i in range(3):
        for j in range(3):
            c[i+j]+=a[i]*b[j]
    for k in (4,3):
        c[k-1]-=19*c[k]; c[k-2]+=19*c[k]; c[k-3]+=19*c[k]
    return tuple(x%p for x in c[:3])

def power(a,n,p):
    z=ONE
    while n:
        if n%2: z=mul(z,a,p)
        a=mul(a,a,p); n//=2
    return z

def project(a,p):
    c=pow(next(x%p for x in a if x%p),-1,p)
    return tuple(x*c%p for x in a)

def prime(n):
    return n>=2 and all(n%d for d in range(2,isqrt(n)+1))

def factors(n):
    d=2
    while d*d<=n:
        if n%d==0:
            yield d
            while n%d==0: n//=d
        d+=1
    if n>1: yield n

def trac(a,p):
    return (3*a[0]-19*a[1]+399*a[2])%p

def lattice():
    N,h=1,0
    for n,(p,k) in enumerate(DATA):
        S=p*p+p+1; m=S//gcd(S,3)
        assert prime(p) and prime(m) and gcd(N,m)==1
        assert all((t**3+19*t*t-19*t-19)%p for t in range(p))
        assert project(E,p)!=ONE and project(power(E,m,p),p)==ONE
        assert project(power(E,k,p),p)==project(T,p)
        if n>=22:
            assert p%3==2
            assert all(pow(3,(m-1)//d,m)!=1 for d in factors(m-1))
        elif '--full' in sys.argv:
            seen=bytearray(m); x=ONE
            for i in range(m):
                if not trac(x,p):
                    j=i
                    while not seen[j]:
                        seen[j]=1; j=3*j%m
                x=mul(E,x,p)
            assert sum(seen)==m-1 and not seen[0]
        h+=N*((k-h)*pow(N,-1,m)%m); N*=m
    # Necessary mixed-prime rows certified by the mixed-prime verifier.
    for m,k in [(827, 170), (5717, 2537), (442831, 150449), (8676949, 53311), (3690553, 1948595), (6804103, 1928470), (13212967, 3028876), (13279681, 21911), (40432879, 15212592), (15359221, 8616), (65102251, 22090485), (155001523, 137389042), (162718393, 61179805), (55917217, 28006376), (184459489, 144883764), (188796709, 102736586), (91360543, 22263101), (102307693, 56953496), (326672083, 183970904), (476132299, 433328269), (160806973, 118130478), (514791553, 366613010), (779317963, 271732298), (835893649, 451425189), (288923683, 119137600), (895944349, 382819642), (971319901, 512010047), (1099244383, 402987394), (1153276093, 942784979), (391767853, 74651683), (433947523, 169250685), (8460589, 1096523), (42473059, 36263348), (87682513, 64852616), (95590213, 25560133), (139544467, 133833515), (490211467, 376237680), (4657, 4170), (88969, 64229), (123757, 48500), (456679, 354949), (1516363, 1218236), (1767919, 832718), (695263, 105769), (4903243, 4499034), (6571393, 2719704), (7014751, 595805), (7628407, 39416), (33448183, 5617033), (34018669, 20984962), (35759071, 5654718), (41873971, 12929849), (49157533, 21484646), (19339357, 11157766), (73929463, 70851137), (77075749, 72648020), (27340693, 25463160), (30601357, 7501815), (103425079, 25012967), (47634451, 5014751), (48026653, 40890907), (161891071, 95875488), (259112923, 5888300), (87605107, 78958763), (267078793, 33096160), (90279157, 55648125), (301852399, 209743788), (339320263, 101114423), (357655003, 15516254), (119850691, 72464659), (155234923, 100789324), (588042919, 581459294), (245343157, 18621495), (766594579, 319269706), (334906063, 64167256), (1132013203, 225150291), (385459357, 294595561), (432459871, 430762393), (1327101271, 973297854), (550419907, 503474519), (1690068739, 1394597284), (1771794379, 1758632887), (633142507, 587255922), (1905224929, 1883433119), (694492663, 585960269), (2140655449, 1596047346), (725752507, 617042177), (738025963, 707538938), (752444107, 577226010), (2312106271, 392796347), (2335781821, 851909424), (7746203149, 3475884807), (4873770283, 1806255930), (18664510471, 13699147176), (7811881513, 1721835122), (41762650459, 16677457130)]+[(x[1],x[2]) for x in ROWS]:
        assert gcd(N,m)==1
        h+=N*((k-h)*pow(N,-1,m)%m); N*=m
    def dp(a,b): return sum(x*y for x,y in zip(a,b))
    u,v=(N,0),(-h,1)
    while True:
        if dp(v,v)<dp(u,u): u,v=v,u
        d=dp(u,u); q=(2*dp(u,v)+d)//(2*d)
        if not q: break
        v=tuple(x-q*y for x,y in zip(v,u))
    assert abs(u[0]*v[1]-u[1]*v[0])==N
    return u,v

def normform(x,y,z):
    return (x**3-x*x*z-76*x*y*y-38*x*y*z-101*x*z*z
            -228*y**3-532*y*y*z+646*y*z*z+387*z**3)

def units():
    for p in (2,19):
        for x in range(p):
            for y in range(p):
                for z in range(p):
                    b=3*x*x-2*x*z-76*y*y-38*y*z-101*z*z
                    if (3*x-z)%p==b%(p*p)==normform(x,y,z)%(p**3)==0:
                        assert (x,y,z)==(0,0,0)
    solutions=set()
    for y in range(-26,27):
        for z in range(-34,35):
            q=-76*y*y-38*y*z-101*z*z; r=normform(0,y,z)
            for x in range(-230,231):
                if ((x-z)*x+q)*x+r==1: solutions.add((x,y,z))
    assert solutions=={(1,0,0),(-24,5,1),(46,15,-9),(-132,-22,19)}

def interval(a,b=None):
    if b is None:b=a
    return (F(a),F(b))
def add(a,b):return a[0]+b[0],a[1]+b[1]
def neg(a):return -a[1],-a[0]
def prod(a,b):
    s=[x*y for x in a for y in b]
    return min(s),max(s)
def inv(a):
    assert a[0]*a[1]>0
    return 1/a[1],1/a[0]
def sub(a,b):return add(a,neg(b))
def ab(a):
    assert a[0]*a[1]>0
    return a if a[0]>0 else neg(a)
def dot(a,b):return add(prod(a[0],b[0]),prod(a[1],b[1]))
def inv2(M):
    a,b=M[0];c,d=M[1];D=inv(sub(prod(a,d),prod(b,c)))
    return [[prod(t,D) for t in R] for R in [[d,neg(b)],[neg(c),a]]]
def sqrt_ge(a,x):return a>=x*x
def natlog_one_two(x):
    t=(x-1)/(x+1)
    t=F(t); t2=t*t; q=t; s=0;n=18
    for k in range(n):s+=2*q/(2*k+1);q*=t2
    return (s,s+2*q/((2*n+1)*(1-t2)))
logtwo=natlog_one_two(F(2))
def logpoint(x):
    assert x>0;k=0
    while x<1:x*=2;k-=1
    while x>2:x/=2;k+=1
    return add(natlog_one_two(x),prod(interval(k),logtwo))
def logiv(a):
    lo,hi=logpoint(a[0])[0],logpoint(a[1])[1]
    f=10**8
    return interval(F((lo*f).__floor__(),f),F((hi*f).__ceil__(),f))
def box(n):return interval(F(n,10**12),F(n+1,10**12))
theta=[box(t) for t in [-19906514237356,-623731016520,1530245253875]]
def poly(a,coeffs):
    v=interval(0)
    for c in coeffs:v=add(prod(v,a),interval(c))
    return v
for a,b in theta:
    def f(x):return x**3+19*x*x-19*x-19
    assert f(a)*f(b)<0
es=[poly(t,[2,37,-56]) for t in theta]
ts=[poly(t,[2,27,-46]) for t in theta]
logs=[[logiv(ab(e)),logiv(ab(t))] for e,t in zip(es,ts)]
G=[[sum((prod(R[i],R[j])[0] for R in logs),F()),sum((prod(R[i],R[j])[1] for R in logs),F())] for i,j in [(0,0),(1,1),(0,1)]]
g,h,j=G
assert g[0]>47 and (g[0]-47)*(h[0]-47)>max(j[0]**2,j[1]**2)
def box_certificate():
    bounds=(300,71,268)
    for i in range(3):
        e,t=ab(es[i]),ab(ts[i])
        assert max(e[1],1/e[0])*max(t[1],1/t[0]) < bounds[i]**2
    R=[ [interval(1),poly(t,[F(1,3),6,F(-19,3)]),
          poly(t,[F(1,3),7,F(-1,3)])] for t in theta]
    def cofactor(i,j):
        u,v=[k for k in range(3) if k!=i]
        r,s=[k for k in range(3) if k!=j]
        det=sub(prod(R[u][r],R[v][s]),prod(R[u][s],R[v][r]))
        return neg(det) if (i+j)%2 else det
    D=interval(0)
    for j in range(3):D=add(D,prod(R[0][j],cofactor(0,j)))
    for i,maximum in enumerate((230,26,34)):
        assert sum(bounds[j]*ab(prod(inv(D),cofactor(j,i)))[1]
                   for j in range(3))<maximum
    units()

box_certificate()
U,V=lattice()
S=U[0]**2+U[1]**2;D=abs(U[0]*V[1]-U[1]*V[0])
# Binary Gauss reduction: U is a shortest nonzero vector.
assert S<=sum(v*v for v in V)
assert 2*abs(sum(u*v for u,v in zip(U,V)))<=S

from math import isqrt
B=isqrt(int(47*S/((1+F(17,40)**2+F(23,40)**2)*(logiv(interval(10))[1]**2))))
print('S length',len(str(S)),'lower bound log10 beta',str(B)[:10],len(str(B)))
n=len(str(B))-1
t=0
while 2*10**13*3**t<10**n:t+=1
print(n,t,t+1, 'for pure power lower')
t=0
while 2*10**13*3**t<B:t+=1
assert B>514*10**600 and 2*10**13*3**1235<514*10**600 and t==1236
print('Verified degree bound: 3^1237.')

"""Finite arithmetic for the index 3,457,120 family in the cubic field."""
from math import gcd
E=(-56,37,2); T=(-46,27,2); ONE=(1,0,0)
def mul(a,b,p):
 w=a[2]*b[2];v=a[1]*b[2]+a[2]*b[1]-19*w
 return ((a[0]*b[0]+19*v)%p,(a[0]*b[1]+a[1]*b[0]+19*w+19*v)%p,
         (a[0]*b[2]+a[1]*b[1]+a[2]*b[0]+19*w-19*v)%p)
def pw(a,n,p):
 z=ONE
 while n:
  if n%2:z=mul(z,a,p)
  a=mul(a,a,p);n//=2
 return z
def pr(a,p):return tuple(x*pow(next(x%p for x in a if x%p),-1,p)%p for x in a)
def prime(n):return n>1 and all(n%i for i in range(2,1+int(n**.5)))
for p,h,n,f,e,t in [(17,1,288,[2,3],45,145),(2789,6,1296420,[2,3,5,17,31,41],318371,175467)]:
 assert prime(p) and pow(61,(p-1)//2,p)==p-1
 assert n==(p*p-1)//h
 g=pw((4,1,0),h,p)
 assert pr(pw(g,n,p),p)==ONE
 m=n
 for r in f:
  assert prime(r) and not m%r and pr(pw(g,n//r,p),p)!=ONE
  while m%r==0:m//=r
 assert m==1 and pr(pw(g,e,p),p)==pr(E,p) and pr(pw(g,t,p),p)==pr(T,p)
 for u in range(1,3):
  x,y,z=pw(g,(n//3)*u,p)
  assert (3*x-19*y+399*z)%p==0
print("Verified.")

assert 1585643%32==11 and 1585643%432140==289223
assert 3457120%32==3457120%432140==0
assert (45*11+145)%32==0 and (318371*289223+175467)%432140==0


# Additional squareclass-three certificate. Polynomials are in increasing degree.
def check_prime_59():
    p, n = 59, 3541
    assert all((x**3 - 6*x*x - 3*x + 2) % p for x in range(p))
    assert all(n % d for d in range(2, isqrt(n) + 1))
    assert pow(3, 1770, n) == 1
    assert all(pow(3, 1770//q, n) != 1 for q in (2, 3, 5, 59))
    assert pow(p, 1770, n) == 1
    assert [pow(e, 1770, n) for e in (9, 315)] == [1, n-1]
    v = (1, 0, 0)
    targets = {27: ((26, 4, 6), 354), 945: ((11, 21, 51), 2301)}
    for k in range(1, 946):
        a, b, c = v
        v = (-2*c % p, (a + 3*c) % p, (b + 6*c) % p)
        if k == 3:
            assert v[1:] != (0, 0)
        if k in targets:
            w, trace = targets[k]
            assert v == w and 3*v[0] + 6*v[1] + 42*v[2] == trace and trace % p == 0
    print("Verified two-orbit coverage at p=59 in discriminant squareclass 3.")
check_prime_59()
'''

# Source: Mills_Seven_and_Signed_Norms.tex
DESCRIPTIONS['08-seven.py'] = 'Both unit signs, the divisor 7, two ray norms and gap examples.'
DIGESTS['08-seven.py'] = 'c5a9030cf9a9d3c570a639f493b188e6c1c4bf773e789048eaf93384de495684'
PROGRAMS['08-seven.py'] = r'''"""Completely rational, standard-library proof for a cubic field of squareclass 3."""
from fractions import Fraction as F
from itertools import product
from math import isqrt
U=(-796,-24,29);V=(-65,82,20);one=(1,0,0)
def norm(i,j,k):
    return (i**3+i*i*(j+55*k)+i*(-27*j*j-18*j*k+735*k*k)
            -3*(j**3+j*j*k-27*j*k*k-3*k**3))
def mul(v,w):
    a=[0]*5
    for i,vi in enumerate(v):
        for j,wj in enumerate(w):a[i+j]+=vi*wj
    for k in [4,3]:
        z=a[k];a[k-1]+=z;a[k-2]+=27*z;a[k-3]-=3*z
    return tuple(a[:3])
assert norm(*U)==norm(*V)==1
assert mul(mul(U,V),(38,-350,61))==one
assert all(tuple(c%4 for c in mul(w,w))==one for w in (U,V))
assert (3*U[0]+U[1]+55*U[2])%4==3
assert (3*V[0]+V[1]+55*V[2])%4==3
assert [(w[0]%3,sum(w)%3) for w in (U,V)]==[(2,1),(1,1)]
# Candidate integral basis 1, theta, (1+theta^2)/2 of discriminant
# 80688/2^2 = 20172 = 2^2 *3*41^2. Test no integral vector with these denominators.
for p in (2,41):
    for x,y,z in product(range(p),repeat=3):
        i,j,k=2*x+z,2*y,z
        a=3*i+j+55*k
        b=3*i*i+2*i*(j+55*k)-27*j*j-18*j*k+735*k*k
        if a%(2*p)==b%(4*p*p)==norm(i,j,k)%(8*p**3)==0:
            assert (x,y,z)==(0,0,0)
def I(x):return F(x),F(x)
def add(a,b):return a[0]+b[0],a[1]+b[1]
def neg(a):return -a[1],-a[0]
def sub(a,b):return add(a,neg(b))
def times(a,b):
    p=[x*y for x in a for y in b]
    return min(p),max(p)
def inv(a):
    assert a[0]*a[1]>0
    return 1/a[1],1/a[0]
def absiv(a):
    assert a[0]*a[1]>0
    return a if a[0]>0 else neg(a)
def evalpoly(poly,x):
    v=I(0)
    for c in reversed(poly):v=add(times(v,x),I(c))
    return v
g=(3,-27,-1,1)
roots=[]
for t in [-4779926581195511883,110707433684311705,5669219147511200177]:
    r=(F(t,10**18),F(t+1,10**18)); roots.append(r)
    assert evalpoly(g,I(r[0]))[0]*evalpoly(g,I(r[1]))[0]<0
bound=[1100,220,4100]; mat=[[I(1),z,times(z,z)] for z in roots]
for i,r in enumerate(roots):
    u=evalpoly(U,r);v=evalpoly(V,r)
    assert (u[1]<0 and v[1]<0) if i<2 else (u[0]>0 and v[0]>0)
    a,b=absiv(u),absiv(v)
    assert max(a[1],1/a[0])*max(b[1],1/b[0])<bound[i]**2
assert absiv(evalpoly(U,roots[0]))[0]>1 and absiv(evalpoly(V,roots[0]))[1]<1
assert absiv(evalpoly(U,roots[1]))[0]>1 and absiv(evalpoly(V,roots[1]))[0]>1
def cofactor(i,j):
    a,b=[t for t in range(3) if t!=i];c,d=[t for t in range(3) if t!=j]
    v=sub(times(mat[a][c],mat[b][d]),times(mat[b][c],mat[a][d]))
    return neg(v) if (i+j)%2 else v
D=I(0)
for j in range(3):D=add(D,times(mat[0][j],cofactor(0,j)))
for j,q in enumerate((600,1000,240)):
    assert 2*sum(bound[t]*absiv(times(inv(D),cofactor(t,j)))[1] for t in range(3))<q
# Enumerate norm-1 vectors in the coefficient box. For fixed j,k a monic cubic
# in i can be searched by congruence mod a suitable prime before the exact test.
P=97;sol=set()
modroots={}
for j in range(P):
    for k in range(P):
        modroots[j,k]=[i for i in range(P) if (norm(i,j,k)-8)%P==0]
for j in range(-1000,1001,2):
    for k in range(-240,241):
        for r in modroots[j%P,k%P]:
            for i in range(-600+(r+600)%P,601,P):
                if (k-i)%2==0 and norm(i,j,k)==8:
                    sol.add((i,j,k))
assert sol=={(2,0,0),(-130,164,40),(76,-700,122)}
print("Certified the norm-one unit lattice, signs and residues.")
def powmod(v,n,p):
    a=one
    while n:
        if n%2:a=tuple(t%p for t in mul(a,v))
        v=tuple(t%p for t in mul(v,v));n//=2
    return a
def normproj(a,p):
    s=pow(next(x%p for x in a if x%p),-1,p)
    return tuple(x*s%p for x in a)
p=7;g=normproj(powmod(U,3,p),p)
assert g==(1,2,5)
assert normproj(powmod(V,3,p),p)==normproj(powmod(g,5,p),p)==(1,6,3)
assert normproj(powmod(g,16,p),p)==one and normproj(powmod(g,8,p),p)!=one
zero=set()
for i in range(16):
    x,y,z=powmod(g,i,p)
    if (3*x+y+55*z)%p==0:zero.add(i)
assert zero=={1,7}
assert all(zero & {a*3**j%16 for j in range(4)} for a in range(1,16,2))
print("Verified 7 on both odd-exponent orbits; both unit norms.")

for element, expected in [((425,84,0),9192413), ((509,-84,0),14916929)]:
    assert all(evalpoly(element, root)[0] > 0 for root in roots)
    assert all((a-b)%84 == 0 for a,b in zip(element,(5,0,0)))
    assert norm(*element) == expected and expected%84 == 41
    assert all(expected%d for d in range(2,isqrt(expected)+1))
print("Verified both totally positive ray generators and their prime norms.")

# Two independent integer checks on the valuations of the gap quotient.
def valuations(a, b, c, center, count):
    result = []
    for j in range(count):
        a1 = a**3 - 3*a*b + 3*c
        b1 = b**3 - 3*a*b*c + 3*c*c
        c1 = c**3
        h = c - a*b
        quotient, remainder = divmod(c1 - a1*b1, h)
        assert remainder == 0
        numerator = abs(quotient - center*c*c)
        assert numerator and c%3
        valuation = 0
        while numerator % 3 == 0:
            numerator //= 3
            valuation += 1
        result.append(valuation)
        a, b, c = a1, b1, c1
    return result

assert valuations(-817,14927,1,9,4) == [5,7,9,11]
# x^3 - 7x^2 - 8x - 1 has b=1 mod 3 and norm 1.
first = valuations(7,-8,1,-3,6)
assert [y-x for x,y in zip(first[1:],first[2:])] == [1,1,1,1]
print("Verified gap valuations at centers -3 and +9:", first, [5,7,9,11])
'''


def names_for(selector):
    if selector == "all":
        return list(PROGRAMS)
    if selector in PROGRAMS:
        return [selector]
    matching = [name for name in PROGRAMS if name.startswith(selector + "-")]
    if not matching:
        raise SystemExit("Unknown selector " + repr(selector) + ". Use list for names.")
    return matching


def payload(name):
    data = PROGRAMS[name].encode("utf-8")
    if hashlib.sha256(data).hexdigest() != DIGESTS[name]:
        raise SystemExit("Archive integrity failure: " + name)
    return data


def run_one(name, thorough):
    print("RUN " + name, flush=True)
    with tempfile.TemporaryDirectory(prefix="mills-certificate-") as work:
        folder = Path(work)
        source = folder / name
        source.write_bytes(payload(name))
        if source.suffix == ".py":
            cmd = [sys.executable, "-I", str(source)]
            if thorough and name in {"05-lattice.py", "06-early.py", "07-franel.py"}:
                cmd.append("--full")
        else:
            binary = folder / "certificate"
            subprocess.run(["c++", "-std=c++17", "-O2", "-UNDEBUG",
                            str(source), "-o", str(binary)], cwd=folder, check=True)
            cmd = [str(binary)]
        subprocess.run(cmd, cwd=folder, check=True)
    print("PASS " + name, flush=True)


def main():
    parser = argparse.ArgumentParser(description=__doc__)
    commands = parser.add_subparsers(dest="command", required=True)
    commands.add_parser("list", help="show the original certificate filenames")
    ext = commands.add_parser("extract", help="write exact source listings")
    ext.add_argument("selector", help="filename, chapter number such as 04, or all")
    ext.add_argument("directory", type=Path)
    check = commands.add_parser("run", help="run a single listing or a chapter")
    check.add_argument("selector", help="filename, chapter number such as 04, or all")
    allchecks = commands.add_parser("all", help="run all ten recovered listings")
    for cmd in [check, allchecks]:
        cmd.add_argument("--thorough", action="store_true",
                         help="also repeat inherited orbit checks in the Python listings")
    args = parser.parse_args()
    if args.command == "list":
        for name, description in DESCRIPTIONS.items():
            print(name + "  " + description)
        return
    selected = names_for("all" if args.command == "all" else args.selector)
    if args.command == "extract":
        args.directory.mkdir(parents=True, exist_ok=True)
        for name in selected:
            path = args.directory / name
            data = payload(name)
            if path.exists() and path.read_bytes() != data:
                raise SystemExit("Destination has different content: " + str(path))
            path.write_bytes(data)
            print(path)
        return
    for name in selected:
        run_one(name, args.thorough)


if __name__ == "__main__":
    main()
